% Le néocolonialisme — Allemand Tle A — Cours signé OnBuch+
% Programme officiel MINESEC. Compiler avec : tectonic AL27-le-neocolonialisme.tex
\newcommand{\DOCMATIERE}{Allemand}
\newcommand{\DOCNIVEAU}{Tle A}
\newcommand{\DOCMODULE}{Chapitre 5 : Citoyenneté}
\newcommand{\DOCLECON}{AL27}
\newcommand{\DOCTITRE}{Le néocolonialisme}
% OnBuch+ — préambule « cours signés » ENRICHI (compilé avec tectonic / XeLaTeX)
% Système visuel « Pop » d'OnBuch+ : fond crème (jamais blanc), bordures encre
% épaisses, ombres « dures » décalées, titres Archivo Black en capitales.
% Version MATHÉMATIQUES : graphes (pgfplots/tikz), boîtes pédagogiques
% (définition, propriété, méthode, attention, le savais-tu, exercice résolu,
% exercice, TP, bilan), page de garde et signature OnBuch+ ancrée en bas.
%
% Macros attendues dans main.tex AVANT \input{preamble} :
%   \DOCMATIERE  \DOCNIVEAU  \DOCMODULE  \DOCLECON (numéro)  \DOCTITRE
\documentclass[11pt]{article}

\usepackage{fontspec}
\usepackage[ngerman,french]{babel} % français = langue principale ; l'allemand via \all{..} / textebox
\usepackage[a4paper,margin=2cm,top=3.4cm,bottom=2.8cm,headheight=1.4cm,footskip=1.3cm]{geometry}
\usepackage{amsmath,amssymb,amsfonts}
\usepackage{mathtools} % \xmapsto, \coloneqq... souvent utilisés par les modèles
\usepackage{cancel} % simplifications barrées (bilans de forces)
% \abs{z} (module, valeur absolue) et \norm{u} (norme), utilisés par les modèles
\providecommand{\abs}{}\let\abs\relax\DeclarePairedDelimiter{\abs}{\lvert}{\rvert}
\providecommand{\norm}{}\let\norm\relax\DeclarePairedDelimiter{\norm}{\lVert}{\rVert}
\usepackage[table]{xcolor} % [table] : fournit aussi \rowcolors (alternance de lignes)
\usepackage{pagecolor}
\usepackage{tikz}
\usetikzlibrary{arrows.meta,positioning,calc,shapes.geometric,shapes.misc,shapes.arrows,decorations.pathmorphing,decorations.pathreplacing,decorations.markings,patterns,shadows,mindmap,trees,fit,backgrounds,matrix,intersections,angles,quotes}
\usepackage{pgfplots}
\usepackage[version=4]{mhchem} % équations nucléaires \ce{^{235}_{92}U}
\usepackage[european,siunitx]{circuitikz} % schémas électriques
\pgfplotsset{compat=1.18}
\usepgfplotslibrary{fillbetween,groupplots} % aires entre courbes (calcul intégral)
\usepackage{enumitem}
\usepackage{fancyhdr}
% [most] charge toutes les bibliothèques tcolorbox (listings, minted, raster,
% documentation...) et gonfle inutilement la mémoire TeX sur un document long
% (~300 boîtes) : on ne charge que ce qui est réellement utilisé.
\usepackage[skins,breakable]{tcolorbox}
\usepackage{graphicx}
\usepackage{colortbl}
\usepackage{booktabs}
\usepackage{tabularx}
\usepackage{diagbox} % case d'en-tête diagonale des tableaux à double entrée
% Colonnes extensibles centrée / gauche / droite (C, L, R), souvent utilisées par les modèles
\newcolumntype{C}{>{\centering\arraybackslash}X}
\newcolumntype{L}{>{\raggedright\arraybackslash}X}
\newcolumntype{R}{>{\raggedleft\arraybackslash}X}
\usepackage{array}
\usepackage{multicol}
\usepackage{multirow}
\usepackage{siunitx}
% parse-numbers=false : accepte aussi \SI{2,5 \times 10^{-3}}{...}
\sisetup{locale=FR,output-decimal-marker={,},group-minimum-digits=4,parse-numbers=false}
\DeclareSIUnit\ohm{\ensuremath{\symup{\Omega}}} % U+2126 absent de la police
\DeclareSIUnit\division{div} % réglages d'oscilloscope (ms/div, V/div)
\DeclareSIUnit\tr{tr} % tours (vitesse de rotation, ex. \SI{1500}{\tr\per\minute})
\DeclareSIUnit\molar{M} % concentration molaire (ex. \SI{3}{\molar}, \SI{1}{\micro\molar})
\DeclareSIUnit\an{an} % année (le modèle confond parfois avec \year, un compteur TeX) — ex. \SI{5}{\kilo\meter\per\an}
\DeclareSIUnit\FCFA{FCFA} % franc CFA (ex. \SI{500}{\FCFA\per\kilo\gram})
\DeclareSIUnit\degreeBrix{\textdegree Brix} % teneur en sucre d'un sirop/jus (réfractométrie)
\DeclareSIUnit\month{mois} % le modèle utilise parfois \month, un compteur TeX, comme unité de durée
\usepackage{qrcode}
% \degree : commande fréquemment utilisée par les modèles pour un angle ou une
% température en dehors de \SI{}{} (non fournie par les paquets chargés).
\providecommand{\degree}{^{\circ}}
% \qed (fin de démonstration), utilisé par les modèles sans amsthm
\providecommand{\qed}{\hfill$\square$}
% \barre : double barre des valeurs interdites dans un tableau de variations (tkz-tab)
\providecommand{\barre}{\ensuremath{\|}}
% environnement proof (amsthm non chargé) utilisé par les modèles
\ifdefined\proof\else\newenvironment{proof}[1][Démonstration]{\par\noindent\textit{#1.}\ }{\qed\par}\fi
\usepackage{lastpage}
\usepackage{hyperref}
\hypersetup{hidelinks}

% ============================================================
% Palette « Pop » OnBuch+ — mêmes hex que la classe P (plus_ui.dart)
% ============================================================
\definecolor{poporange}{HTML}{F59321}
\definecolor{popdark}{HTML}{E07A0C}
\definecolor{popcream}{HTML}{FFF6E8}
\definecolor{popink}{HTML}{1C1714}
\definecolor{poppurple}{HTML}{7A5AE0}
\definecolor{popgreen}{HTML}{2E9E6A}
\definecolor{poppink}{HTML}{E0457B}
\definecolor{popblue}{HTML}{2F7DE1}
\definecolor{popgold}{HTML}{C98A12}
\definecolor{popmuted}{HTML}{8A7F76}
\definecolor{popline}{HTML}{EADFD2}
\definecolor{popgray}{HTML}{8A7F76}
\colorlet{popgrey}{popgray}
% Alias tolérants (le modèle nomme parfois les couleurs différemment)
\colorlet{popwhite}{white}
\colorlet{popblack}{popink}
\colorlet{popred}{poppink}
\colorlet{popyellow}{popgold}
% Teintes claires (fonds de tableaux/encadrés) — le modèle nomme parfois
% les couleurs avec un suffixe L (ex. popblueL) sans les avoir définies.
\colorlet{poporangeL}{poporange!20}
\colorlet{popdarkL}{popdark!20}
\colorlet{poppurpleL}{poppurple!20}
\colorlet{popgreenL}{popgreen!20}
\colorlet{poppinkL}{poppink!20}
\colorlet{popblueL}{popblue!20}
\colorlet{popgoldL}{popgold!20}
\colorlet{popredL}{popred!20}
\colorlet{popgrayL}{popgray!20}
% Teintes claires pour fonds de tableaux / zones
\colorlet{popblueL}{popblue!10!white}
\colorlet{poporangeL}{poporange!14!white}
\colorlet{popgreenL}{popgreen!12!white}
\colorlet{poppurpleL}{poppurple!12!white}
\colorlet{poppinkL}{poppink!12!white}
\colorlet{popgoldL}{popgold!14!white}

\pagecolor{popcream}

% ============================================================
% Typographie
% ============================================================
\defaultfontfeatures{Ligatures=TeX}
\setmainfont{PlusJakartaSans-Medium.ttf}[Path=../../../fonts/,BoldFont=PlusJakartaSans-Bold.ttf]
% Symboles, API et emojis absents de la police principale : repli sur DejaVu Sans (caractères actifs)
\newfontface\symfont{DejaVuSans.ttf}[Path=../../../fonts/]
\makeatletter
\newcommand\symactive[1]{\catcode#1=13 \begingroup\lccode`\~=#1 \lowercase{\endgroup\def~}{{\symfont\char#1}}}
\newcommand\symblank[1]{\catcode#1=13 \begingroup\lccode`\~=#1 \lowercase{\endgroup\def~}{}}
\newcommand\symhyph[1]{\catcode#1=13 \begingroup\lccode`\~=#1 \lowercase{\endgroup\def~}{\mbox{-}}}
\newcount\symc
\newcommand\symrange[2]{\symc=#1 \@whilenum\symc<#2 \do{\expandafter\symactive\expandafter{\the\symc}\advance\symc by 1 }}
\newcommand\symblankrange[2]{\symc=#1 \@whilenum\symc<#2 \do{\expandafter\symblank\expandafter{\the\symc}\advance\symc by 1 }}
\makeatother
\symrange{"0250}{"02FF}\symactive{"2713}\symactive{"2714}\symactive{"2717}\symactive{"2718}\symactive{"2610}\symactive{"2611}\symactive{"2612}
\symhyph{"2011}\symhyph{"2E1A}\symblankrange{"1F300}{"1FAFF}\symblank{"FE0F}
% Maths en sans-serif (Fira Math) avec chiffres et lettres droites de Plus
% Jakarta, pour que les formules se fondent dans le texte.
\usepackage[math-style=ISO,bold-style=ISO]{unicode-math}
\setmathfont{FiraMath-Regular.otf}[Path=../../../fonts/]
\setmathfont{PlusJakartaSans-Medium.ttf}[Path=../../../fonts/,range={up/num,up/{Latin,latin},\mathup}]
\usepackage{icomma} % virgule décimale sans espace en mode math
% Caractères mathématiques Unicode absents des polices de texte (Plus Jakarta,
% Archivo Black) : rendus par leur équivalent mathématique, en texte comme en
% formule — sinon ils s'affichent en carré vide (ex. « Arithmétique dans ℤ »).
\usepackage{newunicodechar}
\newunicodechar{ℝ}{\ensuremath{\mathbb{R}}}
\newunicodechar{ℤ}{\ensuremath{\mathbb{Z}}}
\newunicodechar{ℂ}{\ensuremath{\mathbb{C}}}
\newunicodechar{ℕ}{\ensuremath{\mathbb{N}}}
\newunicodechar{ℚ}{\ensuremath{\mathbb{Q}}}
\newunicodechar{𝔻}{\ensuremath{\mathbb{D}}}
\newunicodechar{α}{\ensuremath{\alpha}}
\newunicodechar{β}{\ensuremath{\beta}}
\newunicodechar{γ}{\ensuremath{\gamma}}
\newunicodechar{δ}{\ensuremath{\delta}}
\newunicodechar{ε}{\ensuremath{\varepsilon}}
\newunicodechar{θ}{\ensuremath{\theta}}
\newunicodechar{λ}{\ensuremath{\lambda}}
\newunicodechar{μ}{\ensuremath{\mu}}
\newunicodechar{σ}{\ensuremath{\sigma}}
\newunicodechar{τ}{\ensuremath{\tau}}
\newunicodechar{φ}{\ensuremath{\varphi}}
\newunicodechar{ψ}{\ensuremath{\psi}}
\newunicodechar{ω}{\ensuremath{\omega}}
\newunicodechar{Σ}{\ensuremath{\Sigma}}
% majuscules produites par \MakeUppercase dans les titres (α-aminés -> Α)
\newunicodechar{Α}{\ensuremath{\alpha}}
\newunicodechar{Β}{\ensuremath{\beta}}
\newunicodechar{Γ}{\ensuremath{\Gamma}}
\newunicodechar{Δ}{\ensuremath{\Delta}}
\newunicodechar{Ε}{\ensuremath{\varepsilon}}
\newunicodechar{Θ}{\ensuremath{\Theta}}
\newunicodechar{Λ}{\ensuremath{\Lambda}}
\newunicodechar{Μ}{\ensuremath{\mu}}
\newunicodechar{Τ}{\ensuremath{\tau}}
\newunicodechar{Φ}{\ensuremath{\Phi}}
\newunicodechar{Ψ}{\ensuremath{\Psi}}
\newunicodechar{Ω}{\ensuremath{\Omega}}
\newunicodechar{↦}{\ensuremath{\mapsto}}
\newunicodechar{∈}{\ensuremath{\in}}
\newunicodechar{∉}{\ensuremath{\notin}}
\newunicodechar{∧}{\ensuremath{\wedge}}
\newunicodechar{⇔}{\ensuremath{\Leftrightarrow}}
\newunicodechar{⟺}{\ensuremath{\Longleftrightarrow}}
\newunicodechar{⇒}{\ensuremath{\Rightarrow}}
\newunicodechar{⟹}{\ensuremath{\Longrightarrow}}
\newunicodechar{∘}{\ensuremath{\circ}}
\newunicodechar{∪}{\ensuremath{\cup}}
\newunicodechar{∩}{\ensuremath{\cap}}
\newunicodechar{≡}{\ensuremath{\equiv}}
\newunicodechar{⊂}{\ensuremath{\subset}}
\newunicodechar{⊥}{\ensuremath{\perp}}
\newunicodechar{∃}{\ensuremath{\exists}}
\newunicodechar{∀}{\ensuremath{\forall}}
\newunicodechar{∛}{\ensuremath{\sqrt[3]{\,}}}
\newunicodechar{∜}{\ensuremath{\sqrt[4]{\,}}}
\newunicodechar{⃗}{\ensuremath{{}^{\to}}}
\newunicodechar{ⁿ}{\ifmmode^{n}\else\textsuperscript{\ensuremath{n}}\fi}
\newunicodechar{ˣ}{\ifmmode^{x}\else\textsuperscript{\ensuremath{x}}\fi}
\newunicodechar{ᵇ}{\ifmmode^{b}\else\textsuperscript{\ensuremath{b}}\fi}
\newunicodechar{ʳ}{\ifmmode^{r}\else\textsuperscript{\ensuremath{r}}\fi}
\newunicodechar{ᵘ}{\ifmmode^{u}\else\textsuperscript{\ensuremath{u}}\fi}
\newunicodechar{ʸ}{\ifmmode^{y}\else\textsuperscript{\ensuremath{y}}\fi}
\newunicodechar{ᵃ}{\ifmmode^{a}\else\textsuperscript{\ensuremath{a}}\fi}
\newunicodechar{ᵉ}{\ifmmode^{e}\else\textsuperscript{\ensuremath{e}}\fi}
\newunicodechar{ᵏ}{\ifmmode^{k}\else\textsuperscript{\ensuremath{k}}\fi}
\newunicodechar{ᵖ}{\ifmmode^{p}\else\textsuperscript{\ensuremath{p}}\fi}
\newunicodechar{ᴺ}{\ifmmode^{N}\else\textsuperscript{\ensuremath{N}}\fi}
\newunicodechar{ᶜ}{\ifmmode^{c}\else\textsuperscript{\ensuremath{c}}\fi}
\newunicodechar{ʰ}{\ifmmode^{h}\else\textsuperscript{\ensuremath{h}}\fi}
\newunicodechar{ᵠ}{\ifmmode^{q}\else\textsuperscript{\ensuremath{q}}\fi}
\newunicodechar{ₙ}{\ifmmode_{n}\else\textsubscript{\ensuremath{n}}\fi}
\newunicodechar{ₐ}{\ifmmode_{a}\else\textsubscript{\ensuremath{a}}\fi}
\newunicodechar{ᵢ}{\ifmmode_{i}\else\textsubscript{\ensuremath{i}}\fi}
\newunicodechar{ᵣ}{\ifmmode_{r}\else\textsubscript{\ensuremath{r}}\fi}
\newunicodechar{ₖ}{\ifmmode_{k}\else\textsubscript{\ensuremath{k}}\fi}
\newunicodechar{ᵦ}{\ifmmode_{\beta}\else\textsubscript{\ensuremath{\beta}}\fi}
\newunicodechar{ₓ}{\ifmmode_{x}\else\textsubscript{\ensuremath{x}}\fi}
\newfontfamily\popdisplay{ArchivoBlack-Regular.ttf}[Path=../../../fonts/]
\newfontfamily\poplogo{SpaceGrotesk-Bold.ttf}[Path=../../../fonts/]
\newfontfamily\popsemi{PlusJakartaSans-SemiBold.ttf}[Path=../../../fonts/]
\newfontfamily\popxbold{PlusJakartaSans-ExtraBold.ttf}[Path=../../../fonts/]
\newfontfamily\popxboldls{PlusJakartaSans-ExtraBold.ttf}[Path=../../../fonts/,LetterSpace=3.0]

\setlist[enumerate]{leftmargin=1.7em,itemsep=5pt,topsep=4pt}
\setlist[itemize]{leftmargin=1.7em,itemsep=4pt,topsep=4pt,label={\color{poporange}\textbullet}}
\setlength{\parindent}{0pt}
\setlength{\parskip}{5pt}
\renewcommand{\arraystretch}{1.25}

% pgfplots : style Pop par défaut
\pgfplotsset{
  every axis/.append style={axis line style={popink,line width=1pt},tick style={popink},
    grid style={popline,line width=0.5pt},label style={font=\small},tick label style={font=\footnotesize}},
  popcurve/.style={poporange,line width=1.8pt,no markers,smooth},
}
% Même style disponible dans un \draw TikZ simple (hors pgfplots)
\tikzset{popcurve/.style={poporange,line width=1.8pt}}
\tikzset{popgrid/.style={popline,very thin}}
\colorlet{popcurve}{poporange} % le nom du style est parfois utilisé comme couleur

% ============================================================
% Pill / squiggle
% ============================================================
\newcommand{\poppill}[3]{%
  \tikz[baseline=-0.55ex]\node[draw=popink,line width=1.2pt,rounded corners=2.6pt,%
    fill=#2,inner xsep=6.5pt,inner ysep=3pt]{%
    \popxboldls\fontsize{7}{7}\selectfont\color{#3}\MakeUppercase{#1}};%
}
\newcommand{\popsquiggle}[1]{%
  \tikz[baseline=-0.4ex]{
    \draw[#1,line width=2.6pt,line cap=round]
      plot [smooth] coordinates {(0,0.09) (0.34,0.01) (0.68,0.21) (1.02,0.01) (1.36,0.10)};
  }%
}
% Mot-clé surligné (marqueur orange sous le texte)
\newcommand{\cle}[1]{{\popxbold\color{popdark}#1}}

% ============================================================
% En-tête / pied
% ============================================================
\pagestyle{fancy}
\fancyhf{}
\renewcommand{\headrulewidth}{0pt}
\renewcommand{\footrulewidth}{0pt}
\lhead{\raisebox{-0.15ex}{\poplogo\fontsize{13}{13}\selectfont\color{popink}OnBuch\color{poporange}+}}
\rhead{\poppill{\DOCMATIERE\ \textbullet\ \DOCNIVEAU\ \textbullet\ Leçon \DOCLECON}{white}{popink}}
\lfoot{\popxbold\fontsize{7.6}{7.6}\selectfont\color{popmuted}\MakeUppercase{Cours signé OnBuch+}\ \textbullet\ {\popsemi\DOCTITRE}}
\rfoot{\tikz[baseline=-0.6ex]\node[rounded corners=8.5pt,fill=popink,minimum height=17pt,minimum width=17pt,inner xsep=5pt,inner sep=0]{\popxbold\fontsize{8}{8}\selectfont\color{popcream}\thepage\,/\,\pageref*{LastPage}};}

% ============================================================
% Titres
% ============================================================
\newcommand{\chaptitre}[2]{%
  \begin{tcolorbox}[enhanced,colback=poporange,colframe=popink,boxrule=2.6pt,arc=11pt,
      drop shadow=popink,left=15pt,right=15pt,top=12pt,bottom=11pt,before skip=3pt,after skip=18pt]
    \poppill{#2}{popink}{popcream}\par\vspace{9pt}
    {\raggedright\hyphenpenalty=10000\exhyphenpenalty=10000\popdisplay\fontsize{22}{24}\selectfont\color{popcream}\MakeUppercase{#1}\par}
  \end{tcolorbox}
}
% Section numérotée : pastille orange avec le numéro + titre Archivo
\newcounter{popsec}
\newcommand{\coursec}[1]{%
  \stepcounter{popsec}\setcounter{popsub}{0}%
  \par\vspace{16pt}\noindent
  \begin{minipage}{\linewidth}
  \tikz[baseline=-0.7ex]\node[draw=popink,line width=1.6pt,fill=poporange,rounded corners=4pt,minimum size=22pt,inner sep=0]{\popdisplay\fontsize{12}{12}\selectfont\color{popcream}\thepopsec};%
  \hspace{8pt}{\popdisplay\fontsize{15}{17}\selectfont\color{popink}\MakeUppercase{#1}}\par
  \vspace{4pt}\noindent{\color{poporange}\rule{36pt}{3.6pt}}
  \end{minipage}\par\nopagebreak\vspace{8pt}
}
\newcounter{popsub}
\newcommand{\courssub}[1]{%
  \stepcounter{popsub}%
  \par\vspace{9pt}\noindent{\popxbold\fontsize{11.5}{13}\selectfont\color{popink}{\color{poporange}\thepopsec.\thepopsub}\hspace{6pt}#1}\par\nopagebreak\vspace{4pt}
}

% ============================================================
% Boîtes pédagogiques — même recette : fond blanc, bordure épaisse colorée,
% ombre dure encre, badge plein. Titre optionnel [#1].
% ============================================================
\tcbset{popbase/.style={breakable,enhanced,colback=white,boxrule=2.3pt,arc=9pt,
  shadow={3pt}{-3pt}{0pt}{popink},left=14pt,right=14pt,top=11pt,bottom=11pt,before skip=11pt,after skip=15pt,
  fontupper=\popsemi\fontsize{10.6}{14.5}\selectfont\color{popink},
  fontlower=\popsemi\fontsize{10.6}{14.5}\selectfont\color{popink}}}
% Coche dessinée (le glyphe ✓ n'existe pas dans les polices du document)
\newcommand{\popcheck}[1]{\tikz[baseline=-0.5ex]\draw[#1,line width=1.6pt,line cap=round,line join=round](-0.13,0.0)--(-0.04,-0.09)--(0.13,0.1);}
\providecommand{\coche}{\popcheck{popgreen}}
\providecommand{\resultat}[1]{\par\smallskip\noindent\fcolorbox{popgreen}{popgreenL}{\parbox{\dimexpr\linewidth-2\fboxsep-2\fboxrule\relax}{\textbf{Résultat.} #1}}\par\smallskip}
\newcommand{\popboxhead}[3]{\poppill{#1}{#2}{white}\ifx\relax#3\relax\else\hspace{6pt}{\popxbold\fontsize{10.6}{12}\selectfont\color{popink}#3}\fi\par\vspace{7pt}}

\newtcolorbox{aretenir}[1][]{popbase,colframe=popgreen,before upper={\popboxhead{À retenir}{popgreen}{#1}}}
% Texte en allemand (document de lecture, dialogue, extrait) : cadre
% d'étude. Un titre optionnel (source, auteur) s'affiche dans l'en-tête.
\newtcolorbox{textebox}[1][]{popbase,colframe=popblue,before upper={\popboxhead{Text}{popblue}{#1}\begin{otherlanguage}{ngerman}},after upper={\end{otherlanguage}}}
% Marques ✓ / ✗ (forme correcte / fautive) souvent utilisées par les modèles
\providecommand{\cross}{\textcolor{poppink}{\ensuremath{\times}}}
\providecommand{\xmark}{\cross}\providecommand{\crossmark}{\cross}\providecommand{\cmark}{\ensuremath{\checkmark}}
% Ligne de réponse / justification à compléter (inventée par les modèles)
\providecommand{\justification}[1]{\par\noindent\textit{Justification~:} \dotfill\par}
% \textipa (package tipa non chargé) : la police ne couvre pas l'API -> simple passage
\providecommand{\textipa}[1]{#1}
% Soulignement (ulem non chargé) : \uline, \ul
\providecommand{\uline}[1]{\underline{#1}}\providecommand{\ul}[1]{\underline{#1}}
% \glossaire{\item ...} inventée par les modèles : liste de glossaire
\providecommand{\glossaire}[1]{\begin{itemize}#1\end{itemize}}
% \alert (beamer) inventée par les modèles : texte mis en évidence
\providecommand{\alert}[1]{\textbf{\textcolor{poppink}{#1}}}
% variantes ASCII d'ordinaux écrites en macro
\providecommand{\iere}{\textsuperscript{re}}\providecommand{\ieme}{\textsuperscript{e}}\providecommand{\ere}{\textsuperscript{re}}\providecommand{\eme}{\textsuperscript{e}}\providecommand{\er}{\textsuperscript{er}}\providecommand{\ie}{\textsuperscript{e}}
% Ordinaux français écrits en macro par les modèles : 1\ière, 3\ième
\newcommand{\ière}{\textsuperscript{re}}\newcommand{\ième}{\textsuperscript{e}}
% \exemple (inventée par les modèles) : étiquette « Exemple : »
\providecommand{\exemple}{\textbf{Exemple~:}\ }
% \square utilisable aussi hors mode math (cases à cocher)
\AtBeginDocument{\let\squareMath\square\renewcommand{\square}{\ensuremath{\squareMath}}}
% Mot ou phrase en allemand à l'intérieur d'un paragraphe français
\newcommand{\all}[1]{\ifmmode\text{\textit{#1}}\else\begingroup\selectlanguage{ngerman}\itshape #1\endgroup\fi}
\let\esp\all % alias toléré
\newtcolorbox{exemplebox}[1][]{popbase,colframe=poporange,before upper={\popboxhead{Exemple}{poporange}{#1}}}
\newtcolorbox{definition}[1][]{popbase,colframe=popblue,before upper={\popboxhead{Définition}{popblue}{#1}}}
\newtcolorbox{propriete}[1][]{popbase,colframe=poppurple,before upper={\popboxhead{Propriété}{poppurple}{#1}}}
\newtcolorbox{methode}[1][]{popbase,colframe=popgold,before upper={\popboxhead{Méthode}{popgold}{#1}}}
\newtcolorbox{attention}[1][]{popbase,colframe=poppink,before upper={\popboxhead{Attention}{poppink}{#1}}}
\newtcolorbox{experience}[1][]{popbase,colframe=popblue,colback=popblueL,before upper={\popboxhead{Expérience}{popblue}{#1}}}
% « Le savais-tu ? » : carte violette pleine (contexte, culture, Cameroun)
\newtcolorbox{savaistu}[1][]{popbase,colback=poppurple,colframe=popink,
  fontupper=\popsemi\fontsize{10.4}{14.2}\selectfont\color{white},
  before upper={\poppill{Le savais-tu\,?}{white}{poppurple}\ifx\relax#1\relax\else\hspace{6pt}{\popxbold\color{white}#1}\fi\par\vspace{7pt}}}
% Exercice résolu : énoncé en haut, \tcblower puis la solution sur fond crème
\newcounter{popexo}\newcounter{popexr}
\newtcolorbox{exoresolu}[1][]{popbase,colframe=poporange,colbacklower=popcream,
  segmentation style={popink,line width=1.2pt,dashed},
  before upper={\stepcounter{popexr}\popboxhead{Exercice résolu \thepopexr}{poporange}{#1}},
  before lower={\poppill{Solution}{popink}{popcream}\par\vspace{6pt}}}
% Exercice (non résolu en place) — la correction est dans la partie Corrigés
\newtcolorbox{exercice}[1][]{popbase,colframe=popink,
  before upper={\stepcounter{popexo}\popboxhead{Exercice \thepopexo}{popink}{#1}}}
% Corrigé
\newtcolorbox{corrige}[1][]{popbase,colframe=popgreen,colback=popgreenL,
  before upper={\popboxhead{Corrigé}{popgreen}{#1}}}
% Objectifs / prérequis / bilan
\newtcolorbox{objectifs}{popbase,colframe=popink,colback=poporangeL,
  before upper={\popboxhead{Objectifs de la leçon}{popink}{}}}
\newtcolorbox{prerequis}{popbase,colframe=popink,colback=popblueL,
  before upper={\popboxhead{Prérequis}{popblue}{}}}
\newtcolorbox{bilan}{popbase,colframe=popink,colback=white,boxrule=2.8pt,
  before upper={\popboxhead{Fiche bilan}{poporange}{L'essentiel à connaître pour l'examen}}}
% Figure encadrée avec légende
\newtcolorbox{popfigure}[1][]{enhanced,breakable=false,colback=white,colframe=popline,boxrule=1.4pt,arc=8pt,
  left=10pt,right=10pt,top=10pt,bottom=8pt,before skip=10pt,after skip=12pt,halign=center,
  lower separated=false,halign lower=center,
  fontlower=\popsemi\fontsize{9}{11.5}\selectfont\color{popmuted},
  #1}
\newcommand{\legende}[1]{\par\vspace{6pt}{\popsemi\fontsize{9}{11.5}\selectfont\color{popmuted}#1}}

% ============================================================
% Signature OnBuch+
% ============================================================
\newcommand{\poplogoinline}[1]{{\poplogo\fontsize{#1}{#1}\selectfont\color{popink}OnBuch\color{poporange}+}}
\newcommand{\signatureonbuch}{%
  \begin{tcolorbox}[enhanced,colback=popink,colframe=popink,boxrule=2.2pt,arc=10pt,
      drop shadow=poporange,left=14pt,right=14pt,top=10pt,bottom=10pt,before skip=0pt,after skip=0pt]
    \tikz[baseline=-0.7ex]\node[circle,fill=popgreen,draw=popcream,line width=1.3pt,minimum size=22pt,inner sep=0]{\popcheck{white}};%
    \hspace{9pt}\begin{minipage}[c]{0.62\linewidth}
      {\popxbold\fontsize{12}{13}\selectfont\color{popcream}Signé\ \textbullet\ {\poplogo OnBuch\color{poporange}+}}\par\vspace{2pt}
      {\raggedright\popsemi\fontsize{8}{10.5}\selectfont\color{popline}\MakeUppercase{Équipe pédagogique OnBuch+}\\\MakeUppercase{Conforme au programme officiel MINESEC}\par}
    \end{minipage}\hfill
    {\poplogo\fontsize{22}{22}\selectfont\color{popcream}OnBuch\color{poporange}+}
  \end{tcolorbox}%
}

% ============================================================
% Page de garde
% #1 titre · #2 matière · #3 niveau · #4 module · #5 n° leçon · #6 durée ·
% #7 description / objectifs (texte ou itemize)
% ============================================================
\newcommand{\pagegarde}[7]{%
  \thispagestyle{empty}
  \newgeometry{a4paper,margin=1.7cm,top=1.6cm,bottom=1.4cm}%
  \begin{tikzpicture}[remember picture,overlay]
    \fill[poporange!18] ([xshift=-3.2cm,yshift=-3.6cm]current page.north east) circle (4.6cm);
    \fill[poppurple!14] ([xshift=2.4cm,yshift=7.6cm]current page.south west) circle (3.8cm);
  \end{tikzpicture}%
  {\poplogo\fontsize{30}{30}\selectfont\color{popink}OnBuch\color{poporange}+}\hfill
  \raisebox{0.6ex}{\poppill{Cours signé \textbullet\ Premium}{popink}{popcream}}\par
  \vspace{1pt}\popsquiggle{poppurple}\par
  \vspace{8pt}
  \begin{tcolorbox}[enhanced,colback=poporange,colframe=popink,boxrule=2.8pt,arc=13pt,
      drop shadow=popink,left=18pt,right=18pt,top=12pt,bottom=13pt,after skip=10pt]
    \poppill{#2 \textbullet\ #3}{popink}{popcream}\hspace{5pt}\poppill{#4}{white}{popink}\par\vspace{8pt}
    {\popdisplay\fontsize{14}{14}\selectfont\color{popink}LEÇON #5}\par\vspace{4pt}
    {\raggedright\hyphenpenalty=10000\exhyphenpenalty=10000\popdisplay\fontsize{25}{27}\selectfont\color{popcream}\MakeUppercase{#1}\par}
    \vspace{8pt}
    {\popxbold\fontsize{9.5}{9.5}\selectfont\color{popink}\MakeUppercase{Durée conseillée : #6}}
  \end{tcolorbox}
  \begin{tcolorbox}[enhanced,colback=white,colframe=popink,boxrule=2.2pt,arc=10pt,
      drop shadow=popink,left=16pt,right=16pt,top=10pt,bottom=10pt,after skip=8pt]
    \poppill{Dans cette leçon}{poppurple}{white}\par\vspace{5pt}
    \popsemi\fontsize{9.3}{12.6}\selectfont\color{popink}#7
  \end{tcolorbox}
  \noindent\begin{minipage}[t]{0.487\linewidth}
    \begin{tcolorbox}[enhanced,colback=popgreen,colframe=popink,boxrule=2.2pt,arc=10pt,
        drop shadow=popink,left=11pt,right=11pt,top=10pt,bottom=10pt,height=3.1cm,valign=center]
      \tcbox[colback=white,colframe=popink,boxrule=1.3pt,arc=5pt,boxsep=0pt,nobeforeafter,
        left=3pt,right=3pt,top=3pt,bottom=3pt,valign=center]{\qrcode[height=1.9cm]{https://play.google.com/store/apps/details?id=com.luvvix.onbuch&hl=fr}}%
      \hspace{8pt}\begin{minipage}[c]{0.5\linewidth}
        {\popdisplay\fontsize{11}{12.5}\selectfont\color{white}\MakeUppercase{Télécharge OnBuch}}\par\vspace{4pt}
        {\raggedright\popsemi\fontsize{8.4}{11}\selectfont\color{white}Gratuit \textbullet\ Google Play\\Tuteur IA, annales, concours, résultats\par}
      \end{minipage}
    \end{tcolorbox}
  \end{minipage}\hfill
  \begin{minipage}[t]{0.487\linewidth}
    \begin{tcolorbox}[enhanced,colback=poppurple,colframe=popink,boxrule=2.2pt,arc=10pt,
        drop shadow=popink,left=11pt,right=11pt,top=10pt,bottom=10pt,height=3.1cm,valign=center]
      \tikz[baseline=-0.7ex]\node[circle,fill=white,draw=popink,line width=1.3pt,minimum size=1.6cm,inner sep=0]{\popdisplay\fontsize{24}{24}\selectfont\color{poppurple}+};%
      \hspace{8pt}\begin{minipage}[c]{0.6\linewidth}
        {\popdisplay\fontsize{11}{12.5}\selectfont\color{white}\MakeUppercase{Passe à OnBuch+}}\par\vspace{4pt}
        {\raggedright\popsemi\fontsize{8.4}{11}\selectfont\color{white}Tous les cours signés, Tuteur IA illimité, fiches PDF — onglet OnBuch+ de l'app\par}
      \end{minipage}
    \end{tcolorbox}
  \end{minipage}
  \vspace{8pt}
  \popcontenu
  \vfill
  \signatureonbuch
  \restoregeometry
  \newpage
}

% Rangée « Ce cours contient » (page de garde)
\newcommand{\poptile}[3]{% #1 couleur #2 gros texte #3 légende
  \begin{tcolorbox}[enhanced,colback=#1,colframe=popink,boxrule=2pt,arc=9pt,shadow={2.5pt}{-2.5pt}{0pt}{popink},
      left=6pt,right=6pt,top=9pt,bottom=9pt,halign=center,valign=center,height=1.75cm,width=\linewidth,nobeforeafter]
    {\popdisplay\fontsize{12.5}{13}\selectfont\color{white}\MakeUppercase{#2}}\par\vspace{3pt}
    {\centering\popsemi\fontsize{7.8}{9.5}\selectfont\color{white}#3\par}
  \end{tcolorbox}}
\newcommand{\popcontenu}{%
  \par\noindent{\popxboldls\fontsize{7.5}{7.5}\selectfont\color{popmuted}CE COURS SIGNÉ CONTIENT}\par\vspace{7pt}
  \noindent\begin{minipage}[t]{0.235\linewidth}\poptile{popblue}{Cours}{complet, illustré, pas à pas}\end{minipage}\hfill
  \begin{minipage}[t]{0.235\linewidth}\poptile{poporange}{Méthodes}{et exercices résolus}\end{minipage}\hfill
  \begin{minipage}[t]{0.235\linewidth}\poptile{poppink}{Exercices}{type examen + corrigés}\end{minipage}\hfill
  \begin{minipage}[t]{0.235\linewidth}\poptile{popgreen}{Fiche bilan}{l'essentiel à retenir}\end{minipage}\par}

% Sommaire
\newtcolorbox{sommaire}{enhanced,colback=white,colframe=popink,boxrule=2.2pt,arc=10pt,
  drop shadow=popink,left=18pt,right=18pt,top=13pt,bottom=13pt,before skip=6pt,after skip=20pt,
  before upper={\poppill{Sommaire}{poppurple}{white}\par\vspace{9pt}\popsemi\fontsize{10.6}{14.5}\selectfont\color{popink}}}

% Signature de fin de cours — poussée tout en bas de la dernière page
\newcommand{\findecours}{%
  \par\vfill\nopagebreak
  \begin{center}\popsquiggle{poporange}\end{center}\vspace{4pt}
  \signatureonbuch
}
\AtBeginDocument{\def\llbracket{\mathopen{[\mkern-3.5mu[}}\def\rrbracket{\mathclose{]\mkern-3.5mu]}}} % intervalles d'entiers (glyphe ⟦ absent de la police)
% Glyphes absents de Fira Math : redéfinis à partir de glyphes disponibles
\AtBeginDocument{%
  \def\setminus{\mathbin{\backslash}}\let\smallsetminus\setminus
  \def\vdots{\mathord{\vbox{\baselineskip3.2pt\lineskiplimit0pt\kern4pt\hbox{.}\hbox{.}\hbox{.}}}}%
  \def\ddots{\mathinner{\mkern1mu\raise6.5pt\hbox{.}\mkern2mu\raise3.6pt\hbox{.}\mkern2mu\raise0.7pt\hbox{.}\mkern1mu}}%
  \def\triangle{\mathord{\scalebox{0.9}{$\symup{\Delta}$}}}%
  \def\nabla{\mathord{\raisebox{\depth}{\scalebox{1}[-1]{$\symup{\Delta}$}}}}%
  \def\checkmark{\popcheck{popdark}}%
  \def\star{\mathbin{\ast}}\def\bigstar{\mathord{\ast}}%
  \def\diamond{\mathbin{\scalebox{0.6}{$\Diamond$}}}%
  \def\bigcup{\mathop{\mathchoice{\raisebox{-0.3em}{\scalebox{1.7}{$\cup$}}}{\scalebox{1.25}{$\cup$}}{\cup}{\cup}}\displaylimits}%
  \def\bigcap{\mathop{\mathchoice{\raisebox{-0.3em}{\scalebox{1.7}{$\cap$}}}{\scalebox{1.25}{$\cap$}}{\cap}{\cap}}\displaylimits}%
  \def\lesseqqgtr{\mathrel{\lessgtr}}\def\bigoplus{\mathop{\oplus}\displaylimits}%
}
\providecommand{\ssi}{si et seulement si} % abréviation fréquente
\AtBeginDocument{\providecommand{\wideparen}{\overparen}} % arc de cercle

\begin{document}

\pagegarde{Le néocolonialisme}{Allemand}{Tle A}{Chapitre 5 : Citoyenneté}{AL27}{2 h}{Lecture et commentaire d'un texte allemand sur le néocolonialisme, enrichissement du vocabulaire politique et économique (Kolonialismus, Abhängigkeit, Ausbeutung, Rohstoffe), et étude grammaticale de la nominalisation et du génitif avec ses prépositions. La leçon se termine par un commentaire guidé et une production écrite courte de type Bac.
\vspace{4pt}
\begin{multicols}{2}\small
\begin{itemize}
\item À la fin de cette leçon, je sais comprendre globalement puis en détail un texte allemand sur le néocolonialisme.
\item À la fin de cette leçon, je sais définir et employer correctement le vocabulaire thématique : die Kolonie, der Kolonialismus, die Abhängigkeit, die Ausbeutung, die Unabhängigkeit, der Rohstoff.
\item À la fin de cette leçon, je sais former et utiliser les familles de mots du thème (abhängig / die Abhängigkeit, ausbeuten / die Ausbeutung, unabhängig / die Unabhängigkeit).
\item À la fin de cette leçon, je sais transformer une phrase verbale en construction nominale (nominalisation) et inversement.
\item À la fin de cette leçon, je sais utiliser le génitif et les prépositions suivies du génitif (wegen, trotz, während, innerhalb, aufgrund).
\item À la fin de cette leçon, je sais répondre par écrit et oralement à des questions de compréhension sur un texte.
\end{itemize}\end{multicols}}

\begin{sommaire}
\begin{multicols}{2}
\begin{enumerate}
\item Texte support : Der Neokolonialismus in Afrika
\item Compréhension du texte
\item Lexique thématique : Kolonialismus und Abhängigkeit
\item Grammaire 1 : la nominalisation (die Nominalisierung)
\item Grammaire 2 : le génitif et les prépositions avec le génitif
\item Commentaire guidé et production
\item Activité d'intégration
\item Méthodes et exercices résolus
\item Exercices
\item Corrigés des exercices
\item Fiche bilan
\end{enumerate}
\end{multicols}
\end{sommaire}

\begin{objectifs}
\begin{itemize}
\item À la fin de cette leçon, je sais comprendre globalement puis en détail un texte allemand sur le néocolonialisme.
\item À la fin de cette leçon, je sais définir et employer correctement le vocabulaire thématique : die Kolonie, der Kolonialismus, die Abhängigkeit, die Ausbeutung, die Unabhängigkeit, der Rohstoff.
\item À la fin de cette leçon, je sais former et utiliser les familles de mots du thème (abhängig / die Abhängigkeit, ausbeuten / die Ausbeutung, unabhängig / die Unabhängigkeit).
\item À la fin de cette leçon, je sais transformer une phrase verbale en construction nominale (nominalisation) et inversement.
\item À la fin de cette leçon, je sais utiliser le génitif et les prépositions suivies du génitif (wegen, trotz, während, innerhalb, aufgrund).
\item À la fin de cette leçon, je sais répondre par écrit et oralement à des questions de compréhension sur un texte.
\item À la fin de cette leçon, je sais résumer un texte en quelques phrases avec des connecteurs adaptés.
\item À la fin de cette leçon, je sais rédiger un court paragraphe d'opinion sur le néocolonialisme en réutilisant le lexique et la grammaire de la leçon.
\end{itemize}
\end{objectifs}

\begin{prerequis}
\begin{itemize}
\item Les articles et les pluriels des noms (der/die/das, pluriels en -e, -en, -er) vus en Première.
\item Le vocabulaire de base de l'histoire et de la politique (der Staat, die Regierung, die Wirtschaft, das Land).
\item La déclinaison de l'adjectif et des articles aux quatre cas, notamment le génitif simple (des Vaters, der Mutter).
\item Les connecteurs logiques de base (weil, denn, deshalb, aber, obwohl).
\item La lecture de textes documentaires et la réponse à des questions de compréhension (méthode vue en Première).
\end{itemize}
\end{prerequis}

\newpage

\coursec{Texte support : Der Neokolonialismus in Afrika}

\courssub{Situation d'accroche et problématique}
Imaginez : le Cameroun est indépendant depuis 1960, mais une grande partie de son cacao, de son pétrole et de son bois est achetée et transformée par des entreprises étrangères, souvent à des prix fixés ailleurs. En Allemagne, on débat aussi de ce sujet : des chocolatiers suisses et allemands dépendent du cacao africain. Peut-on parler d'une vraie indépendance quand l'économie reste contrôlée de l'extérieur ? C'est exactement ce que les politologues appellent \all{der Neokolonialismus}. Dans cette leçon, nous allons lire un texte allemand sur ce phénomène et apprendre à en parler avec un vocabulaire précis.

\courssub{Lecture globale du texte}
\begin{methode}[Lire un texte en trois étapes]
\begin{enumerate}
\item \textbf{Titre et première phrase :} identifiez le thème général et la thèse principale.
\item \textbf{Mots-clés :} soulignez les noms récurrents (souvent des nominalisations) et les connecteurs logiques (\all{deshalb}, \all{wegen}, \all{während}).
\item \textbf{Relecture avec glossaire :} lisez une deuxième fois en vous aidant du glossaire pour saisir les détails et les nuances.
\end{enumerate}
\end{methode}

Appliquez cette méthode au texte ci-dessous. Lisez-le une première fois en silence, puis écoutez la lecture à voix haute de l'enseignant. Repérez : le thème, la thèse (l'argument principal) et le ton (informatif, critique, analytique).

\begin{textebox}[Der Neokolonialismus in Afrika]
Viele afrikanische Länder sind seit den 1960er Jahren politisch unabhängig. Trotz der Unabhängigkeit bleiben sie wirtschaftlich von den früheren Kolonialmächten abhängig. Man spricht vom Neokolonialismus: Die Ausbeutung der Rohstoffe geht weiter, weil die Preise auf den Weltmärkten im Norden festgelegt werden. Wegen dieser Abhängigkeit können viele Staaten ihre eigene Industrie kaum entwickeln. Während der Kolonialzeit herrschte offene Gewalt; heute herrscht die Abhängigkeit durch Schulden, Handelsverträge und internationale Konzerne. Aufgrund dieser Strukturen bleibt der Reichtum der Rohstoffe oft im Ausland. Eine echte Unabhängigkeit verlangt deshalb eine faire Zusammenarbeit und die Verarbeitung der Rohstoffe im eigenen Land.
\end{textebox}

\courssub{Glossaire en marge}
Voici les mots difficiles du texte, avec leur article, leur pluriel et leur traduction. Apprenez-les par cœur ; ils constituent le socle lexical du thème.

\begin{center}
\begin{tabularx}{\linewidth}{|X|X|X|}
\hline
\rowcolor{popblueL}
\textbf{Allemand (article + pluriel)} & \textbf{Français} & \textbf{Classe / Note} \\
\hline
\all{die Kolonialmacht, die -mächte} & la puissance coloniale & nom féminin \\
\hline
\all{festlegen (legt fest, hat festgelegt)} & fixer, déterminer & verbe séparable \\
\hline
\all{die Schulden (Plural only)} & les dettes & nom féminin (pluriel seulement) \\
\hline
\all{der Handelsvertrag, die -verträge} & l'accord commercial & nom masculin \\
\hline
\all{der Konzern, die -zerne} & la multinationale, le groupe & nom masculin \\
\hline
\all{die Verarbeitung, die -en} & la transformation (industrielle) & nom féminin \\
\hline
\all{fair} & équitable, juste & adjectif (invariable) \\
\hline
\end{tabularx}
\end{center}

\courssub{Repérage du vocabulaire du programme dans le texte}
Le programme exige la maîtrise de cinq noms clés. Voici où ils apparaissent dans le texte, avec leur article, leur pluriel et leur traduction.

\begin{center}
\begin{tabularx}{\linewidth}{|X|X|X|X|}
\hline
\rowcolor{popblueL}
\textbf{Mot du programme} & \textbf{Forme dans le texte} & \textbf{Article + Pluriel} & \textbf{Traduction} \\
\hline
\cle{die Abhängigkeit} & \all{abhängig} (adj.), \all{Abhängigkeit} (nom) & \all{die Abhängigkeit, die -en} & la dépendance \\
\hline
\cle{die Ausbeutung} & \all{die Ausbeutung} & \all{die Ausbeutung, die -en} & l'exploitation \\
\hline
\cle{die Unabhängigkeit} & \all{unabhängig} (adj.), \all{Unabhängigkeit} (nom) & \all{die Unabhängigkeit, die -en} & l'indépendance \\
\hline
\cle{der Rohstoff} & \all{die Rohstoffe} (pluriel) & \all{der Rohstoff, die -e} & la matière première \\
\hline
\cle{der Kolonialismus} & \all{Neokolonialismus} (composé) & \all{der Kolonialismus} (pas de pluriel usuel) & le colonialisme \\
\hline
\end{tabularx}
\end{center}

\begin{aretenir}[Familles de mots utiles]
\begin{itemize}
\item \all{abhängig} (adj.) $\rightarrow$ \all{die Abhängigkeit} (nom), \all{abhängig machen} (rendre dépendant).
\item \all{ausbeuten} (verbe) $\rightarrow$ \all{die Ausbeutung} (nom), \all{der Ausbeuter} (l'exploiteur).
\item \all{unabhängig} (adj.) $\rightarrow$ \all{die Unabhängigkeit} (nom), \all{Unabhängigkeitserklärung} (déclaration d'indépendance).
\item \all{der Rohstoff} $\rightarrow$ \all{roh} (adj. = brut), \all{verarbeiten} (transformer), \all{die Verarbeitung}.
\item \all{kolonial} (adj.) $\rightarrow$ \all{der Kolonialismus}, \all{die Kolonialmacht}, \all{kolonisieren} (coloniser).
\end{itemize}
\end{aretenir}



\courssub{Repérage des exemples de grammaire dans le texte}
Observez deux phénomènes grammaticaux centraux de cette leçon : la \textbf{nominalisation} (transformation verbes/adjectifs $\rightarrow$ noms) et les \textbf{prépositions régissant le génitif}.

\begin{exemplebox}[Nominalisations repérées]
\begin{itemize}
\item \all{die Ausbeutung der Rohstoffe} (l'exploitation des matières premières) : nominalisation du verbe \all{ausbeuten} + génitif de l'objet.
\item \all{die Verarbeitung der Rohstoffe} (la transformation des matières premières) : nominalisation du verbe \all{verarbeiten} + génitif de l'objet.
\item \all{die Unabhängigkeit} (l'indépendance) : nominalisation de l'adjectif \all{unabhängig} (suffixe \all{-keit}).
\item \all{die Abhängigkeit} (la dépendance) : nominalisation de l'adjectif \all{abhängig} (suffixe \all{-keit}).
\end{itemize}
\end{exemplebox}

\begin{exemplebox}[Prépositions + génitif repérées]
\begin{itemize}
\item \all{Trotz der Unabhängigkeit} (malgré l'indépendance) : \all{trotz} + génitif.
\item \all{Wegen dieser Abhängigkeit} (à cause de cette dépendance) : \all{wegen} + génitif.
\item \all{Während der Kolonialzeit} (pendant l'époque coloniale) : \all{während} + génitif.
\item \all{Aufgrund dieser Strukturen} (en raison de ces structures) : \all{aufgrund} + génitif.
\end{itemize}
\end{exemplebox}

\begin{attention}[Piège fréquent pour francophones]
En français, « malgré », « à cause de », « pendant », « en raison de » sont suivis d'un nom sans marque de cas (ou d'une proposition). En allemand, ces prépositions \textbf{exigent le génitif} (surtout à l'écrit soutenu). Ne les confondez pas avec les prépositions à datif (\all{mit}, \all{von}, \all{zu}) ou à accusatif (\all{durch}, \all{für}, \all{gegen}). Apprenez-les par cœur : \cle{trotz, wegen, während, aufgrund, statt, anstatt, innerhalb, außerhalb, jenseits, diesseits}.
\end{attention}

\courssub{Définition du concept clé}
\begin{definition}[der Neokolonialismus]
\all{der Neokolonialismus} = die wirtschaftliche und politische Abhängigkeit eines offiziell unabhängigen Landes von früheren Kolonialmächten oder ausländischen Konzernen.
\end{definition}

\courssub{Le saviez-vous ?}
\begin{savaistu}[Kwame Nkrumah et le terme « Neokolonialismus »]
Le terme \all{Neokolonialismus} a été popularisé par \textbf{Kwame Nkrumah} (1909–1972), premier président du Ghana, dans son ouvrage \emph{Neo-Colonialism: The Last Stage of Imperialism} (1965). Il y décrit comment les anciennes puissances coloniales maintiennent leur emprise par des moyens économiques (dettes, contrôle des prix des matières premières) plutôt que militaires. Ce concept est encore aujourd'hui au cœur des débats dans les universités allemandes (ex. \all{Entwicklungssoziologie}, \all{Postkoloniale Studien}) et dans la presse (ex. \all{Die Zeit}, \all{Der Spiegel}) lorsqu'il s'agit des relations Europe–Afrique.
\end{savaistu}

\courssub{Carte mentale lexicale : le champ du néocolonialisme}
La figure ci-dessous résume les liens sémantiques entre les notions du texte. Utilisez-la pour réviser le vocabulaire et préparer une prise de parole.

\begin{popfigure}
\begin{tikzpicture}[
  mindmap,
  concept color=popblue,
  level 1 concept/.append style={level distance=3.5cm, sibling angle=72},
  level 2 concept/.append style={level distance=2.5cm},
  every node/.style={font=\small},
  grow cyclic
]
\node[concept] {Neokolonialismus}
  child[concept color=poporange] {node[concept, align=center]{politische\\Unabhängigkeit}}
  child[concept color=popgreen] {node[concept, align=center]{wirtschaftliche\\Abhängigkeit}}
  child[concept color=poppurple] {node[concept, align=center]{Ausbeutung der\\Rohstoffe}}
  child[concept color=poppink] {node[concept, align=center]{Schulden \&\\Handelsverträge}}
  child[concept color=popgold] {node[concept, align=center]{internationale\\Konzerne}};
\end{tikzpicture}
\legende{Carte mentale : champs lexicaux autour du néocolonialisme (vocabulaire du programme en gras).}
\end{popfigure}

\courssub{Exercice résolu : compréhension globale}
\begin{exoresolu}[Question 1 : thèse et ton du texte]
\textbf{Énoncé :} Lisez le texte \emph{Der Neokolonialismus in Afrika}. Quelle est la thèse principale ? Quel est le ton adopté ? Justifiez par deux éléments du texte.

\tcblower
\textbf{Solution détaillée :}
\begin{enumerate}
\item \textbf{Thèse :} L'indépendance politique (depuis les années 1960) n'a pas entraîné l'indépendance économique ; les pays africains restent dépendants par le commerce, la dette et les multinationales — c'est le néocolonialisme.
\item \textbf{Ton :} Analytique et critique. Le texte constate des faits (\all{bleiben sie wirtschaftlich abhängig}) et dénonce une structure injuste (\all{Preise ... im Norden festgelegt}, \all{Reichtum ... oft im Ausland}).
\item \textbf{Justificatifs textuels :}
\begin{itemize}
\item \all{Wegen dieser Abhängigkeit können viele Staaten ihre eigene Industrie kaum entwickeln.} (conséquence économique)
\item \all{Aufgrund dieser Strukturen bleibt der Reichtum der Rohstoffe oft im Ausland.} (constat structurel)
\end{itemize}
\end{enumerate}
\end{exoresolu}

\courssub{Exercice résolu : repérage grammatical}
\begin{exoresolu}[Question 2 : transformer des verbes en noms (nominalisation)]
\textbf{Énoncé :} Transformez les verbes/adjectifs suivants en noms (nominalisation) et donnez l'article + le pluriel. Traduisez.
\begin{enumerate}
\item \all{abhängig} (adj.)
\item \all{ausbeuten} (verbe)
\item \all{verarbeiten} (verbe)
\item \all{kolonisieren} (verbe)
\end{enumerate}

\tcblower
\textbf{Solution détaillée :}
\begin{center}
\begin{tabular}{|l|l|l|l|}
\hline
\rowcolor{popblueL}
\textbf{Base} & \textbf{Nominalisation} & \textbf{Article + Pluriel} & \textbf{Français} \\
\hline
\all{abhängig} & \all{die Abhängigkeit} & \all{die Abhängigkeit, die -en} & la dépendance \\
\hline
\all{ausbeuten} & \all{die Ausbeutung} & \all{die Ausbeutung, die -en} & l'exploitation \\
\hline
\all{verarbeiten} & \all{die Verarbeitung} & \all{die Verarbeitung, die -en} & la transformation \\
\hline
\all{kolonisieren} & \all{die Kolonisierung} & \all{die Kolonisierung, die -en} & la colonisation \\
\hline
\end{tabular}
\end{center}
\end{exoresolu}

\begin{attention}[Rappel orthographe]
Les nominalisations en \all{-ung}, \all{-keit}, \all{-heit}, \all{-ion}, \all{-tion}, \all{-sion} sont \textbf{toujours féminines} (\all{die}) et font leur pluriel en \all{-en}. Majuscule obligatoire car ce sont des noms.
\end{attention}


\courssub{Consolidation : vocabulaire actif}
Pour finir cette première section, écrivez en allemand cinq phrases courtes qui réutilisent \textbf{chacun des cinq mots du programme} (\all{Abhängigkeit, Ausbeutung, Unabhängigkeit, Rohstoff, Kolonialismus}) et \textbf{au moins deux prépositions au génitif} (\all{trotz, wegen, während, aufgrund}). Vous pourrez vous en servir comme brouillon pour l'expression guidée de la section 6.

\begin{exercice}[Production écrite préparatoire]
Écrivez 5 phrases en allemand (niveau B1/B2) en utilisant le vocabulaire ci-dessus.
\end{exercice}

\coursec{Texte support : Der Neokolonialismus in Afrika}

\begin{textebox}[Der Neokolonialismus in Afrika]
Seit den sechziger Jahren sind viele afrikanische Länder politisch unabhängig. Trotz dieser Unabhängigkeit bleiben sie wirtschaftlich abhängig. Die ehemaligen Kolonialmächte und die großen Industrieländer des Nordens kaufen die Rohstoffe Afrikas — Kaffee, Kakao, Kupfer, Erdöl — zu niedrigen Preisen, denn die Preise werden nicht in Afrika, sondern im Norden festgelegt. Während die Länder des Südens billige Rohstoffe verkaufen, importieren sie teure Industrieprodukte aus dem Norden. So entsteht ein ungerechter Kreislauf: Der Kontinent verkauft viel, aber er verdient wenig. Außerdem zwingen internationale Banken die afrikanischen Staaten oft zu Krediten, die sie später kaum zurückzahlen können. Viele Politiker und Intellektuelle nennen diese Situation „Neokolonialismus“, das heißt eine neue Form des Kolonialismus ohne Soldaten und ohne Flaggen. Eine echte Unabhängigkeit verlangt deshalb mehr als politische Freiheit: Sie verlangt die Kontrolle der eigenen Rohstoffe, die Industrialisierung des Kontinents und die Zusammenarbeit der afrikanischen Länder untereinander.
\end{textebox}

\begin{definition}[Lexique de base : Kolonialismus und Abhängigkeit]
\begin{tabular}{lll}
\hline
\rowcolor{popblueL} \textbf{Deutsch (Artikel, Pluriel)} & \textbf{Français} & \textbf{Famille de mots / Derivation} \\
\hline
\all{die Kolonie, -n} & la colonie & \all{kolonial} (adj.), \all{kolonisieren} (v.), \all{der Kolonialismus} \\
\all{der Kolonialismus} & le colonialisme & \all{kolonialistisch}, \all{der Neokolonialismus} \\
\all{die Unabhängigkeit} & l'indépendance & \all{unabhängig} (adj.), \all{unabhängig werden}, \all{die Unabhängigkeitserklärung} \\
\all{die Abhängigkeit} & la dépendance & \all{abhängig} (adj.), \all{abhängig sein von + Dat.}, \all{der Abhängigkeitsvertrag} \\
\all{die Ausbeutung} & l'exploitation & \all{ausbeuten} (v.), \all{der Ausbeuter}, \all{ausbeuterisch} \\
\all{der Rohstoff, -e} & la matière première & \all{roh} (adj.), \all{der Rohstoffexport}, \all{rohstoffreich} \\
\all{die Kolonialmacht, -mächte} & la puissance coloniale & \all{die Großmacht}, \all{die Vormachtstellung} \\
\all{der Kreislauf, -läufe} & le cycle, le circuit & \all{kreisläufig}, \all{der Teufelskreis} \\
\all{der Kredit, -e} & le crédit, l'emprunt & \all{kreditieren}, \all{der Kreditnehmer}, \all{der Schuldner} \\
\all{die Industrialisierung} & l'industrialisation & \all{industrialisieren}, \all{die Industrie}, \all{industriell} \\
\all{die Zusammenarbeit} & la coopération & \all{zusammenarbeiten}, \all{der Partner}, \all{partnerschaftlich} \\
\hline
\end{tabular}
\end{definition}

\coursec{Compréhension du texte}

\courssub{Compréhension globale : Worum geht es im Text?}

\begin{methode}[Répondre à une question de compréhension en 3 temps]
\begin{enumerate}
\item \cle{Repérer} la ligne du texte qui contient l'information (souligner les \all{Schlüsselwörter}).
\item \cle{Reformuler} avec ses propres mots : ne jamais recopier la phrase du texte mot à mot.
\item \cle{Rédiger} une phrase complète : sujet + verbe conjugué en 2\up{ème} position + compléments.
\end{enumerate}
\end{methode}

\textbf{Question globale :} \all{Worum geht es im Text? Was ist die These des Autors?}

\textbf{Réponse attendue :} \all{Der Text handelt vom Neokolonialismus in Afrika. Die These des Autors ist: Die afrikanischen Länder sind zwar politisch frei, aber wirtschaftlich noch abhängig vom Norden. Echte Unabhängigkeit braucht auch wirtschaftliche Freiheit.} « Le texte traite du néocolonialisme en Afrique. La thèse de l'auteur : les pays africains sont certes politiquement libres, mais encore économiquement dépendants du Nord. Une vraie indépendance exige aussi la liberté économique. »

\begin{propriete}[La place du verbe dans la réponse (Position 2)]
Dans une phrase principale allemande, le verbe conjugué occupe \textbf{toujours la 2\up{ème} position}, même si la phrase commence par un complément de temps, de lieu ou d'objet. Le sujet passe alors après le verbe (inversion).
\end{propriete}

\begin{exemplebox}[Exemples]
\all{Die Preise werden im Norden festgelegt.} (Sujet en position 1)\\
\all{Im Norden werden die Preise festgelegt.} (Complément de lieu en position 1 $\rightarrow$ inversion)
\end{exemplebox}


\begin{popfigure}
\begin{tikzpicture}[node distance=0.55cm, every node/.style={font=\small}]
\node[draw=popblue, fill=popblueL, rounded corners, align=center, minimum width=4.2cm] (a) {\textbf{1. Lesen}\\ Lire le texte 2 fois};
\node[draw=popgreen, fill=popgreenL, rounded corners, align=center, minimum width=4.2cm, below=of a] (b) {\textbf{2. Schlüsselwörter finden}\\ Souligner les mots-clés};
\node[draw=poporange, fill=poporangeL, rounded corners, align=center, minimum width=4.2cm, below=of b] (c) {\textbf{3. Fragen beantworten}\\ Répondre en phrases complètes};
\node[draw=poppurple, fill=poppurpleL, rounded corners, align=center, minimum width=4.2cm, below=of c] (d) {\textbf{4. Zusammenfassen}\\ Résumer avec connecteurs};
\draw[-{Stealth}, thick, popdark] (a) -- (b);
\draw[-{Stealth}, thick, popdark] (b) -- (c);
\draw[-{Stealth}, thick, popdark] (c) -- (d);
\end{tikzpicture}
\legende{Organigramme de la méthode de compréhension d'un texte.}
\end{popfigure}

\courssub{Compréhension détaillée : Fragen und Antworten}

\begin{exoresolu}[Fragen zum Text — Lösungen]
\textbf{Énoncé.} Réponds par des phrases complètes en reformulant (pas de copier-coller).
\begin{enumerate}
\item \all{Seit wann sind viele afrikanische Länder unabhängig?}
\item \all{Wovon bleiben sie wirtschaftlich abhängig?}
\item \all{Wer legt die Preise der Rohstoffe fest?}
\item \all{Wodurch herrscht heute die Abhängigkeit?}
\item \all{Was verlangt eine echte Unabhängigkeit?}
\end{enumerate}
\tcblower
\textbf{Solution détaillée.}
\begin{enumerate}
\item \all{Viele afrikanische Länder sind seit den sechziger Jahren politisch unabhängig.} (Reprise du complément temporel, verbe \all{sind} en position 2.)
\item \all{Sie bleiben von der Wirtschaft des Nordens und von den ehemaligen Kolonialmächten abhängig.} (Reformulation : \all{abhängig von + Dativ}.)
\item \all{Nicht Afrika, sondern die Industrieländer des Nordens legen die Preise der Rohstoffe fest.} (Verbe à particule \all{festlegen} $\rightarrow$ \all{fest} en fin de phrase ; structure \all{nicht ... sondern}.)
\item \all{Die Abhängigkeit herrscht heute durch die ungerechten Preise, die teuren Importe und die Kredite der internationalen Banken.} (\all{Wodurch} appelle \all{durch + Akkusativ}.)
\item \all{Eine echte Unabhängigkeit verlangt die Kontrolle der Rohstoffe, die Industrialisierung und die Zusammenarbeit der afrikanischen Länder.} (Reformulation de la dernière phrase sans la recopier.)
\end{enumerate}
\end{exoresolu}

\begin{attention}[Le piège de la subordonnée avec « weil / dass / während »]
Ne commence jamais une réponse par une conjonction de subordination (\all{weil, dass, während, wenn}) en laissant le verbe en 2\up{ème} position !
\begin{itemize}
\item Faux : \all{Weil die Preise werden im Norden festgelegt.}
\item Correct (subordonnée) : \all{Weil die Preise im Norden festgelegt \textbf{werden}, verdient Afrika wenig.}
\item Correct (phrase principale, plus sûr) : \all{Die Preise werden im Norden festgelegt.}
\end{itemize}
\cle{Règle} : Dans une subordonnée, le verbe conjugué va \textbf{à la fin}.
\end{attention}

\courssub{Richtig oder Falsch? Begründe mit einem Zitat}

\begin{exercice}[Richtig / Falsch + Zitat]
Entscheide : richtig (R) oder falsch (F)? Begründe mit einem wörtlichen Zitat aus dem Text (guillemets „…“).
\begin{enumerate}
\item \all{Die afrikanischen Länder sind seit den neunziger Jahren unabhängig.}
\item \all{Die Preise der Rohstoffe werden in Afrika festgelegt.}
\item \all{Afrika verkauft billige Rohstoffe und kauft teure Industrieprodukte.}
\item \all{Der Neokolonialismus ist ein Kolonialismus mit Soldaten und Flaggen.}
\item \all{Politische Freiheit allein reicht für eine echte Unabhängigkeit nicht aus.}
\end{enumerate}
\end{exercice}

\begin{corrige}[Exercice : Richtig / Falsch]
\begin{enumerate}
\item \textbf{Falsch.} Zitat : \all{„Seit den sechziger Jahren sind viele afrikanische Länder politisch unabhängig.“}
\item \textbf{Falsch.} Zitat : \all{„Die Preise werden nicht in Afrika, sondern im Norden festgelegt.“}
\item \textbf{Richtig.} Zitat : \all{„Die Länder des Südens verkaufen billige Rohstoffe“} und \all{„importieren teure Industrieprodukte aus dem Norden.“}
\item \textbf{Falsch.} Zitat : \all{„eine neue Form des Kolonialismus ohne Soldaten und ohne Flaggen.“}
\item \textbf{Richtig.} Zitat : \all{„Eine echte Unabhängigkeit verlangt mehr als politische Freiheit.“}
\end{enumerate}
\end{corrige}

\courssub{Les moyens d'expression : l'opposition indépendance / dépendance}

Le texte oppose constamment liberté politique et dépendance économique. Repérons les opérateurs logiques.

\begin{center}
\begin{tabularx}{\linewidth}{|l|X|X|}
\hline
\rowcolor{popblueL} \textbf{Mot / Structure} & \textbf{Fonction} & \textbf{Exemple du texte} \\
\hline
\all{trotz + Genitiv} & concession (malgré) & \all{Trotz dieser Unabhängigkeit bleiben sie abhängig.} \\
\hline
\all{aber} & opposition coordonnée (mais) & \all{Der Kontinent verkauft viel, aber er verdient wenig.} \\
\hline
\all{während + Verbe à la fin} & contraste (tandis que) & \all{Während die Länder des Südens billige Rohstoffe verkaufen, importieren sie teure Produkte.} \\
\hline
\all{nicht ... sondern} & correction (non pas ... mais) & \all{Die Preise werden nicht in Afrika, sondern im Norden festgelegt.} \\
\hline
\all{heute} & repère temporel (vs. passé colonial) & \all{Wodurch herrscht heute die Abhängigkeit?} \\
\hline
\end{tabularx}
\end{center}

\begin{exemplebox}[Exemples traduits]
\all{Trotz der Unabhängigkeit bleiben sie wirtschaftlich abhängig.} « Malgré l'indépendance, ils restent économiquement dépendants. »\\
\all{Afrika ist politisch frei, aber wirtschaftlich nicht frei.} « L'Afrique est politiquement libre, mais pas économiquement. »\\
\all{Während der Süden Rohstoffe verkauft, verkauft der Norden Industrieprodukte.} « Tandis que le Sud vend des matières premières, le Nord vend des produits industriels. »
\end{exemplebox}

\begin{attention}[Faux amis et pièges syntaxiques]
\begin{itemize}
\item \all{trotz} régit le \textbf{génitif} : \all{trotz der Unabhängigkeit} (jamais *\emph{trotz die Unabhängigkeit}).
\item \all{abhängig} s'emploie avec \all{von + Dativ} : \all{abhängig von den Kolonialmächten} (jamais *\emph{abhängig aus}).
\item \all{während} (conjonction) envoie le verbe conjugué \textbf{à la fin} de la subordonnée.
\item \all{seit} (préposition temporelle « depuis ») + \textbf{Datif} : \all{seit den sechziger Jahren}.
\end{itemize}
\end{attention}

\coursec{Grammaire 1 : La nominalisation (die Nominalisierung)}

\begin{definition}[La nominalisation]
La \cle{nominalisation} consiste à transformer un verbe ou un adjectif en nom (substantif). En allemand, tout nom s'écrit avec une \textbf{majuscule} et est précédé de son article (\all{der, die, das}). C'est une caractéristique majeure de la langue écrite (style administratif, scientifique, journalistique).
\end{definition}

\begin{propriete}[Procédés principaux de nominalisation]
\begin{enumerate}
\item \textbf{Verbe $\rightarrow$ Substantif (souvent féminin \all{die} avec suffixes) :}
\begin{itemize}
\item \all{-ung} : \all{unabhängig werden} $\rightarrow$ \all{die Unabhängigkeit} ; \all{industrialisieren} $\rightarrow$ \all{die Industrialisierung} ; \all{zusammenarbeiten} $\rightarrow$ \all{die Zusammenarbeit} ; \all{ausbeuten} $\rightarrow$ \all{die Ausbeutung}.
\item \all{-heit / -keit} (adjectif $\rightarrow$ nom) : \all{frei} $\rightarrow$ \all{die Freiheit} ; \all{abhängig} $\rightarrow$ \all{die Abhängigkeit} ; \all{möglich} $\rightarrow$ \all{die Möglichkeit}.
\item \all{-nis} : \all{bitter} $\rightarrow$ \all{die Bitterkeit} (rare), \all{erfahren} $\rightarrow$ \all{die Erfahrung}.
\item \all{-tion / -sion} (mots savants) : \all{kolonisieren} $\rightarrow$ \all{die Kolonisation} ; \all{exploitieren} $\rightarrow$ \all{die Exploitation}.
\item \all{-schaft} : \all{Partner} $\rightarrow$ \all{die Partnerschaft} ; \all{Freund} $\rightarrow$ \all{die Freundschaft}.
\end{itemize}
\item \textbf{Infinitif substantivé (neutre \all{das}) :} \all{das Lesen, das Schreiben, das Festlegen, das Importieren, das Verkaufen}. Souvent après préposition : \all{beim Lesen} (en lisant), \all{zum Festlegen} (pour fixer).
\item \textbf{Adjectif substantivé (neutre \all{das}) :} \all{das Gute, das Neue, das Wichtige, das Politische, das Wirtschaftliche}.
\end{enumerate}
\end{propriete}

\begin{exemplebox}[Nominalisations repérées dans le texte]
\begin{tabular}{lll}
\hline
\rowcolor{popblueL} \textbf{Forme verbale / adj.} & \textbf{Nominalisation (article + majuscule)} & \textbf{Contexte} \\
\hline
\all{unabhängig (adj.)} & \all{die Unabhängigkeit} & \all{politisch unabhängig / die Unabhängigkeit} \\
\all{abhängig (adj.)} & \all{die Abhängigkeit} & \all{wirtschaftlich abhängig / die Abhängigkeit} \\
\all{festlegen (v. sep.)} & \all{die Festlegung} & \all{die Preise werden festgelegt / die Festlegung der Preise} \\
\all{industrialisieren (v.)} & \all{die Industrialisierung} & \all{die Industrialisierung des Kontinents} \\
\all{zusammenarbeiten (v.)} & \all{die Zusammenarbeit} & \all{zusammenarbeiten / die Zusammenarbeit} \\
\all{ausbeuten (v.)} & \all{die Ausbeutung} & (thème voisin) \\
\all{kolonisieren (v.)} & \all{der Kolonialismus / die Kolonisation} & \all{der Neokolonialismus} \\
\hline
\end{tabular}
\end{exemplebox}

\begin{exoresolu}[Exercice : Nominaliser]
\textbf{Énoncé.} Transforme les phrases en remplaçant le verbe/souligné par une nominalisation (style écrit).
\begin{enumerate}
\item \all{Die Länder exportieren Rohstoffe, aber sie importieren Maschinen.} $\rightarrow$ \all{Der Export von Rohstoffen und der Import von Maschinen ...}
\item \all{Weil die Banken Kredite geben, sind die Staaten verschuldet.} $\rightarrow$ \all{Durch die Kreditvergabe der Banken sind die Staaten verschuldet.}
\item \all{Afrikanische Politiker fordern, dass man die Preise fair festlegt.} $\rightarrow$ \all{Afrikanische Politiker fordern eine faire Festlegung der Preise.}
\end{enumerate}
\tcblower
\textbf{Solution.}
\begin{enumerate}
\item \all{Der Export von Rohstoffen und der Import von Maschinen prägen den Handel.} (Verbes $\rightarrow$ \all{der Export, der Import} + \all{von + Dativ} pour l'objet).
\item \all{Durch die Kreditvergabe der Banken sind die Staaten verschuldet.} (\all{Kredit geben} $\rightarrow$ \all{die Kreditvergabe} ; \all{weil} $\rightarrow$ \all{durch + Nominalisation + Génitif}).
\item \all{Afrikanische Politiker fordern eine faire Festlegung der Preise.} (\all{festlegen} $\rightarrow$ \all{die Festlegung} ; \all{dass}-phrase $\rightarrow$ nom + adjectif attributif).
\end{enumerate}
\end{exoresolu}

\coursec{Grammaire 2 : Le génitif et les prépositions avec le génitif}

\begin{propriete}[Formation du génitif (der Genitiv)]
Le génitif marque l'appartenance, la possession ou la relation entre deux noms (complément du nom). Il répond à la question \all{Wessen?} (De qui ? De quoi ?).
\begin{center}
\begin{tabular}{|c|c|c|c|c|}
\hline
\rowcolor{popblueL} & \textbf{Masculin} & \textbf{Féminin} & \textbf{Neutre} & \textbf{Pluriel} \\
\hline
\all{Article défini} & \all{des Vaters} & \all{der Unabhängigkeit} & \all{des Landes} & \all{der Länder} \\
\hline
\all{Article indéfini} & \all{eines Vaters} & \all{einer Unabhängigkeit} & \all{eines Landes} & \all{---} \\
\hline
\all{Article négatif} & \all{keines Vaters} & \all{keiner Unabhängigkeit} & \all{keines Landes} & \all{keiner Länder} \\
\hline
\all{Adjectif (faible)} & \all{des neuen Vaters} & \all{der neuen Unabhängigkeit} & \all{des neuen Landes} & \all{der neuen Länder} \\
\hline
\end{tabular}
\end{center}
\cle{Règle noms masculins/neutres :} Au génitif singulier, ils prennent \textbf{-(e)s} (\all{des Kontinents, des Rohstoffes/ Rohstoffs}). Noms en \all{-s, -ß, -x, -z} $\rightarrow$ \all{-es} (\all{des Preises}). Noms faibles (\all{der Name, der Staat}) $\rightarrow$ \all{des Namens, des Staates}.
\end{propriete}

\begin{propriete}[Prépositions régissant le génitif (niveau Tle)]
Ces prépositions sont du style soutenu / écrit. À l'oral, on utilise souvent le datif (ex. \all{wegen dem Wetter}), mais à l'écrit et au bac, le \textbf{génitif est obligatoire}.
\begin{center}
\begin{tabularx}{\linewidth}{|l|X|l|}
\hline
\rowcolor{popblueL} \textbf{Préposition} & \textbf{Sens} & \textbf{Exemple (contexte leçon)} \\
\hline
\all{trotz} & malgré & \all{trotz dieser Unabhängigkeit} (texte) \\
\hline
\all{während} & pendant / tandis que & \all{während der Kolonialzeit} ; \all{während der Sitzung} \\
\hline
\all{wegen} & à cause de & \all{wegen der Ausbeutung} ; \all{wegen der ungerechten Preise} \\
\hline
\all{statt / anstatt} & au lieu de & \all{statt der Zusammenarbeit} ; \all{anstatt der Unabhängigkeit} \\
\hline
\all{innerhalb / außerhalb} & à l'intérieur / à l'extérieur de & \all{innerhalb Afrikas} ; \all{außerhalb Europas} \\
\hline
\all{jenseits / diesseits} & au-delà / de ce côté de & \all{jenseits der Grenzen} \\
\hline
\end{tabularx}
\end{center}
\end{propriete}

\begin{attention}[Génitif vs Datif après ces prépositions]
\begin{itemize}
\item \textbf{Écrit standard (Bac, presse, admin) :} \all{wegen des Regens, trotz der Hitze, während der Pause}.
\item \textbf{Oral familier :} \all{wegen dem Regen, trotz der Hitze (datif fém. = identique), während der Pause (datif fém. = identique)}.
\item \cle{Consigne} : Dans vos productions écrites (thème, version, rédaction), utilisez \textbf{systématiquement le génitif} après ces prépositions.
\end{itemize}
\end{attention}

\begin{exoresolu}[Exercice : Le génitif en contexte]
\textbf{Énoncé.} Complète avec la forme correcte du génitif (article + nom + adjectif éventuel).
\begin{enumerate}
\item \all{Trotz \_\_\_\_\_\_\_\_\_\_\_\_ (die politische Freiheit) bleibt die Armut.}
\item \all{Wegen \_\_\_\_\_\_\_\_\_\_\_\_ (der niedrige Preis) verdient das Land wenig.}
\item \all{Während \_\_\_\_\_\_\_\_\_\_\_\_ (die Kolonialzeit) gab es keine Unabhängigkeit.}
\item \all{Anstatt \_\_\_\_\_\_\_\_\_\_\_\_ (die Zusammenarbeit) herrscht Konkurrenz.}
\item \all{Die Kontrolle \_\_\_\_\_\_\_\_\_\_\_\_ (die eigenen Rohstoffe) ist nötig.}
\end{enumerate}
\tcblower
\textbf{Solution.}
\begin{enumerate}
\item \all{Trotz \textbf{der politischen Freiheit} bleibt die Armut.} (Fém. : article \all{der}, adj. \all{-en}, nom inchangé).
\item \all{Wegen \textbf{des niedrigen Preises} verdient das Land wenig.} (Masc. : article \all{des}, adj. \all{-en}, nom \all{-s}).
\item \all{Während \textbf{der Kolonialzeit} gab es keine Unabhängigkeit.} (Fém. : \all{der Kolonialzeit}).
\item \all{Anstatt \textbf{der Zusammenarbeit} herrscht Konkurrenz.} (Fém. : \all{der Zusammenarbeit}).
\item \all{Die Kontrolle \textbf{der eigenen Rohstoffe} ist nötig.} (Pluriel : article \all{der}, adj. \all{-en}, nom \all{-e} pluriel).
\end{enumerate}
\end{exoresolu}

\coursec{Commentaire guidé et production écrite / orale}

\courssub{Résumé structuré : der Kurzvortrag / die Zusammenfassung}

\begin{methode}[Rédiger un résumé cohérent (4-5 phrases)]
\begin{enumerate}
\item Utilise les connecteurs logiques dans l'ordre : \all{Zuerst} (d'abord) $\rightarrow$ \all{Dann / Anschließend} (ensuite) $\rightarrow$ \all{Außerdem / Darüber hinaus} (de plus) $\rightarrow$ \all{Schließlich / Fazit} (finalement / conclusion).
\item Une idée principale par phrase.
\item Verbe conjugué en position 2 (phrase principale) ou à la fin (subordonnée \all{dass, weil, während}).
\item Vocabulaire du thème + nominalisations (\all{die Unabhängigkeit, die Abhängigkeit, die Festlegung, die Industrialisierung}).
\end{enumerate}
\end{methode}

\begin{exemplebox}[Modèle de résumé écrit (Musterlösung)]
\all{Zuerst stellt der Text fest, dass viele afrikanische Länder seit den sechziger Jahren politisch unabhängig sind. Dann erklärt er, dass die wirtschaftliche Abhängigkeit durch die Festlegung der Rohstoffpreise im Norden und den Import teurer Industrieprodukte fortbesteht. Außerdem verschärfen die Kredite internationaler Banken die Schuldenkrise. Schließlich fordert der Autor eine echte Unabhängigkeit durch die Kontrolle der Rohstoffe, die Industrialisierung und die afrikanische Zusammenarbeit.}
\end{exemplebox}

\courssub{Production guidée : Stellungnahme / Kurzessay}

\begin{experience}[Débat structuré : „Afrika braucht faire Preise, keine Almosen“]
\textbf{Consigne :} En groupes de 3. Préparez 3 arguments \textbf{pour} et 3 arguments \textbf{contre} la thèse : \all{„Entwicklungshilfe hilft Afrika nicht, fairer Handel schon.“} Utilisez le vocabulaire de la leçon et au moins 2 nominalisations et 2 prépositions avec génitif (\all{trotz, wegen, statt}) par argument. Présentez au tableau.
\begin{itemize}
\item \textit{Argument Pro :} \all{Trotz der Entwicklungshilfe wächst die Armut, weil die Handelsbedingungen ungerecht sind. Wegen der niedrigen Rohstoffpreise fehlt das Kapital für Industrialisierung. Statt Almosen braucht Afrika faire Verträge.}
\item \textit{Argument Contra :} \all{Wegen der fehlenden Infrastruktur ist Hilfe nötig. Trotz aller Probleme sichern Impfprogramme das Überleben. Statt nur Handel braucht es Bildung und Gesundheit.}
\end{itemize}
\end{experience}

\begin{exercice}[Expression écrite : Thème / Version court]
\textbf{Thème (Fr $\rightarrow$ De) :} Traduisez en allemand.
\begin{enumerate}
\item Malgré l'indépendance politique, l'exploitation économique continue.
\item À cause des prix bas des matières premières, l'Afrique ne peut pas s'industrialiser.
\item Au lieu de la coopération, on voit souvent la concurrence entre les pays africains.
\item Le contrôle des propres ressources exige une forte industrie.
\end{enumerate}
\textbf{Version (De $\rightarrow$ Fr) :} Traduisez en français.
\all{„Die ehemalige Kolonialmacht sichert sich den Zugang zu Rohstoffen durch ungerechte Verträge. Trotz formaler Souveränität entscheiden oft ausländische Konzerne über die Förderung. Eine echte Souveränität verlangt die Verarbeitung der Rohstoffe im eigenen Land.“}
\end{exercice}

\begin{corrige}[Exercice : Thème / Version]
\textbf{Thème :}
\begin{enumerate}
\item \all{Trotz der politischen Unabhängigkeit dauert die wirtschaftliche Ausbeutung an.} (\all{trotz + Gen.}, nominalisation \all{Ausbeutung}).
\item \all{Wegen der niedrigen Rohstoffpreise kann Afrika sich nicht industrialisieren.} (\all{wegen + Gen.}, \all{sich industrialisieren}).
\item \all{Anstatt der Zusammenarbeit sieht man oft Konkurrenz zwischen den afrikanischen Ländern.} (\all{anstatt + Gen.}, nominalisation \all{Zusammenarbeit, Konkurrenz}).
\item \all{Die Kontrolle der eigenen Rohstoffe erfordert eine starke Industrie.} (\all{Génitif possession}, \all{erfordern + Akk.}).
\end{enumerate}
\textbf{Version :}
« L'ancienne puissance coloniale s'assure l'accès aux matières premières par des contrats injustes. Malgré une souveraineté formelle, ce sont souvent des groupes étrangers qui décident de l'exploitation. Une vraie souveraineté exige la transformation des matières premières dans le propre pays. »
\end{corrige}

\begin{savaistu}[Cameroun \& Allemagne : Partenariat et histoire]
Le Cameroun a été un \all{Schutzgebiet} (protectorat) allemand de 1884 à 1916 (\all{Kamerun}). Après la Première Guerre mondiale, il fut partagé entre la France et la Grande-Bretagne (mandat SDN). Aujourd'hui, la coopération allemande (\all{GIZ, KfW, DAAD}) soutient l'agriculture (cacao, café), la décentralisation, l'énergie renouvelable et la formation professionnelle (ex. \all{Ausbildung} duale). Le néocolonialisme est débattu au Cameroun notamment sur les accords de pêche, l'exportation de bois brut vs transformation locale, et la dette extérieure. Le concept de \all{„Afrika wächst zusammen“} (UA/Agenda 2063) vise à briser la dépendance \all{„von außen“} par l'intégration \all{„nach innen“} (ZLECAf / AfCFTA).
\end{savaistu}

\begin{aretenir}[Synthèse de la leçon AL27 — Neokolonialismus]
\begin{itemize}
\item \textbf{Texte} : Indépendance politique (années 60) $\neq$ indépendance économique. Mécanismes : prix fixés au Nord, importations chères, dette (Banques).
\item \textbf{Vocabulaire} : \all{die Kolonie, der Kolonialismus, die Unabhängigkeit, die Abhängigkeit, die Ausbeutung, der Rohstoff, die Industrialisierung, die Zusammenarbeit} + familles de mots.
\item \textbf{Grammaire 1 — Nominalisation} : Verbe/Adj $\rightarrow$ Nom (majuscule, article). Suffixes \all{-ung, -heit, -keit, -tion, -schaft} (\all{die} fém.) ; Infinitif \all{das ...en} ; Adj. \all{das ...e}. Style écrit.
\item \textbf{Grammaire 2 — Génitif \& Prépositions} : \all{Wessen?} Formation : \all{des Vaters, der Unabhängigkeit, des Landes, der Länder}. Prépositions : \all{trotz, während, wegen, statt, anstatt, innerhalb, außerhalb, jenseits, diesseits} + \textbf{Génitif} (écrit standard).
\item \textbf{Expression} : Comprendre (global/détaillé), justifier (Zitat), résumer (connecteurs \all{zuerst, dann, außerdem, schließlich}), produire (argumentation, thème/version).
\item \textbf{Pièges francophones} : \all{trotz + Gen.} (pas Dat.), \all{abhängig von + Dat.} (pas \emph{aus}), \all{während/weil} $\rightarrow$ verbe à la fin, V2 en phrase principale, majuscules noms, \all{ß} après voyelle longue/diphtongue (\all{Preise, maßgeblich}).
\end{itemize}
\end{aretenir}

\coursec{Lexique thématique : Kolonialismus und Abhängigkeit}

Après avoir compris globalement et en détail le texte \all{Der Neokolonialismus in Afrika}, nous allons maintenant nous arrêter sur les mots-clés qui permettent de parler de ce sujet. Un bon lexique est la condition d'une bonne expression, à l'oral comme à l'écrit, et surtout au baccalauréat. Nous allons organiser ce vocabulaire en trois champs lexicaux, étudier les familles de mots, puis nous entraîner à employer ces termes dans des phrases.

\courssub{Observation : les mots du texte}

Commençons par relire quelques phrases tirées du texte support et observons les mots en gras.

\begin{textebox}[Extraits du texte support]
\all{\textbf{Die Kolonialzeit} ist vorbei, aber viele afrikanische Staaten bleiben wirtschaftlich \textbf{abhängig von} den früheren \textbf{Kolonialmächten}. Die \textbf{Ausbeutung der Rohstoffe} geht weiter: große \textbf{Konzerne} kaufen Kakao, Kaffee und Öl zu niedrigen \textbf{Preisen} und verkaufen die verarbeiteten Produkte teuer auf dem \textbf{Weltmarkt}. Die \textbf{Unabhängigkeit} ist politisch, aber nicht wirtschaftlich. Viele Länder verlangen heute eine faire Zusammenarbeit.}

\emph{Glossaire :} die Kolonialmacht (-mächte) = la puissance coloniale ; der Konzern (-e) = la grande entreprise, le groupe ; verarbeitet = transformé ; verlangen = exiger.
\end{textebox}

Observons : les noms sont tous accompagnés de leur article (\all{die, der}), certains sont composés (\all{Kolonial + Macht}, \all{Roh + Stoff}), et l'adjectif \all{abhängig} est suivi de \all{von} + datif. C'est exactement ce que nous allons systématiser.

\courssub{Champ lexical 1 : Kolonialismus und Unabhängigkeit}

\begin{definition}[Le vocabulaire de la colonisation et de l'indépendance]
\begin{itemize}
\item \cle{die Kolonie} (-n) : la colonie. \all{Kamerun war eine deutsche Kolonie.} « Le Cameroun était une colonie allemande. »
\item \cle{der Kolonialismus} (sans pluriel) : le colonialisme. Nom abstrait masculin en \all{-ismus}, comme \all{der Kapitalismus}, \all{der Kommunismus}.
\item \cle{der Neokolonialismus} (sans pluriel) : le néocolonialisme. Préfixe \all{neo-} = nouveau, comme en français.
\item \cle{die Kolonialmacht} (-mächte) : la puissance coloniale. Mot composé : \all{Kolonial-} + \all{die Macht} (le pouvoir).
\item \cle{die Kolonialzeit} (sans pluriel) : l'époque coloniale. \all{Zeit} = le temps, l'époque.
\item \cle{die Unabhängigkeit} (sans pluriel) : l'indépendance.
\item \cle{unabhängig} : indépendant. \all{Kamerun ist seit 1960 unabhängig.} « Le Cameroun est indépendant depuis 1960. »
\item \cle{die Befreiung} (-en) : la libération. Verbe : \all{befreien} (libérer).
\end{itemize}
\end{definition}

\begin{exemplebox}[Phrases modèles]
\all{Nach der Kolonialzeit erkämpften viele Völker ihre Befreiung.} « Après l'époque coloniale, beaucoup de peuples ont conquis leur libération. »

\all{Der Neokolonialismus ist eine neue Form der Abhängigkeit.} « Le néocolonialisme est une nouvelle forme de dépendance. »

\all{Die ehemaligen Kolonialmächte haben noch großen Einfluss.} « Les anciennes puissances coloniales ont encore une grande influence. »
\end{exemplebox}

\courssub{Champ lexical 2 : Abhängigkeit und Ausbeutung}

\begin{definition}[Die Ausbeutung]
\cle{die Ausbeutung} (-en) = l'exploitation, des personnes ou des ressources. Le verbe correspondant est \cle{ausbeuten} (verbe à particule séparable : \all{aus} + \all{beuten}) : \all{Die Konzerne beuten die Rohstoffe aus.} « Les groupes exploitent les matières premières. » La personne qui exploite est \all{der Ausbeuter} (-) ; la personne exploitée est \all{der/die Ausgebeutete} (-n), adjectif nominalisé qui se décline comme un adjectif.
\end{definition}

\begin{definition}[Les autres mots du champ]
\begin{itemize}
\item \cle{die Abhängigkeit} (-en) : la dépendance.
\item \cle{abhängig} : dépendant. Construction : \all{abhängig von + Dativ}.
\item \cle{die Ungerechtigkeit} (-en) : l'injustice. Adjectif : \all{ungerecht} (injuste) ; contraire : \all{gerecht} / \all{die Gerechtigkeit}.
\item \cle{die Schulden} (pluriel seulement) : les dettes. Au singulier, \all{die Schuld} signifie « la faute, la culpabilité » : attention au sens !
\end{itemize}
\end{definition}

\begin{attention}[La construction de « abhängig »]
\all{abhängig} se construit avec \cle{von + datif} : \all{abhängig von den Kolonialmächten} « dépendant des puissances coloniales ». Ne dites jamais \all{abhängig auf} : le français « dépendre de » se rend par \all{von}, pas par \all{auf}. De même : \all{Es hängt vom Wetter ab.} « Cela dépend du temps. »
\end{attention}

\begin{exemplebox}[Exemples traduits]
\all{Die Ausbeutung der Rohstoffe geht weiter.} « L'exploitation des matières premières continue. »

\all{Viele Staaten sind von ihren Schulden abhängig.} « Beaucoup d'États dépendent de leurs dettes. »

\all{Die Bauern sind von den Preisen auf dem Weltmarkt abhängig.} « Les paysans dépendent des prix sur le marché mondial. »

\all{Die Ungerechtigkeit im Handel muss enden.} « L'injustice dans le commerce doit cesser. »
\end{exemplebox}

\courssub{Champ lexical 3 : Wirtschaft und Rohstoffe}

\begin{definition}[Économie et matières premières]
\begin{itemize}
\item \cle{der Rohstoff} (-e) : la matière première. Mot composé : \all{roh} (brut) + \all{der Stoff} (la matière).
\item \cle{die Ressource} (-n) : la ressource. Souvent au pluriel : \all{natürliche Ressourcen}.
\item \cle{der Weltmarkt} (-märkte) : le marché mondial.
\item \cle{der Preis} (-e) : le prix. Attention : \all{der Preis} signifie aussi « le prix » au sens de récompense.
\item \cle{der Handelsvertrag} (-verträge) : le traité/accord commercial. \all{der Handel} = le commerce ; \all{der Vertrag} = le contrat.
\item \cle{der Konzern} (-e) : la grande entreprise, le groupe industriel.
\item \cle{die Industrie} (-n) : l'industrie. Prononciation : « in-dous-tri ».
\item \cle{die Verarbeitung} (-en) : la transformation. Verbe : \all{verarbeiten} (transformer).
\item \cle{exportieren / importieren} : exporter / importer. Verbes faibles réguliers : \all{Kamerun exportiert Kakao und importiert Maschinen.} « Le Cameroun exporte du cacao et importe des machines. »
\end{itemize}
\end{definition}

\begin{exemplebox}[Phrases modèles]
\all{Die Preise für Rohstoffe werden oft im Ausland festgelegt.} « Les prix des matières premières sont souvent fixés à l'étranger. »

\all{Ohne Industrie gibt es keine Verarbeitung vor Ort.} « Sans industrie, il n'y a pas de transformation sur place. »

\all{Faire Handelsverträge sind wichtig für beide Seiten.} « Des accords commerciaux équitables sont importants pour les deux parties. »
\end{exemplebox}

\begin{center}
\begin{tabularx}{\linewidth}{|X|X|X|}
\hline
\rowcolor{popblueL} Champ lexical & Mot allemand (article, pluriel) & Traduction française \\
\hline
Colonialisme & die Kolonie (-n) & la colonie \\
\hline
Colonialisme & der Kolonialismus (sg.) & le colonialisme \\
\hline
Colonialisme & der Neokolonialismus (sg.) & le néocolonialisme \\
\hline
Colonialisme & die Kolonialmacht (-mächte) & la puissance coloniale \\
\hline
Colonialisme & die Kolonialzeit (sg.) & l'époque coloniale \\
\hline
Colonialisme & die Unabhängigkeit (sg.) & l'indépendance \\
\hline
Colonialisme & unabhängig & indépendant \\
\hline
Colonialisme & die Befreiung (-en) & la libération \\
\hline
Dépendance & die Abhängigkeit (-en) & la dépendance \\
\hline
Dépendance & abhängig (von + Dat.) & dépendant (de) \\
\hline
Dépendance & die Ausbeutung (-en) & l'exploitation \\
\hline
Dépendance & ausbeuten & exploiter \\
\hline
Dépendance & der/die Ausgebeutete (-n) & l'exploité(e) \\
\hline
Dépendance & die Ungerechtigkeit (-en) & l'injustice \\
\hline
Dépendance & die Schulden (pl.) & les dettes \\
\hline
Économie & der Rohstoff (-e) & la matière première \\
\hline
Économie & die Ressource (-n) & la ressource \\
\hline
Économie & der Weltmarkt (-märkte) & le marché mondial \\
\hline
Économie & der Preis (-e) & le prix \\
\hline
Économie & der Handelsvertrag (-verträge) & l'accord commercial \\
\hline
Économie & der Konzern (-e) & le groupe, la grande entreprise \\
\hline
Économie & die Industrie (-n) & l'industrie \\
\hline
Économie & die Verarbeitung (-en) & la transformation \\
\hline
Économie & exportieren / importieren & exporter / importer \\
\hline
\end{tabularx}
\end{center}

\courssub{Les familles de mots}

En allemand, les familles de mots se construisent par préfixes (\all{un-}, \all{ab-}, \all{aus-}) et suffixes (\all{-heit}, \all{-keit}, \all{-ung}, \all{-ismus}). Les connaître permet de deviner le sens et le genre des mots nouveaux.

\begin{propriete}[Suffixes et genre]
\begin{itemize}
\item \all{-heit} et \all{-keit} forment des noms abstraits à partir d'adjectifs : \all{abhängig} $\rightarrow$ \all{die Abhängigkeit} ; \all{unabhängig} $\rightarrow$ \all{die Unabhängigkeit}.
\item \all{-ung} forme des noms d'action à partir de verbes : \all{ausbeuten} $\rightarrow$ \all{die Ausbeutung} ; \all{befreien} $\rightarrow$ \all{die Befreiung} ; \all{verarbeiten} $\rightarrow$ \all{die Verarbeitung}.
\item \all{-ismus} forme des noms masculins de doctrines ou de systèmes : \all{der Kolonialismus}, \all{der Neokolonialismus}.
\item Le préfixe \all{un-} donne le contraire : \all{abhängig} $\rightarrow$ \all{unabhängig} ; \all{gerecht} $\rightarrow$ \all{ungerecht}.
\end{itemize}
\end{propriete}

\begin{savaistu}[Un truc sûr pour l'article]
En allemand, les noms abstraits en \all{-heit} et \all{-keit} sont \textbf{toujours féminins} : \all{die Abhängigkeit}, \all{die Unabhängigkeit}, \all{die Möglichkeit} (la possibilité), \all{die Freiheit} (la liberté). C'est aussi vrai pour \all{-ung}, \all{-schaft} et \all{-ion} : \all{die Ausbeutung}, \all{die Wirtschaft} (l'économie), \all{die Nation}. Un moyen sûr de deviner l'article sans dictionnaire !
\end{savaistu}

\begin{popfigure}
\begin{tikzpicture}[level distance=1.6cm,
  level 1/.style={sibling distance=3.2cm},
  every node/.style={draw, rounded corners, align=center, fill=popblueL, font=\small},
  edge from parent/.style={draw=popblue, thick}]
\node[fill=poporangeL, align=center]{\all{hängen}\\ (accrocher)}
  child { node[align=center]{\all{abhängig}\\ dépendant}
    child { node[fill=popgreenL, align=center]{\all{die Abhängigkeit}\\ la dépendance} }
  }
  child { node[align=center]{\all{unabhängig}\\ indépendant}
    child { node[fill=popgreenL, align=center]{\all{die Unabhängigkeit}\\ l'indépendance} }
  };
\end{tikzpicture}
\legende{Arbre de la famille de mots autour de \all{hängen} : préfixes et suffixes.}
\end{popfigure}

De la même façon, observons la famille de \all{ausbeuten} : \all{ausbeuten} (exploiter) $\rightarrow$ \all{die Ausbeutung} (l'exploitation) $\rightarrow$ \all{der Ausbeuter} (l'exploitant) $\rightarrow$ \all{der/die Ausgebeutete} (l'exploité). Et celle de \all{Kolonie} : \all{die Kolonie} $\rightarrow$ \all{kolonial} (colonial) $\rightarrow$ \all{der Kolonialismus} $\rightarrow$ \all{die Kolonialmacht}.

\begin{attention}[Die Unabhängigkeit : pas de pluriel !]
\all{die Unabhängigkeit} est un nom abstrait qui n'a \textbf{pas de pluriel}, comme \all{die Freiheit} ou \all{die Kolonialzeit}. Pour dire « l'indépendance du Cameroun », on utilise le génitif : \all{die Unabhängigkeit Kameruns}. Exemple : \all{Die Unabhängigkeit Kameruns wurde 1960 proklamiert.} « L'indépendance du Cameroun a été proclamée en 1960. »
\end{attention}

\courssub{Expressions utiles pour le Bac}

\begin{aretenir}[Cinq expressions à savoir par cœur]
\begin{enumerate}
\item \all{von etwas abhängig sein} : dépendre de quelque chose. \all{Afrika ist noch von den Rohstoffpreisen abhängig.}
\item \all{die Rohstoffe ausbeuten} : exploiter les matières premières.
\item \all{die Preise festlegen} : fixer les prix. \all{Wer legt die Preise fest?} « Qui fixe les prix ? »
\item \all{wirtschaftlich abhängig bleiben} : rester économiquement dépendant.
\item \all{eine faire Zusammenarbeit verlangen} : exiger une coopération équitable.
\end{enumerate}
\end{aretenir}

\begin{center}
\begin{tabularx}{\linewidth}{|X|X|X|}
\hline
\rowcolor{popblueL} Français & Deutsch & Remarque \\
\hline
dépendre de & abhängig sein von + Dat. & jamais \all{auf} ! \\
\hline
exploiter les richesses & die Rohstoffe ausbeuten & verbe séparable \\
\hline
fixer les prix & die Preise festlegen & verbe séparable \\
\hline
rester dépendant & abhängig bleiben & \all{bleiben} + adjectif \\
\hline
exiger la justice & Gerechtigkeit verlangen & \all{verlangen} + Akk. \\
\hline
\end{tabularx}
\end{center}

\courssub{Entraînement}

\begin{exoresolu}[Classement dans les trois champs lexicaux]
Énoncé : classez les mots suivants dans les trois champs lexicaux (colonialisme / dépendance / économie) : \all{der Rohstoff, die Kolonie, die Ausbeutung, der Weltmarkt, die Unabhängigkeit, die Schulden, der Konzern, die Befreiung, die Ungerechtigkeit, die Verarbeitung}.
\tcblower
Solution détaillée :
\begin{itemize}
\item Colonialisme et indépendance : \all{die Kolonie}, \all{die Unabhängigkeit}, \all{die Befreiung}.
\item Dépendance et exploitation : \all{die Ausbeutung}, \all{die Schulden}, \all{die Ungerechtigkeit}.
\item Économie et matières premières : \all{der Rohstoff}, \all{der Weltmarkt}, \all{der Konzern}, \all{die Verarbeitung}.
\end{itemize}
Méthode : on se pose la question « de quoi parle ce mot : du système colonial, du rapport de domination, ou de l'échange économique ? ». Certains mots, comme \all{die Ausbeutung}, peuvent toucher deux champs : on choisit le champ dominant.
\end{exoresolu}

\begin{exercice}[Phrases à compléter]
Complétez avec le mot qui convient (article et forme corrects) : \all{Ausbeutung, abhängig, Rohstoffe, Unabhängigkeit, Preise, Handelsvertrag}.
\begin{enumerate}
\item Viele afrikanische Länder sind von den \all{\dots} auf dem Weltmarkt \all{\dots}.
\item Die \all{\dots} der Arbeiter ist eine große Ungerechtigkeit.
\item Kamerun exportiert viele \all{\dots} wie Kakao und Öl.
\item Die \all{\dots} Afrikas begann in den 1960er Jahren.
\item Die beiden Länder haben einen fairen \all{\dots} unterschrieben.
\end{enumerate}
\end{exercice}

\begin{corrige}[Exercice « Phrases à compléter »]
\begin{enumerate}
\item \all{Preisen} (datif pluriel après \all{von}) ; \all{abhängig}. « Beaucoup de pays africains dépendent des prix sur le marché mondial. »
\item \all{Ausbeutung}. « L'exploitation des travailleurs est une grande injustice. »
\item \all{Rohstoffe} (accusatif pluriel). « Le Cameroun exporte beaucoup de matières premières comme le cacao et le pétrole. »
\item \all{Unabhängigkeit}. « L'indépendance de l'Afrique a commencé dans les années 1960. »
\item \all{Handelsvertrag} (accusatif masculin : \all{einen}). « Les deux pays ont signé un accord commercial équitable. »
\end{enumerate}
\end{corrige}

\begin{attention}[Erreurs fréquentes des francophones]
\begin{itemize}
\item Ne dites pas \all{abhängig auf} : toujours \all{abhängig von + Dativ}.
\item \all{die Schuld} (la faute) $\neq$ \all{die Schulden} (les dettes) : le sens change avec le nombre !
\item N'oubliez pas la majuscule à tous les noms : \all{die ausbeutung} est fautif, il faut \all{die Ausbeutung}.
\item \all{der Preis} ne se traduit pas toujours par « le prix » : dans \all{die Preise festlegen}, il s'agit bien du prix de vente ; mais \all{einen Preis gewinnen} = « gagner un prix (récompense) ».
\item Le verbe \all{ausbeuten} est séparable : \all{Man beutet die Natur aus.} La particule \all{aus} va en fin de phrase.
\end{itemize}
\end{attention}

Ce lexique est désormais votre boîte à outils : il vous servira dans la partie grammaire (la nominalisation et le génitif) puis dans la production écrite et orale, où vous devrez résumer le texte et donner votre avis sur le néocolonialisme avec des phrases correctes et un vocabulaire précis.

\coursec{Grammaire 1 : la nominalisation (die Nominalisierung)}

Dans la section précédente, nous avons découvert le vocabulaire du colonialisme et de la dépendance : \all{die Ausbeutung}, \all{die Abhängigkeit}, \all{die Unabhängigkeit}, \all{die Entwicklung}. Avez-vous remarqué quelque chose ? Tous ces noms sont construits à partir de \textbf{verbes} ou d'\textbf{adjectifs} : \all{ausbeuten}, \all{abhängig}, \all{unabhängig}, \all{entwickeln}. Ce n'est pas un hasard : c'est un procédé grammatical capital de la langue allemande, la \cle{nominalisation} (\all{die Nominalisierung}), que nous allons maintenant étudier en détail.

\courssub{Observation : du verbe au nom}

Reprenons deux phrases tirées de notre texte support et comparons-les :

\begin{exemplebox}[Deux façons de dire la même chose]
\all{Man beutet die Rohstoffe weiter aus.} « On continue à exploiter les matières premières. » (phrase \textbf{verbale})

\all{Die Ausbeutung der Rohstoffe geht weiter.} « L'exploitation des matières premières continue. » (phrase \textbf{nominale})
\end{exemplebox}

Les deux phrases expriment exactement la même idée. Mais dans la seconde, le verbe \all{ausbeuten} est devenu un nom, \all{die Ausbeutung}, et le complément d'objet direct \all{die Rohstoffe} est devenu un complément du nom au génitif : \all{der Rohstoffe}. C'est toute la mécanique de la nominalisation.

\begin{definition}[La nominalisation]
La \cle{nominalisation} (\all{die Nominalisierung}) est un procédé qui transforme un \textbf{verbe} ou un \textbf{adjectif} en \textbf{nom}. Elle est extrêmement fréquente en allemand, surtout dans les textes journalistiques, politiques, économiques et administratifs. Elle permet de présenter une action comme un fait, une chose, un concept.
\end{definition}

\begin{savaistu}[Le style nominal, marque de la langue écrite allemande]
Les Allemands parlent du \all{Nominalstil} (style nominal) : dans la presse (\all{Die Zeit}, \all{Der Spiegel}), dans les discours politiques et les documents officiels, on préfère dire \all{die Erhöhung der Preise} (l'augmentation des prix) plutôt que \all{die Preise erhöhen sich} (les prix augmentent). Au Bac, les textes d'actualité sont pleins de nominalisations : savoir les reconnaître et les « défaire » mentalement est la clé de la compréhension de l'écrit.
\end{savaistu}

\courssub{La nominalisation des verbes}

Il existe trois grandes façons de transformer un verbe en nom.

\textbf{1. Verbe + suffixe \all{-ung} : nom féminin}

C'est le procédé le plus productif. On prend le radical du verbe et on ajoute \all{-ung}. Le nom obtenu est \textbf{toujours féminin} (\all{die}) et son pluriel est en \all{-ungen}.

\begin{propriete}[Verbe + -ung → nom féminin]
\all{ausbeuten} → \all{die Ausbeutung} (l'exploitation) ; \all{verarbeiten} → \all{die Verarbeitung} (la transformation) ; \all{entwickeln} → \all{die Entwicklung} (le développement) ; \all{regieren} → \all{die Regierung} (le gouvernement) ; \all{abhängen} → \all{die Abhängung} (rare, sens concret : suspension) ; pour la dépendance, on utilise l'adjectif \all{abhängig} → \all{die Abhängigkeit}.

Attention aux verbes à particule : la particule reste devant : \all{einführen} → \all{die Einführung} (l'introduction) ; \all{auswandern} → \all{die Auswanderung} (l'émigration).
\end{propriete}

\textbf{2. L'infinitif nominalisé : \all{das} + infinitif}

Tout infinitif allemand peut devenir un nom \textbf{neutre} (\all{das}) avec une majuscule. Ce nom désigne l'action elle-même, le fait de faire quelque chose :

\all{das Ausbeuten} (le fait d'exploiter), \all{das Handeln} (le fait d'agir, le commerce), \all{das Lesen} (la lecture, le fait de lire), \all{das Exportieren} (le fait d'exporter).

\all{Das schnelle Wachsen der Städte ist ein Problem.} « La croissance rapide des villes est un problème. »

\textbf{3. Noms dérivés sans suffixe (conversion)}

Certains noms sont formés directement sur le radical du verbe, sans suffixe apparent. Ils sont souvent \textbf{masculins} et doivent s'apprendre par cœur avec leur genre et leur pluriel :

\all{exportieren} → \all{der Export, -e} (l'exportation) ; \all{handeln} → \all{der Handel} (le commerce, pas de pluriel usuel) ; \all{kämpfen} → \all{der Kampf, -e} (la lutte) ; \all{drucken} → \all{der Druck, -e} (la pression).

\begin{center}
\begin{tabularx}{\linewidth}{|X|X|X|}
\hline
\rowcolor{popblueL} Verbe & Nom (article, pluriel) & Français \\
\hline
ausbeuten & die Ausbeutung, -en & l'exploitation \\
\hline
verarbeiten & die Verarbeitung, -en & la transformation \\
\hline
entwickeln & die Entwicklung, -en & le développement \\
\hline
regieren & die Regierung, -en & le gouvernement \\
\hline
exportieren & der Export, -e & l'exportation \\
\hline
handeln & der Handel & le commerce \\
\hline
kämpfen & der Kampf, -e & la lutte \\
\hline
drucken & der Druck, -e & la pression \\
\hline
\end{tabularx}
\end{center}

\courssub{La nominalisation des adjectifs}

Les adjectifs se transforment en noms abstraits grâce aux suffixes \all{-heit} et \all{-keit}. Ces noms sont \textbf{toujours féminins} (\all{die}) et font leur pluriel en \all{-en}.

\begin{propriete}[Adjectif + -heit / -keit → nom féminin]
\begin{itemize}
\item \all{-heit} après la plupart des adjectifs simples (radicaux courts) : \all{frei} → \all{die Freiheit, -en} (la liberté) ; \all{krank} → \all{die Krankheit, -en} (la maladie) ; \all{wahr} → \all{die Wahrheit, -en} (la vérité) ; \all{selbstständig} → \all{die Selbstständigkeit} (l'autonomie).
\item \all{-keit} après les adjectifs se terminant en \all{-ig}, \all{-lich}, \all{-bar}, \all{-sam}, \all{-haft} : \all{abhängig} → \all{die Abhängigkeit, -en} (la dépendance) ; \all{unabhängig} → \all{die Unabhängigkeit, -en} (l'indépendance) ; \all{möglich} → \all{die Möglichkeit, -en} (la possibilité) ; \all{verantwortlich} → \all{die Verantwortlichkeit} (la responsabilité).
\end{itemize}
\end{propriete}

Comparaison avec le français : le français fait de même avec \emph{-té} et \emph{-ment} (libre → liberté ; développer → développement), mais l'allemand utilise ce procédé beaucoup plus systématiquement, y compris à l'oral soutenu.

\courssub{La transformation des compléments : le point capital}

C'est ici que se joue l'essentiel de la grammaire. Quand on nominalise, que deviennent le sujet et le complément d'objet ?

\textbf{1. Le complément d'objet direct (COD) devient un génitif}

\all{Man beutet die Rohstoffe aus.} → \all{die Ausbeutung \textbf{der Rohstoffe}}

\all{Man verarbeitet den Kaffee im eigenen Land.} → \all{die Verarbeitung \textbf{des Kaffees} im eigenen Land}

\textbf{2. Le sujet (l'agent) devient \all{durch} + accusatif (ou un génitif possessif)}

\all{Die Konzerne beuten die Rohstoffe aus.} → \all{die Ausbeutung der Rohstoffe \textbf{durch die Konzerne}} (l'exploitation des matières premières par les multinationales)

Si le sujet est un possesseur ou un acteur évident, on peut aussi utiliser le génitif : \all{die Politik \textbf{der Regierung}} (la politique du gouvernement).

\begin{popfigure}
\begin{center}
\begin{tikzpicture}[node distance=1.2cm and 2.6cm]
\node[draw=popblue, fill=popblueL, rounded corners, thick, align=center, text width=4.6cm] (v) {\textbf{Phrase verbale}\\[2pt] \all{Die Konzerne beuten}\\ \all{die Rohstoffe aus.}};
\node[draw=poporange, fill=poporangeL, rounded corners, thick, align=center, text width=4.8cm, right=of v] (n) {\textbf{Construction nominale}\\[2pt] \all{die Ausbeutung der Rohstoffe}\\ \all{durch die Konzerne}};
\draw[-{Stealth[length=3mm]}, very thick, popdark] (v) -- (n) node[midway, above, align=center, font=\small] {\all{Nominalisierung}};
\node[below=0.5cm of v, align=center, font=\small, text=popmuted] {Sujet + verbe + COD};
\node[below=0.5cm of n, align=center, font=\small, text=popmuted] {Nom + génitif + \all{durch} + acc.};
\end{tikzpicture}
\end{center}
\legende{Schéma de transformation : du verbe au nom}
\end{popfigure}

\begin{exemplebox}[Exemples traduits]
\all{Die Entwicklung der Industrie ist langsam.} « Le développement de l'industrie est lent. » ↔ \all{Die Industrie entwickelt sich langsam.} « L'industrie se développe lentement. »

\all{Der Export des Kakaos bringt viel Geld.} « L'exportation du cacao rapporte beaucoup d'argent. » ↔ \all{Kamerun exportiert Kakao.} « Le Cameroun exporte du cacao. »

\all{Nach der Unabhängigkeit der afrikanischen Staaten blieb die Abhängigkeit vom Ausland.} « Après l'indépendance des États africains, la dépendance vis-à-vis de l'étranger est restée. »
\end{exemplebox}

\begin{attention}[Le génitif, pas « von » !]
Après une nominalisation, le complément se met au \textbf{génitif} dans la langue soignée : on écrit \all{die Ausbeutung \textbf{der Rohstoffe}} et non \all{die Ausbeutung \textbf{von den Rohstoffen}}. La construction \all{von} + datif existe dans la langue parlée, mais elle est considérée comme négligée à l'écrit : au Bac, elle vous coûtera des points. Autre piège : n'oubliez pas la \textbf{majuscule} au nom nominalisé (\all{die Entwicklung}, jamais \all{die entwicklung}).
\end{attention}

\courssub{La dénominalisation : une stratégie de lecture}

Face à un texte difficile, le bon lecteur fait le chemin inverse : il transforme mentalement les constructions nominales en phrases verbales. C'est la \cle{dénominalisation}.

\begin{methode}[Comprendre une longue construction nominale]
\begin{enumerate}
\item Repérer le nom nominalisé (souvent en \all{-ung}, \all{-heit}, \all{-keit}, ou infinitif avec \all{das}).
\item Retrouver le verbe ou l'adjectif d'origine : \all{die Verarbeitung} → \all{verarbeiten}.
\item Transformer le génitif en complément d'objet direct : \all{der Rohstoffe} → \all{die Rohstoffe}.
\item Reconstruire la phrase verbale : \all{die Verarbeitung der Rohstoffe im eigenen Land} → \all{Man verarbeitet die Rohstoffe im eigenen Land.} « On transforme les matières premières dans son propre pays. »
\item Traduire alors la phrase simple obtenue.
\end{enumerate}
\end{methode}

\courssub{Texte d'application}

\begin{textebox}[Rohstoffe und Abhängigkeit]
\all{Seit der Unabhängigkeit vieler afrikanischer Staaten in den sechziger Jahren hat sich die wirtschaftliche Situation nur langsam verändert. Die Ausbeutung der Rohstoffe durch internationale Konzerne geht weiter. Kamerun exportiert Kakao, Kaffee und Holz, aber die Verarbeitung dieser Produkte findet meistens im Ausland statt. Das bedeutet: die Gewinne bleiben nicht im Land. Die Abhängigkeit vom Weltmarkt ist groß, denn die Preise der Rohstoffe steigen und fallen schnell. Viele Experten fordern deshalb die Industrialisierung Afrikas: die Verarbeitung der Rohstoffe im eigenen Land würde Arbeitsplätze schaffen und die Entwicklung der Wirtschaft beschleunigen. Auch die Zusammenarbeit der afrikanischen Staaten untereinander ist wichtig. Der Handel zwischen den Ländern des Kontinents ist noch schwach. Die Gründung neuer Fabriken und die Ausbildung junger Fachkräfte sind erste Schritte in die richtige Richtung. Die Freiheit der Wirtschaft beginnt mit der Kontrolle der eigenen Ressourcen.}

\textbf{Glossaire} : der Rohstoff, -e (la matière première) ; der Konzern, -e (la multinationale) ; der Gewinn, -e (le profit) ; der Weltmarkt, die Weltmärkte (le marché mondial) ; fordern (exiger, réclamer) ; der Arbeitsplatz, die Arbeitsplätze (l'emploi) ; beschleunigen (accélérer) ; die Fachkraft, die Fachkräfte (le travailleur qualifié) ; die Ressource, -n (la ressource).
\end{textebox}

\textbf{Questions de compréhension :}
\begin{enumerate}
\item \all{Was exportiert Kamerun?}
\item \all{Wo findet die Verarbeitung der Produkte meistens statt?}
\item \all{Was fordern viele Experten?}
\item Relevez dans le texte cinq nominalisations et donnez le verbe ou l'adjectif d'origine.
\end{enumerate}

\courssub{Exercices d'application}

\begin{exoresolu}[Nominaliser]
Transformez la phrase verbale en construction nominale : \all{Die Konzerne beuten die Rohstoffe Afrikas aus.}
\tcblower
Étape 1 : le verbe \all{ausbeuten} devient le nom \all{die Ausbeutung}. Étape 2 : le COD \all{die Rohstoffe Afrikas} devient un génitif : \all{der Rohstoffe Afrikas}. Étape 3 : le sujet \all{die Konzerne} devient \all{durch die Konzerne}. Résultat : \all{Die Ausbeutung der Rohstoffe Afrikas durch die Konzerne geht weiter.} « L'exploitation des matières premières africaines par les multinationales continue. »
\end{exoresolu}

\begin{exoresolu}[Dénominaliser]
Transformez en phrase verbale : \all{Die Verarbeitung des Kakaos im eigenen Land schafft Arbeitsplätze.}
\tcblower
\all{die Verarbeitung} → verbe \all{verarbeiten} ; \all{des Kakaos} (génitif) → COD \all{den Kakao}. Phrase verbale : \all{Man verarbeitet den Kakao im eigenen Land. Das schafft Arbeitsplätze.} « On transforme le cacao dans son propre pays. Cela crée des emplois. »
\end{exoresolu}

\begin{exercice}[À vous de jouer]
\textbf{A. Nominalisez} les phrases suivantes (donnez la construction nominale complète) :
\begin{enumerate}
\item \all{Man entwickelt die Industrie des Landes.}
\item \all{Kamerun exportiert Kakao und Kaffee.}
\item \all{Die Regierung verbessert die Ausbildung der Jugend.}
\item \all{Die Staaten arbeiten zusammen.} (\all{zusammenarbeiten} → \all{die Zusammenarbeit})
\item \all{Die Preise steigen schnell.} (\all{steigen} → \all{der Anstieg} / \all{das Steigen})
\end{enumerate}
\textbf{B. Dénominalisez} les constructions suivantes tirées du texte :
\begin{enumerate}
\item \all{die Gründung neuer Fabriken}
\item \all{die Ausbildung junger Fachkräfte}
\item \all{die Kontrolle der eigenen Ressourcen}
\item \all{die Abhängigkeit vom Weltmarkt}
\item \all{die Industrialisierung Afrikas}
\end{enumerate}
\end{exercice}

\begin{corrige}[Exercice « À vous de jouer »]
\textbf{A.} 1. \all{die Entwicklung der Industrie des Landes} ; 2. \all{der Export von Kakao und Kaffee durch Kamerun} (ou mieux : \all{Kameruns Export von Kakao und Kaffee}) ; 3. \all{die Verbesserung der Ausbildung der Jugend durch die Regierung} ; 4. \all{die Zusammenarbeit der Staaten} ; 5. \all{der schnelle Anstieg der Preise}.

\textbf{B.} 1. \all{Man gründet neue Fabriken.} ; 2. \all{Man bildet junge Fachkräfte aus.} ; 3. \all{Man kontrolliert die eigenen Ressourcen.} ; 4. \all{Die Wirtschaft hängt vom Weltmarkt ab.} ; 5. \all{Man industrialisiert Afrika. / Afrika wird industrialisiert.}
\end{corrige}

\begin{aretenir}[L'essentiel]
\begin{itemize}
\item La nominalisation transforme un verbe ou un adjectif en nom : verbe + \all{-ung} → \all{die Ausbeutung} ; adjectif + \all{-heit}/\all{-keit} → \all{die Abhängigkeit} ; infinitif → \all{das Ausbeuten}.
\item Le COD devient un \textbf{génitif} : \all{die Ausbeutung der Rohstoffe} ; le sujet devient \all{durch} + accusatif : \all{durch die Konzerne}.
\item À l'écrit du Bac : le génitif, jamais \all{von} + datif.
\item Pour comprendre un texte difficile : dénominalisez mentalement (nom → verbe, génitif → COD).
\end{itemize}
\end{aretenir}

\coursec{Grammaire 2 : le génitif et les prépositions avec le génitif}

Dans la section précédente, nous avons appris à transformer un verbe en nom : \all{ausbeuten} devient \all{die Ausbeutung}, \all{verarbeiten} devient \all{die Verarbeitung}. Mais une question se pose immédiatement : comment dire \emph{l'exploitation \textbf{des matières premières}}, \emph{la transformation \textbf{du cacao}} ? En français, la petite préposition « de » fait tout le travail. En allemand, c'est un \cle{cas} qui s'en charge : le \cle{génitif} (\all{der Genitiv}). C'est exactement le cas dont nous avons besoin pour compléter nos nominalisations — et c'est aussi le cas exigé par des prépositions très fréquentes dans notre texte sur le néocolonialisme : \all{trotz}, \all{wegen}, \all{während}, \all{aufgrund}.

\courssub{Observation : le génitif dans le texte}

Relisons quelques phrases clés du texte support \emph{Der Neokolonialismus in Afrika} et observons les groupes de mots soulignés par la grammaire.

\begin{textebox}[Extraits du texte support]
Trotz der Unabhängigkeit bleiben viele afrikanische Staaten wirtschaftlich abhängig. Die Ausbeutung der Rohstoffe liegt oft in den Händen ausländischer Konzerne. Wegen dieser Abhängigkeit kann der Staat kaum in Bildung und Gesundheit investieren. Während der Kolonialzeit wurden die Strukturen der Wirtschaft nur für den Export gebaut. Aufgrund dieser Strukturen exportiert das Land den Kakao des Landes unverarbeitet. Die Verarbeitung des Kakaos findet außerhalb Afrikas statt.
\end{textebox}

\begin{definition}[Glossaire des extraits]
\begin{itemize}
\item \all{die Unabhängigkeit, -en} : l'indépendance
\item \all{die Ausbeutung, -en} : l'exploitation (au sens d'abus)
\item \all{der Rohstoff, -e} : la matière première
\item \all{die Abhängigkeit, -en} : la dépendance
\item \all{der Konzern, -e} : le groupe industriel, la multinationale
\item \all{die Kolonialzeit} (nur Singular) : l'époque coloniale
\item \all{die Verarbeitung, -en} : la transformation (industrielle)
\item \all{unverarbeitet} : non transformé, brut
\end{itemize}
\end{definition}

Observons : après \all{trotz}, \all{wegen}, \all{während}, \all{aufgrund}, et après les noms \all{Ausbeutung}, \all{Verarbeitung}, l'article change de forme : \all{der Unabhängigkeit}, \all{der Rohstoffe}, \all{des Kakaos}, \all{der Kolonialzeit}. Ce n'est ni le nominatif (le sujet) ni l'accusatif (le complément direct) : c'est le \cle{génitif}, le cas de la possession et du complément du nom.

\courssub{Les formes du génitif}

\begin{methode}[La question magique : \all{wessen} ?]
Pour identifier ou vérifier un génitif, pose toujours la question \all{wessen ?} (« de qui ? de quoi ? »).
\begin{enumerate}
\item Repère le nom principal : \all{die Ausbeutung} (l'exploitation).
\item Pose la question : l'exploitation — \all{wessen} ? — de quoi ?
\item La réponse est au génitif : \all{der Rohstoffe} (des matières premières).
\item Vérifie la forme de l'article : \all{der} (féminin/pluriel) ou \all{des} (masculin/neutre).
\end{enumerate}
\end{methode}

\begin{propriete}[Formation du génitif]
\begin{itemize}
\item \textbf{Masculin et neutre} : l'article devient \all{des} et le nom prend la terminaison \textbf{-(e)s} : \all{der Rohstoff} $\rightarrow$ \all{des Rohstoffs} ; \all{das Land} $\rightarrow$ \all{des Landes} ; \all{der Vater} $\rightarrow$ \all{des Vaters} ; \all{das Kind} $\rightarrow$ \all{des Kindes}.
\item \textbf{Féminin et pluriel} : l'article devient \all{der} et le nom \textbf{ne change pas} : \all{die Mutter} $\rightarrow$ \all{der Mutter} ; \all{die Kolonie} $\rightarrow$ \all{der Kolonie} ; \all{die Kinder} $\rightarrow$ \all{der Kinder}.
\item \textbf{Adjectif} : au génitif, l'adjectif prend toujours la désinence faible \textbf{-en} après l'article : \all{des kleinen Kindes}, \all{der reichen Länder}.
\item \textbf{Terminaison du nom (masc./neutre)} : \textbf{-es} pour les noms d'une syllabe (\all{des Landes}, \all{des Kindes}) ; \textbf{-s} pour les noms de plusieurs syllabes (\all{des Vaters}, \all{des Rohstoffs}). Les noms en \all{-s, -ß, -z, -x} prennent toujours \textbf{-es} : \all{das Haus} $\rightarrow$ \all{des Hauses} ; \all{der Preis} $\rightarrow$ \all{des Preises}.
\item \textbf{N-Deklination (noms faibles)} : certains masculins (\all{der Name, der Mensch, der Student, der Herr, der Nachbar, der Bauer}) prennent \textbf{-n} / \textbf{-en} à tous les cas sauf nominatif singulier : \all{der Name} $\rightarrow$ \all{des Namens} ; \all{der Mensch} $\rightarrow$ \all{des Menschen}.
\end{itemize}
\end{propriete}

\begin{center}
\begin{tabular}{|l|l|l|l|}
\hline
\rowcolor{popblueL} & \textbf{Article} & \textbf{Adjectif} & \textbf{Nom} \\ \hline
Masculin & des & reichen & Rohstoffs \\ \hline
Féminin & der & wirtschaftlichen & Abhängigkeit \\ \hline
Neutre & des & kleinen & Landes \\ \hline
Pluriel & der & reichen & Länder \\ \hline
\multicolumn{4}{|l|}{\textit{Ex. N-Deklination : \all{des Namens}, \all{des Menschen} (masc. faible)}} \\ \hline
\end{tabular}
\end{center}

\begin{exemplebox}[Exemples traduits]
\begin{itemize}
\item \all{Die Unabhängigkeit Kameruns wurde 1960 proklamiert.} « L'indépendance du Cameroun a été proclamée en 1960. »
\item \all{Die Rohstoffe Afrikas sind sehr wertvoll.} « Les matières premières de l'Afrique sont très précieuses. »
\item \all{Die Ausbeutung der Rohstoffe bringt dem Land wenig Geld.} « L'exploitation des matières premières rapporte peu d'argent au pays. »
\item \all{Die Verarbeitung des Kakaos findet in Europa statt.} « La transformation du cacao a lieu en Europe. »
\item \all{Der Preis des Kaffees fällt jedes Jahr.} « Le prix du café baisse chaque année. »
\item \all{Die Wirtschaft der Schweiz ist sehr stark.} « L'économie de la Suisse est très forte. »
\end{itemize}
\end{exemplebox}

\courssub{Emploi 1 : la possession et l'appartenance}

Le génitif exprime d'abord la \cle{possession} : il répond à « à qui ? de quoi ? ». Le français utilise la préposition « de » ; l'allemand utilise la déclinaison, \textbf{sans préposition}.

\begin{center}
\begin{tabularx}{\linewidth}{|X|X|X|}
\hline
\rowcolor{popblueL} \textbf{Français} & \textbf{Deutsch} & \textbf{Remarque} \\ \hline
l'indépendance du Cameroun & die Unabhängigkeit Kameruns & nom propre sans article + -s \\ \hline
les richesses de l'Afrique & die Reichtümer Afrikas & nom propre sans article + -s \\ \hline
l'exploitation des richesses & die Ausbeutung der Reichtümer & article \all{der} au pluriel \\ \hline
le prix du cacao & der Preis des Kakaos & article \all{des} + -s \\ \hline
l'économie de la Suisse & die Wirtschaft der Schweiz & nom propre fém. avec article \\ \hline
\end{tabularx}
\end{center}

\begin{propriete}[Les noms propres au génitif]
\begin{itemize}
\item Les noms propres \textbf{sans article} prennent simplement un \textbf{-s} final, placé après le nom qu'ils complètent : \all{Kamerun} $\rightarrow$ \all{Kameruns} ; \all{Afrika} $\rightarrow$ \all{Afrikas} ; \all{Deutschland} $\rightarrow$ \all{Deutschlands}. Exemple : \all{Die Geschichte Deutschlands ist komplex.} « L'histoire de l'Allemagne est complexe. »
\item Les noms de pays \textbf{avec article} (féminins ou pluriels : \all{die Schweiz}, \all{die Türkei}, \all{die USA}, \all{der Iran}, \all{der Libanon}) se déclinent normalement : \all{die Wirtschaft der Schweiz}, \all{die Politik des Iran}.
\item Exception : les noms propres se terminant déjà par un \emph{s} (ou un son sifflant \all{-ß, -z, -x}) prennent une apostrophe (sans espace) ou la construction avec \all{von} : \all{Max Frischs Werke} ou \all{die Werke von Max Frisch} « les œuvres de Max Frisch ».
\end{itemize}
\end{propriete}

\courssub{Emploi 2 : le complément du nom après une nominalisation}

C'est le lien direct avec la section précédente. Quand on nominalise un verbe transitif, le complément d'objet direct (accusatif) du verbe devient un \textbf{génitif} attribut du nom déverbal :

\begin{center}
\begin{tabularx}{\linewidth}{|X|X|}
\hline
\rowcolor{popblueL} \textbf{Phrase avec verbe (Akkusativ)} & \textbf{Nominalisation (Genitiv)} \\ \hline
Man beutet die Rohstoffe aus. & die Ausbeutung der Rohstoffe \\ \hline
Man verarbeitet den Kakao. & die Verarbeitung des Kakaos \\ \hline
Man exportiert die Rohstoffe. & der Export der Rohstoffe \\ \hline
Man kontrolliert die Preise. & die Kontrolle der Preise \\ \hline
\end{tabularx}
\end{center}

\begin{attention}[L'erreur n° 1 des francophones : le -s oublié]
Au génitif masculin et neutre, le nom prend \textbf{obligatoirement} -(e)s. Le français ne marque rien sur le nom (« l'exploitation \emph{de la} matière première »), donc l'élève francophone oublie la terminaison :
\begin{itemize}
\item \all{die Ausbeutung der Rohstoff} ✗ (faux : article \all{der} au lieu de \all{des} + pas de -s)
\item \all{die Ausbeutung des Rohstoffs} ✓
\item \all{die Verarbeitung der Kakao} ✗ (faux : article \all{der} + pas de -s)
\item \all{die Verarbeitung des Kakaos} ✓
\end{itemize}
Réflexe de survie : dès que tu vois \all{des}, vérifie que le nom qui suit porte bien son \textbf{-s} ou \textbf{-es}.
\end{attention}

\courssub{Les prépositions suivies du génitif}

Un groupe de prépositions très utiles à l'écrit exige le génitif. Elles sont toutes présentes dans notre texte sur le néocolonialisme.

\begin{center}
\begin{tabularx}{\linewidth}{|l|X|X|}
\hline
\rowcolor{popblueL} \textbf{Präposition} & \textbf{Français} & \textbf{Beispiel} \\ \hline
wegen & à cause de & wegen dieser Abhängigkeit \\ \hline
trotz & malgré & trotz der Unabhängigkeit \\ \hline
während & pendant & während der Kolonialzeit \\ \hline
aufgrund & en raison de & aufgrund dieser Strukturen \\ \hline
innerhalb & à l'intérieur de & innerhalb des Landes \\ \hline
außerhalb & en dehors de & außerhalb Afrikas \\ \hline
anstatt / statt & au lieu de & statt einer eigenen Industrie \\ \hline
\end{tabularx}
\end{center}

\begin{exemplebox}[Exemples traduits en contexte]
\begin{itemize}
\item \all{Trotz der Unabhängigkeit bleibt das Land arm.} « Malgré l'indépendance, le pays reste pauvre. »
\item \all{Wegen der Schulden kann der Staat nicht investieren.} « À cause des dettes, l'État ne peut pas investir. »
\item \all{Während der Kolonialzeit baute man nur Straßen für den Export.} « Pendant l'époque coloniale, on n'a construit que des routes pour l'exportation. »
\item \all{Aufgrund dieser Strukturen bleibt das Land vom Weltmarkt abhängig.} « En raison de ces structures, le pays reste dépendant du marché mondial. »
\item \all{Die Verarbeitung findet außerhalb des Kontinents statt.} « La transformation a lieu en dehors du continent. »
\item \all{Statt einer eigenen Industrie gibt es nur Minen.} « Au lieu d'une industrie propre, il n'y a que des mines. »
\end{itemize}
\end{exemplebox}

\begin{attention}[\all{während} : temps ou opposition ?]
\all{Während} + génitif exprime uniquement le \textbf{temps} (« pendant ») : \all{während der Kolonialzeit} « pendant l'époque coloniale ». Pour exprimer l'\textbf{opposition} (« tandis que »), le français utilise une seule tournure, mais l'allemand exige une \textbf{subordonnée} avec le verbe à la fin : \all{Während der Norden reich ist, bleibt der Süden arm.} « Tandis que le Nord est riche, le Sud reste pauvre. » Ne dis jamais \all{während des Reichtums} pour « tandis que » !
\end{attention}

\begin{attention}[Langue parlée contre langue écrite]
À l'oral, beaucoup de locuteurs allemands disent \all{wegen dem Wetter} (avec le datif). C'est toléré dans la conversation, mais c'est une \textbf{faute à l'écrit} et au baccalauréat. Écris toujours : \all{wegen des Wetters}, \all{trotz der Unabhängigkeit}, \all{während des Krieges}. À l'écrit, le génitif est exigé.
\end{attention}

\begin{savaistu}[Le génitif, un cas menacé ?]
En allemand populaire, le génitif recule devant \all{von} + datif : \all{das Buch von meinem Vater} au lieu de \all{das Buch meines Vaters}. Un livre humoristique célèbre en Allemagne s'intitule même \all{„Der Dativ ist dem Genitiv sein Tod“} (« Le datif est la mort du génitif », Bastian Sick). Mais dans la langue soutenue, journalistique et scolaire — celle du Bac — le génitif reste la norme. Le maîtriser, c'est montrer un allemand élégant et précis.
\end{savaistu}

\begin{popfigure}
\begin{tikzpicture}[mindmap, grow cyclic, every node/.style={concept, execute at begin node=\hskip0pt}, root concept/.append style={concept color=popblue, font=\bfseries}, level 1/.append style={level distance=3.2cm, sibling angle=90}, level 2/.append style={level distance=2.6cm, sibling angle=60}, text=white]
\node[align=center]{der Genitiv\\ wessen ?}
  child[concept color=popgreen] { node {Formen} 
    child { node[concept color=popgreen!70, align=center]{des + -(e)s\\ m / n} }
    child { node[concept color=popgreen!70, align=center]{der\\ f / Pl.} }
    child { node[concept color=popgreen!70, align=center]{N-Deklination\\ des Namens} }
  }
  child[concept color=poporange] { node {Emplois}
    child { node[concept color=poporange!70, align=center]{Possession\\ Kameruns} }
    child { node[concept color=poporange!70, align=center]{Nom + Gen.\\ die Ausbeutung\\ der Rohstoffe} }
  }
  child[concept color=poppurple] { node {Präpositionen}
    child { node[concept color=poppurple!70, align=center]{wegen, trotz\\ während} }
    child { node[concept color=poppurple!70, align=center]{aufgrund, statt\\ innerhalb\\ außerhalb} }
  }
  child[concept color=poppink] { node {Achtung}
    child { node[concept color=poppink!70, align=center]{-s nicht\\ vergessen !} }
    child { node[concept color=poppink!70, align=center]{écrit : Genitiv\\ oral : Dativ} }
  };
\end{tikzpicture}
\legende{Carte mentale : le génitif, ses formes, ses emplois et ses pièges.}
\end{popfigure}

\begin{exoresolu}[Transforme \all{von} + datif en vrai génitif]
\textbf{Énoncé.} Réécris les groupes suivants avec un génitif correct (langue écrite soutenue) : a) die Ausbeutung von dem Rohstoff ; b) die Unabhängigkeit von Kamerun ; c) wegen von dem Wetter ; d) die Verarbeitung von dem Kakao.
\tcblower
\textbf{Solution détaillée.}
a) \all{der Rohstoff} est masculin : \all{des} + nom en -s $\rightarrow$ \all{die Ausbeutung des Rohstoffs}.\\
b) Nom propre sans article : on ajoute -s $\rightarrow$ \all{die Unabhängigkeit Kameruns}.\\
c) \all{das Wetter} est neutre : \all{des} + -s $\rightarrow$ \all{wegen des Wetters}.\\
d) \all{der Kakao} est masculin : \all{des} + -s $\rightarrow$ \all{die Verarbeitung des Kakaos}.\\
Méthode : identifier le genre du nom, choisir \all{des} (m/n) ou \all{der} (f/pl), ajouter -(e)s au masculin et au neutre.
\end{exoresolu}

\begin{exercice}[À toi de jouer]
\begin{enumerate}
\item Décline au génitif : der Kontinent ; die Kolonie ; das Land ; die Rohstoffe (pluriel).
\item Complète avec \all{wegen}, \all{trotz} ou \all{während} au génitif : a) \all{... der Schulden kann der Staat keine Schulen bauen.} b) \all{... der Unabhängigkeit bleibt das Land abhängig.} c) \all{... der Kolonialzeit exportierte man nur Rohstoffe.}
\item Transforme (langue soutenue) : die Kontrolle von den Preisen ; der Export von dem Kaffee ; die Geschichte von Afrika.
\end{enumerate}
\end{exercice}

\begin{corrige}[Exercice : le génitif]
\begin{enumerate}
\item \all{des Kontinents} ; \all{der Kolonie} ; \all{des Landes} ; \all{der Rohstoffe}.
\item a) \all{Wegen der Schulden...} (cause) ; b) \all{Trotz der Unabhängigkeit...} (concession) ; c) \all{Während der Kolonialzeit...} (temps).
\item \all{die Kontrolle der Preise} ; \all{der Export des Kaffees} ; \all{die Geschichte Afrikas}.
\end{enumerate}
\end{corrige}

\begin{aretenir}[L'essentiel du génitif]
\begin{itemize}
\item Question : \all{wessen ?} (de qui ? de quoi ?).
\item Formes : masculin/neutre \all{des} + nom en -(e)s ; féminin/pluriel \all{der}, nom inchangé ; adjectif en -en ; N-Deklination : \all{des Namens}.
\item Noms propres sans article : -s simple (\all{Kameruns}, \all{Afrikas}) ; apostrophe si sifflante finale (\all{Max Frischs}).
\item Après une nominalisation, l'objet direct (accusatif) devient génitif : \all{die Ausbeutung der Rohstoffe}.
\item Prépositions + génitif : \all{wegen, trotz, während, aufgrund, innerhalb, außerhalb, statt}.
\item À l'écrit, jamais de datif après ces prépositions : \all{wegen des Wetters}, pas \all{wegen dem Wetter}.
\end{itemize}
\end{aretenir}

\coursec{Commentaire guidé et production}

Dans la section précédente, nous avons étudié le génitif et les prépositions suivies du génitif (\all{wegen der Abhängigkeit}, \all{trotz der Unabhängigkeit}, \all{während der Kolonialzeit}). Ces outils grammaticaux, ajoutés à la nominalisation, vont maintenant nous servir à \cle{produire} : commenter le texte \all{Der Neokolonialismus in Afrika}, donner notre avis et débattre. C'est l'aboutissement de la leçon : passer de la compréhension à l'expression.

\courssub{La méthode du commentaire de texte en quatre parties}

Au baccalauréat comme en classe, un commentaire de texte allemand (\all{der Kommentar}) suit toujours la même architecture en quatre parties. La voici, avec la problématique modèle de notre leçon : \all{Ist Afrika wirklich unabhängig?} (« L'Afrique est-elle vraiment indépendante ? »).

\begin{methode}[Commenter un texte en quatre étapes]
\begin{enumerate}
\item \textbf{\all{Die Einleitung} (l'introduction)} : présenter l'auteur ou la source, le thème du texte et la problématique. Formule utile : \all{Der Text behandelt das Problem des Neokolonialismus und stellt die Frage: Ist Afrika wirklich unabhängig?} (« Le texte traite du problème du néocolonialisme et pose la question : l'Afrique est-elle vraiment indépendante ? »).
\item \textbf{\all{Die Zusammenfassung} (le résumé)} : reformuler les idées principales avec ses propres mots, sans citer de longs passages. Utiliser \all{Der Autor meint, dass...}, \all{Er erklärt, dass...}.
\item \textbf{\all{Die Analyse} (l'analyse)} : dégager les arguments et les moyens linguistiques : les nominalisations (\all{die Ausbeutung}, \all{die Abhängigkeit}), les oppositions (\all{unabhängig — abhängig}), le champ lexical de la domination.
\item \textbf{\all{Der Schluss} (la conclusion)} : répondre à la problématique et donner son avis personnel justifié (\all{Meiner Meinung nach...}).
\end{enumerate}
\end{methode}

\begin{popfigure}
\begin{tikzpicture}[node distance=0.55cm, every node/.style={font=\small}]
\node[draw=popblue, fill=popblueL, rounded corners, align=center, minimum width=5.2cm, minimum height=0.9cm] (e) {\textbf{1. Einleitung}\\ Autor, Thema, Fragestellung};
\node[draw=popgreen, fill=popgreenL, rounded corners, align=center, minimum width=5.2cm, minimum height=0.9cm, below=of e] (z) {\textbf{2. Zusammenfassung}\\ die Hauptideen};
\node[draw=poporange, fill=poporangeL, rounded corners, align=center, minimum width=5.2cm, minimum height=0.9cm, below=of z] (a) {\textbf{3. Analyse}\\ Argumente, Nominalisierungen, Wortfeld};
\node[draw=poppurple, fill=poppurpleL, rounded corners, align=center, minimum width=5.2cm, minimum height=0.9cm, below=of a] (s) {\textbf{4. Eigene Meinung und Schluss}\\ Meiner Meinung nach...};
\draw[-{Stealth}, thick, popdark] (e) -- (z);
\draw[-{Stealth}, thick, popdark] (z) -- (a);
\draw[-{Stealth}, thick, popdark] (a) -- (s);
\end{tikzpicture}
\legende{Organigramme du commentaire de texte : de l'introduction à la conclusion.}
\end{popfigure}

\courssub{Un commentaire modèle rédigé}

Lisons d'abord ce modèle de commentaire (environ dix lignes), construit sur la problématique \all{Ist Afrika wirklich unabhängig?} et sur deux arguments du texte support : les prix fixés au Nord et la dépendance par les dettes.

\begin{textebox}[Modèle de commentaire : Ist Afrika wirklich unabhängig?]
Der Text behandelt den Neokolonialismus in Afrika und stellt die Frage, ob die afrikanischen Länder wirklich unabhängig sind. Der Autor erklärt, dass die politische Unabhängigkeit seit den sechziger Jahren besteht, dass aber die wirtschaftliche Abhängigkeit weiter existiert. Er nennt zwei Argumente: Erstens werden die Preise der Rohstoffe wie Kakao und Kaffee im Norden festgelegt, nicht in Afrika. Zweitens machen die Schulden die Länder des Südens abhängig von den Banken des Nordens. Auffällig sind die vielen Nominalisierungen wie „die Ausbeutung“ und „die Abhängigkeit“: Sie zeigen, dass es um Strukturen geht, nicht nur um Personen. Meiner Meinung nach ist die Ausbeutung der Rohstoffe eine Form des Neokolonialismus. Ich bin der Ansicht, dass ein Land erst frei ist, wenn es seine Preise selbst bestimmt. Ein Beispiel dafür ist der Kakao: Kamerun produziert ihn, aber andere bestimmen seinen Preis. Man kann sagen, dass die Unabhängigkeit unvollständig bleibt, solange die Wirtschaft abhängig ist.
\end{textebox}

\begin{definition}[Glossaire du commentaire]
\begin{itemize}
\item \all{festlegen} (verbe à particule séparable) : fixer, déterminer.
\item \all{die Schuld, -en} : la dette (aussi « la faute »).
\item \all{auffällig} : frappant, remarquable.
\item \all{die Struktur, -en} : la structure.
\item \all{bestimmen} : déterminer, fixer.
\item \all{unvollständig} : incomplet.
\item \all{solange} : tant que.
\end{itemize}
\end{definition}

\textbf{Questions de compréhension.}
\begin{enumerate}
\item Welche zwei Argumente nennt der Autor? (Quels sont les deux arguments de l'auteur ?)
\item Welche Nominalisierungen findest du im Kommentar? (Quelles nominalisations trouves-tu dans le commentaire ?)
\item Wie begründet der Schreiber seine eigene Meinung? (Comment le rédacteur justifie-t-il son avis personnel ?)
\end{enumerate}

\begin{exemplebox}[Des phrases clés traduites]
\all{Meiner Meinung nach ist die Ausbeutung der Rohstoffe eine Form des Neokolonialismus.} « Selon moi, l'exploitation des matières premières est une forme de néocolonialisme. »

\all{Ich bin der Ansicht, dass ein Land erst frei ist, wenn es seine Preise selbst bestimmt.} « Je suis d'avis qu'un pays n'est libre que lorsqu'il détermine lui-même ses prix. »

\all{Ein Beispiel dafür ist der Kakao.} « Un exemple de cela est le cacao. »
\end{exemplebox}

\courssub{Les expressions pour donner son avis}

Voici le tableau des expressions d'opinion à connaître par cœur. Observez bien la place du verbe dans chaque structure.

\begin{center}
\begin{tabularx}{\linewidth}{|X|X|X|}
\hline
\rowcolor{popblueL} Deutsch & Français & Remarque \\
\hline
\all{Meiner Meinung nach ist...} & Selon moi... / À mon avis... & verbe en 2\textsuperscript{e} position \\
\hline
\all{Ich bin der Ansicht, dass...} & Je suis d'avis que... & \all{dass} + verbe à la fin \\
\hline
\all{Ich denke / glaube, dass...} & Je pense / crois que... & \all{dass} + verbe à la fin \\
\hline
\all{Man kann sagen, dass...} & On peut dire que... & avis général, impersonnel \\
\hline
\all{Ein Beispiel dafür ist...} & Un exemple de cela est... & pour illustrer un argument \\
\hline
\all{Meiner Ansicht nach...} & De mon point de vue... & variante avec génitif \\
\hline
\end{tabularx}
\end{center}

\begin{attention}[Après \all{Meiner Meinung nach}, le verbe reste en 2\textsuperscript{e} position]
L'erreur classique du francophone est de calquer « selon moi, je pense que... ». En allemand, \all{Meiner Meinung nach} occupe la première place : le verbe conjugué vient immédiatement après, puis le sujet. On écrit \all{Meiner Meinung nach ist die Ausbeutung eine Form des Neokolonialismus.} et jamais \all{Meiner Meinung nach, ich denke, die Ausbeutung ist...}. De même : \all{Meiner Meinung nach sind die Preise unfair.} (« Selon moi, les prix sont injustes. »). Une seule idée, un seul verbe en deuxième position.
\end{attention}

\courssub{Production écrite guidée : un paragraphe d'opinion}

\begin{methode}[Un bon paragraphe d'opinion en cinq phrases]
\begin{enumerate}
\item \textbf{Une phrase de position} : \all{Meiner Meinung nach ist Kamerun politisch unabhängig, aber wirtschaftlich abhängig.}
\item \textbf{Un premier argument} : \all{Die Preise der Rohstoffe werden im Norden festgelegt.}
\item \textbf{Un deuxième argument} : \all{Wegen der Schulden bleibt das Land von ausländischen Banken abhängig.}
\item \textbf{Un exemple} : \all{Ein Beispiel dafür ist der Kakao aus der Region Centre.}
\item \textbf{Une phrase de conclusion} : \all{Man kann sagen, dass erst die wirtschaftliche Unabhängigkeit ein Land wirklich frei macht.}
\end{enumerate}
\end{methode}

\begin{exoresolu}[Rédiger un paragraphe : \all{Kamerun: unabhängig, aber abhängig?}]
Énoncé : rédigez un paragraphe de 6 à 8 phrases intitulé \all{Kamerun: unabhängig, aber abhängig?} en utilisant au moins cinq mots du lexique de la leçon, une nominalisation et une préposition suivie du génitif.
\tcblower
Solution détaillée — modèle de production :

\all{Kamerun ist seit 1960 politisch unabhängig, aber ist es auch wirtschaftlich frei? Meiner Meinung nach ist die Abhängigkeit des Landes noch groß. Die Preise der Rohstoffe, zum Beispiel des Kakaos und des Kaffees, werden im Norden festgelegt. Wegen der Schulden muss Kamerun oft die Bedingungen ausländischer Banken akzeptieren. Die Ausbeutung der Rohstoffe bringt wenig Geld für die Bauern, obwohl sie hart arbeiten. Ein Beispiel dafür ist der Kakao: Die Produktion ist kamerunisch, aber der Gewinn geht oft ins Ausland. Ich bin der Ansicht, dass die Industrialisierung des Landes eine Lösung wäre. Man kann sagen, dass die Unabhängigkeit erst vollständig ist, wenn Kamerun seine Produkte selbst verarbeitet.}

Vérification des contraintes : cinq mots du lexique (\all{die Abhängigkeit, der Rohstoff (als Rohstoffe), die Ausbeutung, die Unabhängigkeit, der Kolonialismus}) ; nominalisations formées (\all{die Abhängigkeit, die Ausbeutung, die Produktion, die Industrialisierung}) ; préposition + génitif (\all{wegen der Schulden}). La structure suit le plan : position, deux arguments, exemple, conclusion.
\end{exoresolu}

\begin{exercice}[À vous d'écrire]
Rédigez votre propre paragraphe de 6 à 8 phrases sur le même sujet. Contraintes : au moins cinq mots du lexique (\all{die Kolonie, der Kolonialismus, die Abhängigkeit, die Ausbeutung, die Unabhängigkeit, der Rohstoff}), une nominalisation formée par vous-même (par exemple \all{produzieren} $\rightarrow$ \all{die Produktion}) et une préposition avec le génitif (\all{wegen, trotz, während, statt}).
\end{exercice}

\begin{corrige}[Exercice : production écrite]
Éléments de correction : la position doit apparaître dès la première phrase ; chaque argument doit être introduit par un connecteur (\all{erstens, zweitens, außerdem}) ; le génitif doit être correctement marqué (\all{die Preise der Rohstoffe}, \all{wegen der Schulden}, \all{trotz der Unabhängigkeit}) ; la nominalisation doit porter l'article et la majuscule ; la conclusion doit répondre à la question du titre. Exemple de phrase attendue : \all{Trotz der Unabhängigkeit bleibt die Ausbeutung der Rohstoffe ein Problem.} (« Malgré l'indépendance, l'exploitation des matières premières reste un problème. »)
\end{corrige}

\courssub{Production orale : le débat en binômes}

\begin{experience}[Débat : \all{Ist der Handel zwischen Europa und Afrika fair?}]
Par binômes, l'un défend la position « oui, le commerce est juste », l'autre « non, il est injuste ». Déroulement guidé :
\begin{enumerate}
\item Préparation (3 minutes) : chaque élève note deux arguments et un exemple, avec au moins une expression d'avis et une préposition + génitif.
\item Prise de parole (1 minute chacun) : \all{Meiner Meinung nach ist der Handel fair / unfair, weil...} — \all{Ein Beispiel dafür ist...}
\item Réplique (30 secondes chacun) : répondre à l'argument du partenaire : \all{Ich bin nicht der Ansicht, dass..., denn...}
\item Conclusion commune : formuler ensemble une phrase de synthèse avec \all{Man kann sagen, dass...}.
\end{enumerate}
La classe vote ensuite pour le binôme le plus convaincant, en justifiant en allemand : \all{Die Argumente waren überzeugend, weil...} (« Les arguments étaient convaincants parce que... »).
\end{experience}

\begin{textebox}[Modèle de conclusion]
\all{Zusammenfassend kann man sagen, dass die politische Unabhängigkeit nicht genug ist: Erst die wirtschaftliche Unabhängigkeit macht ein Land wirklich frei.}

« En résumé, on peut dire que l'indépendance politique ne suffit pas : seule l'indépendance économique rend un pays vraiment libre. » Remarquez la nominalisation \all{die Unabhängigkeit} et l'ordre des mots après \all{erst} : le verbe \all{macht} reste en deuxième position.
\end{textebox}

\courssub{Critères d'évaluation de la production}

\begin{center}
\begin{tabularx}{\linewidth}{|X|X|}
\hline
\rowcolor{popblueL} Critère & Ce que le correcteur attend \\
\hline
Exactitude grammaticale & génitif correct (\all{der Rohstoffe}, \all{wegen der Schulden}, \all{trotz der Unabhängigkeit}), nominalisation avec article et majuscule, verbe en 2\textsuperscript{e} position \\
\hline
Richesse du lexique & au moins cinq mots du champ lexical de la leçon, employés avec le bon article \\
\hline
Structure du paragraphe & position + deux arguments + exemple + conclusion, connecteurs variés \\
\hline
Pertinence de l'argumentation & les arguments répondent à la problématique et l'exemple les illustre \\
\hline
\end{tabularx}
\end{center}

\begin{savaistu}[La stratégie payante au baccalauréat]
Au Bac, la production écrite est notée sur trois critères : la communication (le message est-il clair et structuré ?), le vocabulaire (est-il riche et adapté au thème ?) et la grammaire (les structures sont-elles correctes ?). Réutiliser les structures et les mots du texte support — les nominalisations, les prépositions avec le génitif, les expressions d'opinion — est une stratégie payante : vous montrez que vous savez transférer ce que vous avez compris. Un paragraphe simple mais correct vaut toujours mieux qu'une phrase compliquée et fautive.
\end{savaistu}

\begin{aretenir}[L'essentiel de la section]
Un commentaire de texte se construit en quatre parties : introduction, résumé, analyse, conclusion avec avis personnel. Pour donner son avis : \all{Meiner Meinung nach} + verbe en 2\textsuperscript{e} position ; \all{Ich bin der Ansicht, dass...} + verbe à la fin. Un bon paragraphe d'opinion = une phrase de position + deux arguments + un exemple + une conclusion. Les nominalisations et les prépositions avec le génitif, étudiées dans cette leçon, sont les outils grammaticaux qui donnent à votre production son niveau de Terminale.
\end{aretenir}

\newpage

\coursec{Activité d'intégration}

\coursec{Le néocolonialisme}

\courssub{Objectifs de la leçon}
À l'issue de cette leçon, l'élève sera capable de :
\begin{itemize}
\item comprendre un texte authentique traitant du néocolonialisme en Afrique (compréhension globale et détaillée) ;
\item mobiliser le \cle{lexique thématique} (\all{die Kolonie, der Kolonialismus, die Abhängigkeit, die Ausbeutung, die Unabhängigkeit, der Rohstoff, die Zusammenarbeit, die Verarbeitung}) avec articles, pluriels et familles de mots ;
\item maîtriser la \cle{nominalisation} (transformation verbes/adjectifs $\rightarrow$ noms, majuscule, article, complément au génitif) pour accéder au registre formel ;
\item employer correctement le \cle{génitif} (déclinaison noms, articles, adjectifs) et les \cle{prépositions avec le génitif} (\all{trotz, wegen, aufgrund, während}) ;
\item produire un commentaire structuré, un résumé et une prise de position argumentée (écrite puis orale) sur le thème.
\end{itemize}

\courssub{Situation de communication}
Le Parlement du Cameroun débat d'un nouveau contrat : un grand groupe industriel européen, la \all{SchokoWelt AG}, propose d'acheter chaque année 80\,\% du cacao camerounais à prix fixe, de le transformer en Europe et de revendre le chocolat fini — y compris au Cameroun. Une partie des députés y voit une \all{Zusammenarbeit} profitable ; l'autre dénonce un nouveau visage du \all{Neokolonialismus}. Le débat est retransmis à la radio et les élèves du lycée de Yaoundé le rejouent en classe d'allemand.

\courssub{Texte support : Der Neokolonialismus in Afrika}

\begin{textebox}[Pressemitteilung der SchokoWelt AG (extrait)]
„Die SchokoWelt AG bietet der Regierung Kameruns eine langfristige Zusammenarbeit an: Wir kaufen 80\,\% des kamerunischen Kakaos zu einem garantierten Preis. Trotz der Schwankungen des Weltmarktes bleibt dieser Preis stabil. Wegen unserer modernen Fabriken in Europa können wir den Rohstoff schnell verarbeiten. Aufgrund dieser Partnerschaft erhält Kamerun jedes Jahr sichere Einnahmen und neue Arbeitsplätze im Hafen von Douala. Die Verarbeitung des Kakaos im eigenen Land ist zurzeit leider nicht rentabel.“

\emph{Glossaire :} \all{die Schwankung, -en} = la fluctuation ; \all{die Einnahme, -n} = la recette, le revenu ; \all{der Arbeitsplatz, die Arbeitsplätze} = l'emploi ; \all{rentabel} = rentable ; \all{der Weltmarkt, die Weltmärkte} = le marché mondial.
\end{textebox}

\begin{textebox}[Réaction d'un député camerounais (extrait)]
„Seit der Unabhängigkeit exportieren wir unseren Rohstoff und importieren teure Schokolade. Das ist keine Zusammenarbeit, das ist die Fortsetzung der Ausbeutung! Wegen dieser Abhängigkeit bleibt unser Land arm. Wir fordern die Verarbeitung des Kakaos in Kamerun und die Gründung eigener Fabriken.“

\emph{Glossaire :} \all{die Fortsetzung, -en} = la continuation ; \all{fordern} = exiger ; \all{die Gründung, -en} = la création ; \all{eigen} = propre.
\end{textebox}

\courssub{Compréhension du texte et lexique}

\begin{exercice}[Travail sur les documents]
\begin{enumerate}
\item \textbf{Lecture active.} Lis les deux documents. Souligne en \textbf{rouge} tous les mots du lexique du néocolonialisme (\all{Abhängigkeit, Ausbeutung, Rohstoff, Unabhängigkeit, Zusammenarbeit, Verarbeitung, Kolonialismus, Neokolonialismus}) et en \textbf{bleu} toutes les prépositions suivies du génitif (\all{trotz, wegen, aufgrund, während}).
\item \textbf{Questions de compréhension (réponses en allemand, 2 à 3 phrases).}
      \begin{enumerate}
      \item \all{Was verspricht die SchokoWelt AG der Regierung Kameruns?}
      \item \all{Was kritisiert der Abgeordnete an diesem Angebot?}
      \item \all{Welche Forderungen stellt der Abgeordnete?}
      \end{enumerate}
\item \textbf{Travail lexical.} Complète le tableau ci-dessous à partir des textes et de tes connaissances.
\end{enumerate}
\end{exercice}

\begin{corrige}[Exercice 1 -- Compréhension et lexique]
\begin{enumerate}
\item \textbf{Mots soulignés en rouge :} \all{Abhängigkeit, Ausbeutung, Rohstoff, Unabhängigkeit, Zusammenarbeit, Verarbeitung} (dans les deux textes). \textbf{Prépositions + génitif en bleu :} \all{Trotz der Schwankungen, Trotz der Schwierigkeiten, Wegen unserer modernen Fabriken, Wegen dieser Abhängigkeit, Aufgrund dieser Partnerschaft}.
\item \textbf{Réponses modèles :}
      \begin{enumerate}
      \item \all{Die SchokoWelt AG verspricht eine langfristige Zusammenarbeit, den Kauf von 80\,\% des Kakaos zu einem garantierten Preis, sichere Einnahmen und neue Arbeitsplätze im Hafen von Douala.}
      \item \all{Der Abgeordnete kritisiert, dass Kamerun seit der Unabhängigkeit Rohstoffe exportiert und teure Schokolade importiert. Er nennt das keine Zusammenarbeit, sondern Fortsetzung der Ausbeutung.}
      \item \all{Er fordert die Verarbeitung des Kakaos in Kamerun und die Gründung eigener Fabriken.}
      \end{enumerate}
\item \textbf{Tableau lexical complété :}
\begin{center}
\begin{tabularx}{\linewidth}{|l|l|X|}
\hline
\rowcolor{popblueL} \textbf{Nom (article + pluriel)} & \textbf{Verbe / Adjectif} & \textbf{Sens français} \\
\hline
\all{die Kolonie, -n} & \all{kolonisieren} & la colonie / coloniser \\
\hline
\all{der Kolonialismus (Sg.)} & \all{kolonial} & le colonialisme / colonial \\
\hline
\all{der Neokolonialismus (Sg.)} & -- & le néocolonialisme \\
\hline
\all{die Abhängigkeit, -en} & \all{abhängig sein von + Dat.} & la dépendance / dépendre de \\
\hline
\all{die Ausbeutung (Sg.)} & \all{ausbeuten} & l'exploitation / exploiter \\
\hline
\all{die Unabhängigkeit (Sg.)} & \all{unabhängig} & l'indépendance / indépendant \\
\hline
\all{der Rohstoff, -e} & -- & la matière première \\
\hline
\all{die Zusammenarbeit (Sg.)} & \all{zusammenarbeiten} & la coopération / coopérer \\
\hline
\all{die Verarbeitung, -en} & \all{verarbeiten} & la transformation / transformer \\
\hline
\all{die Gründung, -en} & \all{gründen} & la création, la fondation / fonder \\
\hline
\all{der Weltmarkt, die Weltmärkte} & -- & le marché mondial \\
\hline
\end{tabularx}
\end{center}
\end{enumerate}
\end{corrige}

\courssub{Grammaire 1 : La nominalisation (die Nominalisierung)}

\begin{propriete}[Règle : La nominalisation]
La \cle{nominalisation} consiste à transformer un verbe ou un adjectif en nom. En allemand, ce nom prend obligatoirement une \textbf{majuscule} et un \textbf{article} (souvent \all{das} pour les infinitifs substantivés, \all{die} pour les noms en \all{-ung, -heit, -keit, -schaft, -ion, -tät}).
\begin{itemize}
\item \textbf{Verbe $\rightarrow$ Nom :} \all{verarbeiten} (transformer) $\rightarrow$ \all{die Verarbeitung} ; \all{ausbeuten} (exploiter) $\rightarrow$ \all{die Ausbeutung} ; \all{gründen} (fonder) $\rightarrow$ \all{die Gründung}.
\item \textbf{Adjectif $\rightarrow$ Nom :} \all{abhängig} (dépendant) $\rightarrow$ \all{die Abhängigkeit} ; \all{unabhängig} $\rightarrow$ \all{die Unabhängigkeit}.
\item \textbf{Complément d'objet $\rightarrow$ Génitif :} Le complément d'objet direct (accusatif) du verbe devient complément du nom au \textbf{génitif}.
      \all{Man verarbeitet den Kakao.} (accusatif) $\rightarrow$ \all{die Verarbeitung \textbf{des Kakaos}}. (génitif)
      \all{Man exportiert den Rohstoff.} $\rightarrow$ \all{der Export \textbf{des Rohstoffs}}.
\end{itemize}
\emph{Usage :} La nominalisation marque le style formel, administratif, journalistique ou parlementaire (écrit soutenu, discours).
\end{propriete}

\begin{exemplebox}[Exemples de nominalisation (textes support)]
\begin{itemize}
\item \all{Wir verarbeiten den Kakao.} $\rightarrow$ \all{Die Verarbeitung des Kakaos} (texte 1 et 2)
\item \all{Wir kooperieren / arbeiten zusammen.} $\rightarrow$ \all{die Zusammenarbeit} (texte 1)
\item \all{Wir sind abhängig.} $\rightarrow$ \all{die Abhängigkeit} (texte 2)
\item \all{Wir sind unabhängig.} $\rightarrow$ \all{die Unabhängigkeit} (texte 2)
\item \all{Wir gründen Fabriken.} $\rightarrow$ \all{die Gründung eigener Fabriken} (texte 2)
\end{itemize}
\end{exemplebox}

\begin{exercice}[Entraînement : Nominaliser]
Transforme les phrases en remplaçant le verbe souligné par une nominalisation (sujet = \all{das} ou \all{die} + nom). Mets le complément au génitif.
\begin{enumerate}
\item Kamerun \textbf{exportiert} den Kakao.
\item Die Firmen \textbf{ausbeuten} die Arbeiter.
\item Wir \textbf{verarbeiten} den Rohstoff im Land.
\item Das Land \textbf{ist unabhängig} geworden.
\end{enumerate}
\end{exercice}

\begin{corrige}[Exercice 2 -- Nominalisation]
\begin{enumerate}
\item \all{Der Export des Kakaos durch Kamerun} (ou \all{Kameruns Export des Kakaos}).
\item \all{Die Ausbeutung der Arbeiter durch die Firmen}.
\item \all{Die Verarbeitung des Rohstoffs im Land}.
\item \all{Die Unabhängigkeit des Landes} (ou \all{die Erlangung der Unabhängigkeit}).
\end{enumerate}
\end{corrige}

\courssub{Grammaire 2 : Le génitif et les prépositions avec le génitif}

\begin{propriete}[Rappel : Déclinaison du génitif]
Le génitif marque l'appartenance, la possession ou le complément du nom nominalisé.
\begin{center}
\begin{tabular}{|l|c|c|c|}
\hline
\rowcolor{popblueL} & \textbf{Masculin} & \textbf{Féminin} & \textbf{Neutre} / \textbf{Pluriel} \\
\hline
\textbf{Article défini} & \all{des} + \all{-(e)s} & \all{der} & \all{des} / \all{der} \\
\hline
\textbf{Article indéfini} & \all{eines} + \all{-(e)s} & \all{einer} & \all{eines} / -- \\
\hline
\textbf{Adjectif (faible)} & \all{-en} & \all{-en} & \all{-en} / \all{-en} \\
\hline
\end{tabular}
\end{center}
\emph{Exemples :} \all{des Kakaos, des Rohstoffs, des Landes} (masc./neutre + \all{-s}) ; \all{der Unabhängigkeit, der Abhängigkeit, der Verarbeitung} (fém.) ; \all{der Fabriken, der Schwierigkeiten} (pluriel).
\end{propriete}

\begin{propriete}[Les quatre prépositions régissant le génitif]
Ces prépositions placent \textbf{obligatoirement} leur complément au génitif. Elles sont fréquentes dans le style formel.
\begin{center}
\begin{tabularx}{\linewidth}{|l|l|X|}
\hline
\rowcolor{popblueL} \textbf{Préposition} & \textbf{Sens} & \textbf{Exemple (texte ou construit)} \\
\hline
\all{trotz} & malgré & \all{Trotz der Unabhängigkeit bleibt die Abhängigkeit groß.} \\
\hline
\all{wegen} & à cause de & \all{Wegen dieser Abhängigkeit verlieren wir Geld.} \\
\hline
\all{aufgrund} & en raison de, suite à & \all{Aufgrund dieser Partnerschaft gibt es neue Arbeitsplätze.} \\
\hline
\all{während} & pendant & \all{Während der Debatte gab es Tumulte.} \\
\hline
\end{tabularx}
\end{center}
\emph{Attention :} \all{wegen} + génitif (standard soutenu). À l'oral familier, on entend \all{wegen + datif} (\all{wegen dem Wetter}), mais \textbf{à l'écrit et à l'examen, seul le génitif est correct}.
\end{propriete}

\begin{attention}[Pièges fréquents pour francophones]
\begin{itemize}
\item \textbf{*wegen die Abhängigkeit} $\rightarrow$ \textbf{Faux.} Après \all{wegen} : \all{wegen \textbf{der} Abhängigkeit} (génitif fém.).
\item \textbf{*trotz der Unabh\"angigkeit} $\rightarrow$ \textbf{Orthographe :} \all{Unabhängigkeit} (Umlaut \all{ä}).
\item \textbf{Majuscule oubliée :} \all{*die verarbeitung} $\rightarrow$ \textbf{\all{die Verarbeitung}}. Tout nom (nominalisé ou non) prend une majuscule.
\item \textbf{Place du verbe dans \all{dass}/\all{weil} :} \all{Ich bin der Ansicht, *dass dieser Vertrag ungerecht ist.} $\rightarrow$ \all{..., \textbf{dass dieser Vertrag ungerecht \underline{ist}}.} (verbe à la fin).
\item \textbf{Génitif des noms en \all{-s/-ß/-x/-z} :} Ils prennent \all{-es} : \all{der Weltmarkt} $\rightarrow$ \all{des Weltmarkt\textbf{es}} ; \all{der Preis} $\rightarrow$ \all{des Preises}.
\end{itemize}
\end{attention}

\begin{exercice}[Entraînement : Génitif et prépositions]
Complète avec la forme correcte au génitif (article + nom + adjectif si nécessaire).
\begin{enumerate}
\item \all{Trotz \_\_\_\_\_\_\_\_\_\_\_\_ (die große Schwierigkeit) bauen wir Fabriken.}
\item \all{Wegen \_\_\_\_\_\_\_\_\_\_\_\_ (der hohe Weltmarktpreis) ist der Kakao teuer.}
\item \all{Aufgrund \_\_\_\_\_\_\_\_\_\_\_\_ (der neue Vertrag) exportieren wir mehr.}
\item \all{Während \_\_\_\_\_\_\_\_\_\_\_\_ (die lange Debatte) hat niemand gegessen.}
\end{enumerate}
\end{exercice}

\begin{corrige}[Exercice 3 -- Génitif]
\begin{enumerate}
\item \all{Trotz \textbf{der großen Schwierigkeit}} (fém. sing. : \all{der} + adj. \all{-en}).
\item \all{Wegen \textbf{des hohen Weltmarktpreises}} (masc. sing. : \all{des} + \all{-es} nom, adj. \all{-en}).
\item \all{Aufgrund \textbf{des neuen Vertrags}} (masc. sing. : \all{des} + \all{-s} nom, adj. \all{-en}).
\item \all{Während \textbf{der langen Debatte}} (fém. sing. : \all{der} + adj. \all{-en}).
\end{enumerate}
\end{corrige}

\courssub{Production guidée : Débat parlementaire}

\begin{methode}[Déroulement de l'activité d'intégration (2 h)]
\begin{enumerate}
\item \textbf{Préparation écrite (30 min).} Individuellement, rédige ta prise de position (5--6 phrases) selon les consignes ci-dessous. Utilise le lexique, la nominalisation, le génitif et une formule d'opinion.
\item \textbf{Relecture croisée (10 min).} Échange ta feuille avec un voisin. Vérifie : majuscules des noms, génitif (\all{des Kakaos, der Unabhängigkeit, des Weltmarkt\textbf{es}}), place du verbe (\all{dass/weil} $\rightarrow$ verbe final), prépositions + génitif.
\item \textbf{Mise en place du débat (10 min).} La classe forme deux camps : \all{Abgeordnete} (députés camerounais) vs \all{Vertreter der SchokoWelt AG}. L'enseignant est modérateur.
\item \textbf{Débat oral (40 min).} Chaque camp présente un argument. L'adversaire répond \textbf{obligatoirement} par une formule de réaction avant de contre-argumenter. Le modérateur note au tableau les réussites (beau génitif, nominalisation correcte, mot de lexique).
\item \textbf{Conclusion individuelle (10 min).} Rédige 2 phrases personnelles : as-tu changé d'avis ? Justifie avec \all{weil} ou \all{aufgrund} + génitif.
\item \textbf{Prolongement (devoir).} Rédige un court article de presse (8--10 lignes) rendant compte du débat, avec au moins 2 nominalisations et 2 prépositions + génitif.
\end{enumerate}
\end{methode}

\begin{exercice}[Consignes de production écrite (étape 1)]
Rédige une prise de position de \textbf{5 à 6 phrases} pour ton camp. Contraintes obligatoires :
\begin{itemize}
\item Au moins \textbf{4 mots du lexique} : \all{die Abhängigkeit, die Ausbeutung, der Rohstoff, die Unabhängigkeit, die Zusammenarbeit, die Verarbeitung}.
\item Au moins \textbf{une nominalisation} avec article et génitif (ex. \all{die Verarbeitung des Kakaos im eigenen Land}).
\item Au moins \textbf{une préposition avec le génitif} : \all{trotz, wegen, aufgrund, während}.
\item Une \textbf{formule d'opinion} : \all{Ich bin der Ansicht, dass...} / \all{Meiner Meinung nach...} / \all{Ich halte es für falsch/richtig, dass...}.
\item Registre formel (pas de contractions familières, phrases complètes).
\end{itemize}
\end{exercice}

\begin{savaistu}[Formules utiles pour le débat oral (Reaktionen)]
\begin{center}
\begin{tabularx}{\linewidth}{|l|X|}
\hline
\rowcolor{popblueL} \textbf{Deutsch} & \textbf{Français} \\
\hline
\all{Ich bin der Ansicht, dass...} & Je suis d'avis que... \\
\hline
\all{Meiner Meinung nach...} & À mon avis... \\
\hline
\all{Ich bin anderer Meinung, weil...} & Je suis d'un autre avis, car... \\
\hline
\all{Trotz Ihrer Argumente glaube ich, dass...} & Malgré vos arguments, je crois que... \\
\hline
\all{Das mag stimmen, aber...} & Cela peut être vrai, mais... \\
\hline
\all{Aufgrund der Tatsachen muss man sagen, dass...} & Au vu des faits, on doit dire que... \\
\hline
\end{tabularx}
\end{center}
\end{savaistu}

\courssub{Proposition de production et corrigé commenté}

\begin{exoresolu}[Prise de position d'un député (modèle pour le camp \all{Abgeordnete})]
Rédige une prise de position de 5 à 6 phrases pour le camp des députés, avec 4 mots du lexique, une nominalisation et une préposition + génitif.
\tcblower
\textbf{Production modèle :}

\all{Sehr geehrte Damen und Herren! Seit der Unabhängigkeit exportieren wir unseren Rohstoff Kakao und importieren teure Schokolade. Meiner Meinung nach ist das keine Zusammenarbeit, sondern die Fortsetzung der Ausbeutung. Wegen dieser Abhängigkeit verlieren wir jedes Jahr viel Geld. Ich bin der Ansicht, dass die Verarbeitung des Kakaos im eigenen Land unsere Zukunft ist. Trotz der Schwierigkeiten müssen wir eigene Fabriken bauen. Nur so wird die Unabhängigkeit unseres Landes wirklich vollständig.}

\textbf{Commentaire des choix de langue :}
\begin{itemize}
\item \textbf{Lexique (5 mots) :} \all{die Unabhängigkeit, der Rohstoff, die Zusammenarbeit, die Ausbeutung, die Abhängigkeit} -- le contrat est rempli (minimum 4).
\item \textbf{Nominalisations (2) :} \all{die Fortsetzung der Ausbeutung} (verbe \all{fortsetzen} + génitif \all{der Ausbeutung}) et \all{die Verarbeitung des Kakaos im eigenen Land} (verbe \all{verarbeiten} + génitif \all{des Kakaos}, neutre en \all{-s}).
\item \textbf{Prépositions + génitif (2) :} \all{wegen dieser Abhängigkeit} (féminin, démonstratif \all{dieser}) et \all{trotz der Schwierigkeiten} (pluriel, article \all{der}).
\item \textbf{Subordonnées :} Dans \all{Ich bin der Ansicht, dass die Verarbeitung ... unsere Zukunft \underline{ist}}, le verbe \all{ist} est bien rejeté à la fin. Dans \all{weil}-phrases implicites, même règle.
\item \textbf{Style formel :} L'apostrophe \all{Sehr geehrte Damen und Herren!}, l'absence de contractions, l'usage des nominalisations et du génitif donnent le ton attendu dans un hémicycle.
\end{itemize}
\end{exoresolu}

\begin{exoresolu}[Prise de position du groupe \all{SchokoWelt AG} (modèle pour le camp industriel)]
\tcblower
\textbf{Production modèle :}

\all{Sehr geehrte Abgeordnete! Unser Angebot basiert auf einer fairen Partnerschaft. Aufgrund der garantierten Preise sichern wir die Einnahmen der Bauern. Trotz der Kritik an der Rohstoffexportierung investieren wir in den Hafen von Douala. Ich bin der Ansicht, dass die Verarbeitung in Europa derzeit effizienter ist. Wegen der fehlenden Fabriken in Kamerun schaffen wir dort Arbeitsplätze in der Logistik. Eine echte Unabhängigkeit entsteht durch wirtschaftliche Stabilität, nicht durch Isolation.}

\textbf{Commentaire :} Lexique (\all{Partnerschaft, Einnahmen, Rohstoffexport, Unabhängigkeit, Arbeitsplätze}), nominalisation (\all{die Verarbeitung in Europa}, \all{die Rohstoffexportierung}), prépositions + génitif (\all{Aufgrund der garantierten Preise, Trotz der Kritik, Wegen der fehlenden Fabriken}), structure \all{dass}-clause correcte.
\end{exoresolu}

\courssub{Bilan et prolongement}

\begin{aretenir}[Bilan des savoir-faire réinvestis]
\begin{itemize}
\item \textbf{Compréhension :} Exploitation de deux documents contradictoires (communiqué d'entreprise vs discours politique), identification des arguments, du lexique et des marqueurs grammaticaux (génitif).
\item \textbf{Lexique :} Maîtrise des familles de mots (\all{kolonisieren -- Kolonie -- Kolonialismus -- Neokolonialismus}; \all{abhängig -- Abhängigkeit}; \all{verarbeiten -- Verarbeitung}; \all{ausbeuten -- Ausbeutung}) avec genre et pluriel.
\item \textbf{Grammaire :} La nominalisation comme outil du style formel (passage verbe $\rightarrow$ nom + génitif). Le génitif comme marqueur de relation de possession/appartenance et son emploi après \all{trotz, wegen, aufgrund, während}.
\item \textbf{Expression :} Passage de l'écrit préparé (sécurité linguistique) à l'oral interactif (réactivité, écoute, reformulation). Respect des conventions du débat (formules de politesse, de réaction, d'argumentation).
\end{itemize}
\end{aretenir}

\begin{experience}[Prolongement : Article de presse (Devoir / Évaluation)]
\begin{itemize}
\item \textbf{Consigne :} Rédige un article pour le journal scolaire \all{„Yaoundé Heute“} (100--120 mots) relatant le débat.
\item \textbf{Contraintes :} Titre accrocheur ; paragraphe d'introduction (qui, quoi, quand) ; arguments des deux camps ; citation indirecte (\all{Der Abgeordnete betonte, dass...}) ; conclusion personnelle.
\item \textbf{Exigences grammaticales :} Minimum \textbf{2 nominalisations} (ex. \all{die Forderung nach Verarbeitung}, \all{die Ablehnung des Vertrags}) et \textbf{2 prépositions + génitif} (ex. \all{trotz der Einwände, aufgrund der Garantien}).
\item \textbf{Critères d'évaluation :} Respect consignes (20\,\%), correction grammaticale (30\,\%), richesse lexicale (20\,\%), cohérence argumentative (30\,\%).
\end{itemize}
\end{experience}

\newpage

\coursec{Méthodes et exercices résolus}

\coursec{Leçon AL27 : Le néocolonialisme}

\courssub{Texte support : Der Neokolonialismus in Afrika}

\begin{textebox}[Texte original — Niveau Terminale A (LV2)]
\all{Neokolonialismus bezeichnet die fortgesetzte wirtschaftliche und politische Abhängigkeit ehemals kolonialer Gebiete nach ihrer formellen Unabhängigkeit. Viele afrikanische Staaten erlangten in den 1960er Jahren die politische Souveränität, doch die ökonomischen Strukturen blieben oft unverändert: Die Rohstoffe wie Kakao, Erdöl oder Seltene Erden werden weiterhin billig in den Norden exportiert, während teure Fertigprodukte und Technologien importiert werden müssen.}

\all{Ein zentrales Instrument ist die Schuldenfalle: Internationale Finanzinstitutionen vergeben Kredite an Bedingungen, die die wirtschaftliche Entscheidungsfreiheit der Länder einschränken. Hinzu kommt der Einfluss multinationaler Konzerne, die Landrechte erwerben und Gewinne in ihre Heimatländer transferieren, statt lokale Wertschöpfungsketten aufzubauen.}

\all{Trotz dieser Herausforderungen gibt es Ansätze zur Emanzipation: Die Afrikanische Kontinentalfreihandelszone (AfCFTA) soll den innerafrikanischen Handel stärken, und einige Staaten beginnen, Rohstoffe vor Ort zu verarbeiten. Echte Unabhängigkeit erfordert jedoch nicht nur politische, sondern vor allem wirtschaftliche Souveränität.}
\end{textebox}

\begin{definition}[Vocabulaire clé du texte (avec article et pluriel)]
\begin{itemize}
\item \all{die Souveränität, -en} — la souveraineté
\item \all{die ökonomische Struktur, -en} — la structure économique
\item \all{der Rohstoff, -e} — la matière première
\item \all{der Norden} (singulier) — le Nord (pays développés)
\item \all{das Fertigprodukt, -e} — le produit fini
\item \all{die Schuldenfalle, -n} — le piège de la dette
\item \all{die Finanzinstitution, -en} — l'institution financière
\item \all{der Kredit, -e} — le crédit / le prêt
\item \all{die Entscheidungsfreiheit} — la liberté de décision
\item \all{der multinationaler Konzern, -e} — la multinationale
\item \all{der Landrecht, -e} — le droit foncier
\item \all{die Wertschöpfungskette, -n} — la chaîne de valeur ajoutée
\item \all{die Afrikanische Kontinentalfreihandelszone (AfCFTA)} — la ZLECAf (Zone de libre-échange continentale africaine)
\item \all{der innerafrikanische Handel} — le commerce intra-africain
\item \all{die Emanzipation} — l'émancipation
\end{itemize}
\end{definition}

\courssub{Compréhension du texte}

\begin{exercice}[Questions de compréhension globale et détaillée]
\textbf{Consigne} : \all{Beantworten Sie die Fragen auf Deutsch in vollständigen Sätzen. Benutzen Sie möglichst eigene Worte (Paraphrase).}
\begin{enumerate}
\item Was wird im Text unter Neokolonialismus verstanden?
\item Nennen Sie zwei Beispiele für Rohstoffe, die im Text genannt werden.
\item Erklären Sie den Mechanismus der „Schuldenfalle“.
\item Welche Rolle spielen multinationale Konzerne nach dem Text?
\item Nennen Sie eine Maßnahme, die zur wirtschaftlichen Emanzipation beitragen soll.
\item Was ist Ihrer Meinung nach die wichtigste Bedingung für „echte Unabhängigkeit“? Begründen Sie kurz.
\end{enumerate}
\end{exercice}

\begin{corrige}[Exercice — Compréhension du texte]
\begin{enumerate}
\item \all{Neokolonialismus ist die fortgesetzte wirtschaftliche und politische Abhängigkeit ehemals kolonialer Gebiete nach ihrer formellen Unabhängigkeit.}
\item \all{Beispiele für Rohstoffe sind Kakao, Erdöl und Seltene Erden.}
\item \all{Internationale Finanzinstitutionen vergeben Kredite an Bedingungen, die die wirtschaftliche Entscheidungsfreiheit der Länder einschränken.}
\item \all{Multinationale Konzerne erwerben Landrechte, transferieren Gewinne in ihre Heimatländer und bauen keine lokalen Wertschöpfungsketten auf.}
\item \all{Die Afrikanische Kontinentalfreihandelszone (AfCFTA) soll den innerafrikanischen Handel stärken, und einige Staaten verarbeiten Rohstoffe vor Ort.}
\item \all{Meiner Meinung nach ist die wirtschaftliche Souveränität die wichtigste Bedingung, weil politische Unabhängigkeit ohne ökonomische Autonomie nicht nachhaltig ist.}
\end{enumerate}
\end{corrige}

\courssub{Lexique thématique : Kolonialismus und Abhängigkeit}

\begin{methode}[Maîtriser le lexique de la colonisation et de la dépendance]
\begin{enumerate}
\item \textbf{Apprendre chaque nom avec son article, son pluriel et son genre} : 
\all{die Kolonie, -n} (f.) ; \all{der Kolonialismus} (m., pas de pluriel usuel) ; \all{die Abhängigkeit, -en} (f.) ; \all{die Ausbeutung} (f., pas de pluriel usuel) ; \all{die Unabhängigkeit} (f., pas de pluriel usuel) ; \all{der Rohstoff, -e} (m.) ; \all{die Souveränität, -en} (f.) ; \all{der Konzern, -e} (m.).
\item \textbf{Construire les familles de mots (dérivation)} :
\begin{center}
\begin{tabularx}{\linewidth}{|X|X|X|}\hline
\rowcolor{popblueL} \textbf{Verbe / Adjectif} & \textbf{Nom (dérivé)} & \textbf{Construction / Préposition} \\ \hline
\all{abhängig sein} (adj. + sein) & \all{die Abhängigkeit} & \all{abhängig von} + \cle{Datif} ; \all{die Abhängigkeit von} + \cle{Datif} \\ \hline
\all{unabhängig} (adj.) & \all{die Unabhängigkeit} & \all{unabhängig von} + \cle{Datif} \\ \hline
\all{ausbeuten} (verbe) & \all{die Ausbeutung} / \all{das Ausbeuten} & \all{die Ausbeutung} + \cle{Génitif} (obj.) \\ \hline
\all{kolonisieren} (verbe) & \all{die Kolonisierung} / \all{der Kolonialismus} & \all{kolonisieren} + \cle{Accusatif} \\ \hline
\all{exportieren / importieren} & \all{der Export} / \all{der Import} (m.) ; \all{das Exportieren} / \all{das Importieren} (n.) & \all{exportieren} + \cle{Accusatif} \\ \hline
\end{tabularx}
\end{center}
\item \textbf{Mémoriser les collocations (associations figées)} : \all{in die Unabhängigkeit führen} ; \all{aus der Abhängigkeit befreien} ; \all{Rohstoffe ausbeuten} ; \all{eine Kolonie gründen / verlieren} ; \all{wirtschaftliche Souveränität erlangen}.
\item \textbf{Contextualiser} : utiliser ce vocabulaire pour parler du Cameroun (cacao, pétrole, dette, ZLECAf).
\end{enumerate}
\end{methode}

\begin{exoresolu}[Vocabulaire — Complétion et transformation]
\textbf{Énoncé} : Complétez les phrases avec la forme correcte du mot entre parenthèses (nom, adjectif, verbe, préposition, cas).\par
a) Kamerun war bis 1916 eine deutsche \dots\ (Kolonie).\par
b) 1960 erlangte das Land seine \dots\ (Unabhängigkeit).\par
c) Die \dots\ von den Weltmarktpreisen ist riskant. (Abhängigkeit)\par
d) Die \dots\ der Bodenschätze durch Konzerne schadet der Entwicklung. (ausbeuten)\par
e) Das Land ist wirtschaftlich \dots\ von den Importen. (abhängig)\par
f) Wir müssen die \dots\ der Rohstoffe im eigenen Land fördern. (verarbeiten — nominaliser)
\tcblower
\textbf{Solution détaillée} :\par
a) \all{Kolonie} (nom, f., accusatif après \all{eine}). « Le Cameroun était une colonie allemande jusqu'en 1916. »\par
b) \all{Unabhängigkeit} (nom, f., accusatif après \all{seine}). « En 1960, le pays accéda à son indépendance. »\par
c) \all{Abhängigkeit} (nom, f., sujet). Construction : \all{die Abhängigkeit von} + Datif (\all{den Weltmarktpreisen}). « La dépendance aux prix du marché mondial est risquée. »\par
d) \all{Ausbeutung} (nom f. en \all{-ung}, sujet). Génitif de l'objet : \all{der Bodenschätze}. « L'exploitation des ressources minières par les multinationales nuit au développement. »\par
e) \all{abhängig} (adjectif prédicatif). Construction : \all{abhängig von} + Datif (\all{den Importen}). « Le pays est économiquement dépendant des importations. »\par
f) \all{Verarbeitung} (nom f. en \all{-ung}, accusatif après \all{die}). « Nous devons promouvoir la transformation des matières premières dans le propre pays. »
\end{exoresolu}

\begin{savaistu}[Le Cameroun et l'histoire coloniale allemande]
Le Cameroun a été un \all{Schutzgebiet} (protectorat) allemand de 1884 à 1916 (\all{die Kolonialzeit}). Après la Première Guerre mondiale, il a été partagé entre la France et la Grande-Bretagne sous mandat de la SDN (\all{Völkerbundsmandat}). L'héritage linguistique (allemand langue étrangère au lycée), architectural (bâtiments à Yaoundé, Douala, Buea) et infrastructurel (chemin de fer \all{Nordbahn}) est encore visible. Le mot \all{Kolonie} vient du latin \all{colonia} « établissement agricole, colonie » (\all{colere} « cultiver »).
\end{savaistu}

\courssub{Grammaire 1 : La nominalisation (die Nominalisierung)}

\begin{propriete}[Les trois procédés principaux de nominalisation]
La nominalisation permet de transformer un verbe ou un adjectif en nom pour obtenir un style écrit soutenu, compact et précis (typique des résumés, analyses, articles de presse).
\begin{enumerate}
\item \textbf{L'infinitif nominalisé (toujours neutre : \all{das})} : \all{das Ausbeuten}, \all{das Exportieren}, \all{das Verarbeiten}, \all{das Lesen}. \emph{Usage} : action abstraite, générale. Souvent sans complément d'objet.
\item \textbf{Le suffixe \all{-ung} (toujours féminin : \all{die})} : \all{die Ausbeutung}, \all{die Entwicklung}, \all{die Regierung}, \all{die Verarbeitung}, \all{die Kolonisierung}. \emph{Usage} : action concrète, processus, résultat. Le complément d'objet devient \cle{génitif} (\all{die Ausbeutung der Rohstoffe}) ou prépositionnel (\all{die Ausbeutung an den Rohstoffen} — rare).
\item \textbf{Les suffixes \all{-heit} / \all{-keit} (toujours féminin : \all{die})} : 
\begin{itemize}
\item \all{-heit} sur adjectifs courts / radicaux monosyllabiques : \all{die Freiheit}, \all{die Gesundheit}, \all{die Rohheit}.
\item \all{-keit} sur adjectifs en \all{-bar}, \all{-ig}, \all{-lich}, \all{-sam}, \all{-los} et participes passés adjectivés : \all{die Abhängigkeit}, \all{die Unabhängigkeit}, \all{die Möglichkeit}, \all{die Wichtigkeit}, \all{die Hilflosigkeit}.
\end{itemize}
\item \textbf{Autres nominalisations (conversion, suffixes masculins)} : \all{der Export}, \all{der Import}, \all{der Aufbau}, \all{der Abbau}, \all{die Einfuhr}, \all{die Ausfuhr}. Genre à apprendre par cœur.
\end{enumerate}
\end{propriete}

\begin{methode}[Transformer une phrase en structure nominalisée (style écrit)]
\begin{enumerate}
\item Identifier le verbe principal ou l'adjectif clé.
\item Choisir le procédé adapté (contexte : processus \rightarrow \all{-ung} ; qualité/état \rightarrow \all{-keit/heit} ; action générale \rightarrow Infinitif).
\item \textbf{Transformer les compléments} :
\begin{itemize}
\item Sujet de la phrase active $\rightarrow$ \all{von} + Datif (agent) ou Génitif (possesseur) : \all{Die Länder exportieren} $\rightarrow$ \all{der Export der Länder} / \all{der Export von Seiten der Länder}.
\item Objet (accusatif) $\rightarrow$ \cle{Génitif} : \all{exportieren die Rohstoffe} $\rightarrow$ \all{der Export der Rohstoffe}.
\item Complément prépositionnel $\rightarrow$ conserve souvent sa préposition : \all{abhängig von den Märkten} $\rightarrow$ \all{die Abhängigkeit von den Märkten}.
\end{itemize}
\item Mettre la \cle{majuscule} au nominalisé. Vérifier l'article (\all{das}, \all{die}, \all{der}).
\end{enumerate}
\end{methode}

\begin{exoresolu}[Entraînement : Nominalisation guidée]
\textbf{Énoncé} : Nominalisez les éléments soulignés et réécrivez la phrase complète.\par
a) \all{Die afrikanischen Länder \textbf{sind politisch unabhängig}, aber \textbf{hängen wirtschaftlich ab}.}\par
b) \all{\textbf{Weil die Konzerne die Rohstoffe \textbf{ausbeuten}}, bleibt das Land arm.}\par
c) \all{\textbf{Das Land \textbf{verarbeitet} seinen Kakao nicht selbst.}}\par
d) \all{\textbf{Während der Kolonialzeit \textbf{wurden} viele Menschen \textbf{unterdrückt}.}}
\tcblower
\textbf{Solution et raisonnement} :\par
a) Adjectif \all{unabhängig} $\rightarrow$ \all{die Unabhängigkeit} (Génitif : \all{des Landes} / \all{der Länder}). Verbe \all{abhängen} (séparable) $\rightarrow$ \all{die Abhängigkeit}. \all{Die politische Unabhängigkeit der afrikanischen Länder steht im Gegensatz zu ihrer wirtschaftlichen Abhängigkeit.} (Style nominalisé : sujet = noms).\par
b) Verbe \all{ausbeuten} $\rightarrow$ \all{die Ausbeutung} (f.). Sujet \all{die Konzerne} $\rightarrow$ Génitif \all{der Konzerne} (ou \all{von Seiten der Konzerne}). Objet \all{die Rohstoffe} $\rightarrow$ Génitif \all{der Rohstoffe}. \all{Die Ausbeutung der Rohstoffe durch die Konzerne lässt das Land arm bleiben.} (\all{durch} + Accusatif pour l'agent passif).\par
c) Verbe \all{verarbeiten} $\rightarrow$ \all{die Verarbeitung} (f.) ou \all{das Verarbeiten} (n.). \all{Die fehlende Verarbeitung des Kakaos im eigenen Land hindert die Entwicklung.} (Négation par adjectif \all{fehlend} + Génitif \all{des Kakaos}).\par
d) Verbe \all{unterdrücken} $\rightarrow$ \all{die Unterdrückung} (f.). Agent \all{viele Menschen} (objet accusatif original) devient sujet du passif $\rightarrow$ Génitif \all{vieler Menschen} (ou \all{der Menschen}). Temps \all{wurden unterdrückt} (Präteritum Passiv) $\rightarrow$ nom \all{Unterdrückung}. \all{Die Unterdrückung vieler Menschen während der Kolonialzeit ist ein dunkles Kapitel.} (\all{während} + Génitif \all{der Kolonialzeit}).
\end{exoresolu}

\begin{attention}[Pièges fréquents de la nominalisation]
\begin{enumerate}
\item \textbf{Majuscule oubliée} : \all{die abhängigkeit} $\times$ $\rightarrow$ \all{die Abhängigkeit} $\checkmark$. \textbf{Tout nom s'écrit avec une majuscule.}
\item \textbf{Mauvais suffixe} : \all{die Abhängung} $\times$ (n'existe pas) $\rightarrow$ \all{die Abhängigkeit} $\checkmark$ (adj. \all{abhängig} + \all{-keit}). \all{die Entwickelung} $\times$ $\rightarrow$ \all{die Entwicklung} $\checkmark$.
\item \textbf{Génitif manquant ou mal formé} : \all{die Ausbeutung die Rohstoffe} $\times$ $\rightarrow$ \all{die Ausbeutung \textbf{der} Rohstoffe} $\checkmark$ (pluriel : article \all{der}, nom sans \all{-n} si pas de \all{-n} au pluriel nominatif/accusatif ; ici \all{Rohstoffe} $\rightarrow$ \all{der Rohstoffe}).
\item \textbf{Confusion genre/article} : \all{der Abhängigkeit} $\times$ $\rightarrow$ \all{\textbf{die} Abhängigkeit} $\checkmark$ (\all{-keit} = fém.). \all{das Export} $\times$ $\rightarrow$ \all{\textbf{der} Export} $\checkmark$ (masc.) ; \all{das Exportieren} $\checkmark$ (neutre).
\item \textbf{Ordre des mots dans le groupe nominal} : Déterminant + Adjectifs + Nom noyau + Compléments (Génitif / Prépositionnel). Ex: \all{die wirtschaftliche Abhängigkeit \textbf{von den Märkten}}.
\end{enumerate}
\end{attention}

\courssub{Grammaire 2 : Le génitif et les prépositions avec le génitif}

\begin{propriete}[Formation du génitif (Singulier \& Pluriel)]
Le génitif marque l'appartenance, la possession ou la relation entre deux noms (N1 = possédé, N2 = possesseur).
\begin{center}
\begin{tabular}{|l|c|c|c|c|}\hline
\rowcolor{popblueL} \textbf{Cas} & \textbf{Masculin} & \textbf{Féminin} & \textbf{Neutre} & \textbf{Pluriel} \\ \hline
\textbf{Article défini} & \all{des} & \all{der} & \all{des} & \all{der} \\ \hline
\textbf{Article indéfini} & \all{eines} & \all{einer} & \all{eines} & \all{—} (keiner) \\ \hline
\textbf{Nom (terminaison)} & \all{-(e)s} & \all{—} & \all{-(e)s} & \all{—} (souvent \all{-n} si pluriel en \all{-n}) \\ \hline
\textbf{Exemple} & \all{des Landes} & \all{der Unabhängigkeit} & \all{des Rohstoffs} & \all{der Rohstoffe} \\ \hline
\end{tabular}
\end{center}
\emph{Règle des \all{-(e)s}} : Monosyllabes $\rightarrow$ \all{-es} (\all{des Landes}, \all{des Mannes}) ; Polysyllabes $\rightarrow$ \all{-s} (\all{des Computers}, \all{des Exportes}). Noms en \all{-s, -ß, -x, -z} $\rightarrow$ \all{-es} (\all{des Preises}).
\end{propriete}

\begin{propriete}[Prépositions régissant le génitif (Registre soutenu / Écrit / Bac)]
\begin{center}
\begin{tabularx}{\linewidth}{|X|X|X|}\hline
\rowcolor{popblueL} \textbf{Préposition} & \textbf{Sens} & \textbf{Exemple (contexte leçon)} \\ \hline
\all{trotz} & malgré & \all{trotz der Unabhängigkeit} (malgré l'indépendance) \\ \hline
\all{während} & pendant / alors que & \all{während der Kolonialzeit} (pendant la période coloniale) \\ \hline
\all{wegen} & à cause de & \all{wegen der Abhängigkeit} (à cause de la dépendance) \\ \hline
\all{statt / anstatt} & au lieu de & \all{statt der Ausbeutung} (au lieu de l'exploitation) \\ \hline
\all{innerhalb} & à l'intérieur de (délai/espace) & \all{innerhalb eines Jahres} (en un an / dans l'année) \\ \hline
\all{außerhalb} & à l'extérieur de & \all{außerhalb der Zone} (hors de la zone) \\ \hline
\end{tabularx}
\end{center}
\textbf{Attention registre} : À l'oral familier, \all{wegen, trotz} prennent souvent le datif (\all{wegen dem Wetter}). \textbf{À l'écrit (Bac, résumé, analyse), le génitif est obligatoire.}
\end{propriete}

\begin{popfigure}
\begin{tikzpicture}[
  node distance=1.2cm and 1.5cm,
  prep/.style={draw=popblue, thick, fill=popblueL, rounded corners, minimum width=2.2cm, minimum height=0.9cm, align=center, font=\sffamily\bfseries},
  case/.style={draw=poporange, thick, fill=poporangeL, rounded corners, minimum width=2.5cm, minimum height=0.9cm, align=center, font=\sffamily},
  arrow/.style={-Stealth, thick, popdark}
]
\node[prep] (trotz) {trotz};
\node[prep, right=of trotz] (waehrend) {während};
\node[prep, right=of waehrend] (wegen) {wegen};
\node[prep, right=of wegen] (statt) {statt / anstatt};
\node[prep, below=of waehrend] (innerhalb) {innerhalb};
\node[prep, right=of innerhalb] (ausserhalb) {außerhalb};

\node[case, below=1.5cm of trotz, align=center] (gen){GÉNITIF\\ \all{der / des / der / der}};

\draw[arrow] (trotz.south) -- (gen.north);
\draw[arrow] (waehrend.south) -- (gen.north);
\draw[arrow] (wegen.south) -- (gen.north);
\draw[arrow] (statt.south) -- (gen.north);
\draw[arrow] (innerhalb.north) -- (gen.south);
\draw[arrow] (ausserhalb.north) -- (gen.south);
\end{tikzpicture}
\legende{Les six prépositions principales régissant le génitif (niveau Terminale).}
\end{popfigure}

\begin{exoresolu}[Grammaire — Génitif et prépositions : Application]
\textbf{Énoncé} : Mettez le nom entre parenthèses au génitif correct (article + nom).\par
a) \all{Trotz \dots\ (die politische Unabhängigkeit) bleibt die Wirtschaft schwach.}\par
b) \all{Wegen \dots\ (der hohe Schuldenberg) kann das Land nicht investieren.}\par
c) \all{Während \dots\ (die Kolonialzeit) wurden Grenzen willkürlich gezogen.}\par
d) \all{Die Verarbeitung \dots\ (der Kakao) vor Ort schafft Arbeitsplätze.}\par
e) \all{Anstatt \dots\ (der Export) roher Bohnen sollte man Schokolade produzieren.}\par
f) \all{Innerhalb \dots\ (ein Jahrzehnt) soll die AfCFTA wirken.}
\tcblower
\textbf{Solution détaillée} :\par
a) \all{trotz der politischen Unabhängigkeit} (fém. sg : article \all{der}, adj. \all{-en}, nom sans \all{-s}).\par
b) \all{wegen des hohen Schuldenbergs} (masc. sg : article \all{des}, adj. \all{-en}, nom \all{-s} car monosyllabe \all{Berg} $\rightarrow$ \all{Bergs}).\par
c) \all{während der Kolonialzeit} (fém. sg : article \all{der}, nom inchangé).\par
d) \all{des Kakaos} (masc. sg : article \all{des}, nom \all{-s} / \all{-es} : \all{Kakao} $\rightarrow$ \all{Kakaos}).\par
e) \all{des Exports} (masc. sg : article \all{des}, nom \all{-s}).\par
f) \all{eines Jahrzehnts} (neutre sg indéfini : article \all{eines}, nom \all{-s} / \all{-es} : \all{Jahrzehnt} $\rightarrow$ \all{Jahrzehnts}).
\end{exoresolu}

\begin{experience}[Activité orale : « Debat : Rohstoffe verarbeiten oder exportieren? »]
\textbf{Consigne} : En binôme. L'un défend l'exportation brute (revenus rapides, pas d'usines), l'autre la transformation locale (emplois, valeur ajoutée, indépendance). Utilisez : \all{trotz, wegen, statt, die Ausbeutung, die Verarbeitung, die Abhängigkeit, die Souveränität}. Durée : 3 min. Puis changement de rôle.
\end{experience}

\courssub{Commentaire guidé et production écrite}

\begin{methode}[Rédiger un commentaire structuré / résumé argumenté (Type Bac)]
\begin{enumerate}
\item \textbf{Introduction} (1 phrase) : \all{Der Text „Neokolonialismus in Afrika“ handelt von ... / thematisiert ...} + thèse principale.
\item \textbf{Développement} (3-4 phrases) : Enchaîner les idées forces par des connecteurs logiques (\all{zwar ... aber}, \all{außerdem}, \all{deshalb}, \all{folglich}, \all{jedoch}). \textbf{Utiliser la nominalisation} pour condenser : \all{die Ausbeutung}, \all{die Verarbeitung}, \all{die Abhängigkeit}. Employer le \textbf{génitif} (\all{die Verarbeitung der Rohstoffe}) et les \textbf{prépositions + génitif} (\all{trotz der Unabhängigkeit}, \all{wegen der Schuldenfalle}).
\item \textbf{Conclusion / Prise de position personnelle} (1-2 phrases) : \all{Meiner Meinung nach ... / Ich bin der Ansicht, dass ... / Es ist entscheidend, dass ...} + \all{Konjunktiv II} ou \all{müssen/sollen} pour la recommandation.
\item \textbf{Relecture (Check-list)} : Verbe 2ème position ? Majuscules aux noms ? Génitifs corrects ? Prépositions + cas ? Umlaute / ß ? Ponctuation allemande ?
\end{enumerate}
\end{methode}

\begin{exoresolu}[Production écrite — Commentaire / Résumé argumenté (Modèle Bac)]
\textbf{Sujet} : \all{Fassen Sie den Text in ca. 80-100 Wörtern zusammen und äußern Sie sich zur Frage: „Was ist nötig für echte wirtschaftliche Unabhängigkeit Afrikas?“.}
\tcblower
\textbf{Production modèle} :\par
\all{Der Text „Neokolonialismus in Afrika“ thematisiert die anhaltende wirtschaftliche Abhängigkeit afrikanischer Staaten trotz ihrer politischen Unabhängigkeit seit den 1960er Jahren. Zwar sind die Länder souverän, aber die ökonomischen Strukturen der Kolonialzeit bestehen fort: Rohstoffe wie Kakao und Erdöl werden billig exportiert und teure Fertigprodukte importiert. Ein zentrales Problem ist die Schuldenfalle durch internationale Kredite, die die Entscheidungsfreiheit einschränken. Außerdem transferieren multinationale Konzerne Gewinne ab, statt lokale Wertschöpfungsketten aufzubauen.}\par
\all{Meiner Meinung nach ist die Verarbeitung der Rohstoffe im eigenen Land der Schlüssel zur echten Unabhängigkeit. Statt der bloßen Ausbeutung brauchen afrikanische Staaten eigene Industrien und einen starken innerafrikanischen Handel durch die AfCFTA. Nur wirtschaftliche Souveränität sichert langfristig politische Freiheit.}\par
\textbf{Analyse grammaticale du modèle} :\par
\begin{itemize}
\item \all{zwar ... aber} (concession, verbe en 2ème pos. dans les deux propositions).
\item \all{die ökonomischen Strukturen der Kolonialzeit} (Génitif \all{der Kolonialzeit}).
\item \all{werden ... exportiert / importiert} (Passif \all{werden} + Participe II, ordre correct).
\item \all{die Schuldenfalle durch internationale Kredite} (\all{durch} + Accusatif agent).
\item \all{statt der bloßen Ausbeutung} (\all{statt} + Génitif).
\item \all{Meiner Meinung nach ist ... der Schlüssel} (Verbe \all{ist} en 2ème position après adverbial).
\item \all{brauchen ... eigene Industrien} (Accusatif pluriel).
\item \all{Nur wirtschaftliche Souveränität sichert ...} (Inversion : Complément préposé en position 1 $\rightarrow$ Verbe \all{sichert} en position 2).
\end{itemize}
\end{exoresolu}

\begin{attention}[Les 5 erreurs qui coûtent des points au Bac (Synthèse)]
\begin{enumerate}
\item \textbf{Majuscules aux noms} : \all{die abhängigkeit, der export, die verarbeitung} $\times$ $\rightarrow$ \all{die Abhängigkeit, der Export, die Verarbeitung} $\checkmark$. \emph{Règle absolue : tout nom = majuscule.}
\item \textbf{Cas après prépositions (Génitif vs Datif)} : \all{trotz die Unabhängigkeit, wegen den Schulden, während die Zeit} $\times$ $\rightarrow$ \all{trotz \textbf{der} Unabhängigkeit, wegen \textbf{des} Schuldenbergs, während \textbf{der} Kolonialzeit} $\checkmark$. \emph{Apprendre la liste : trotz, während, wegen, statt, anstatt, innerhalb, außerhalb + GÉNITIF.}
\item \textbf{Ordre des mots (V2) après complément en tête} : \all{Trotz der Unabhängigkeit das Land bleibt arm} $\times$ $\rightarrow$ \all{Trotz der Unabhängigkeit \textbf{bleibt} das Land arm} $\checkmark$. \emph{Le verbe conjugué est \textbf{toujours} en 2ème position.}
\item \textbf{Nominalisation mal formée (suffixe / genre)} : \all{die Abhängung, der Ausbeutung, das Unabhängigkeit} $\times$ $\rightarrow$ \all{\textbf{die} Abhängigkeit (\textbf{-keit}), \textbf{die} Ausbeutung (\textbf{-ung}), \textbf{die} Unabhängigkeit (\textbf{-keit})} $\checkmark$. \emph{Vérifier : \all{-ung, -heit, -keit} = fém. (\all{die}) ; Infinitif = neutre (\all{das}) ; Conversion (\all{Export}) = masc. (\all{der}) souvent.}
\item \textbf{Faux amis / Constructions lexicales} : \all{abhängig auf} $\times$ $\rightarrow$ \all{abhängig \textbf{von}} + Datif $\checkmark$. \all{unabhängig von} $\checkmark$. \all{die Kolonie} (le territoire) vs \all{der Kolonialismus} (le système/idéologie). \all{die Unabhängigkeit} (l'indépendance) vs \all{unabhängig} (adj. indépendant).
\end{enumerate}
\end{attention}

\begin{aretenir}[Check-list finale pour le Bac — Leçon AL27]
\begin{itemize}
\item \textbf{Thème} : Neokolonialismus (indépendance politique vs dépendance économique).
\item \textbf{Lexique} : 10 noms clés avec article/pluriel + familles de mots (\all{abhängig $\leftrightarrow$ Abhängigkeit/Unabhängigkeit}, \all{ausbeuten $\leftrightarrow$ Ausbeutung}, \all{exportieren $\leftrightarrow$ Export/Exportieren}).
\item \textbf{Grammaire 1} : Nominalisation (\all{-ung, -heit/-keit, Infinitiv, Conversion}) + majuscule + génitif de l'objet.
\item \textbf{Grammaire 2} : Génitif (formation articles/noms) + 7 prépositions (\all{trotz, während, wegen, statt, anstatt, innerhalb, außerhalb}) + registre écrit.
\item \textbf{Compétences} : Compréhension (globale/détaillée), Résumé structuré (nominalisé, connecteurs), Prise de position (\all{Meiner Meinung nach} + V2).
\end{itemize}
\end{aretenir}

\newpage

\coursec{Exercices}

\courssub{Je vérifie mes connaissances}

\begin{exercice}[QCM : le texte et le thème]
Entoure la bonne réponse.
\begin{enumerate}
\item Der Neokolonialismus bedeutet :
\begin{enumerate}
\item die Rückkehr der Kolonialtruppen nach Afrika
\item die wirtschaftliche Abhängigkeit trotz politischer Unabhängigkeit
\item die Gründung neuer Kolonien in Asien
\end{enumerate}
\item „Die Ausbeutung der Rohstoffe" heißt auf Französisch :
\begin{enumerate}
\item l'exploration des matières premières
\item l'exploitation des matières premières
\item l'exportation des matières premières
\end{enumerate}
\item Welches Wort passt nicht zur Wortfamilie von „abhängig" ?
\begin{enumerate}
\item die Abhängigkeit
\item abhängen
\item der Hangar
\end{enumerate}
\item Die Präposition „wegen" verlangt :
\begin{enumerate}
\item den Akkusativ
\item den Dativ
\item den Genitiv
\end{enumerate}
\end{enumerate}
\end{exercice}

\begin{exercice}[Vrai ou faux ? Justifie ta réponse]
Sag, ob die Aussagen richtig (R) oder falsch (F) sind, und begründe deine Antwort mit einem Satz aus dem Text der Lektion.
\begin{enumerate}
\item Die afrikanischen Länder sind seit den 1960er Jahren völlig unabhängig.
\item Die ehemaligen Kolonialmächte kontrollieren oft noch die Preise der Rohstoffe.
\item Der Neokolonialismus ist nur ein militärisches Problem.
\item Viele junge Afrikaner kennen die Geschichte des Kolonialismus nicht.
\end{enumerate}
\end{exercice}

\begin{exercice}[Vocabulaire : articles et pluriels]
Ergänze die Tabelle mit dem bestimmten Artikel und dem Plural.
\begin{center}
\begin{tabular}{|l|l|l|}
\hline
\rowcolor{popblueL} Substantiv & Artikel & Plural \\
\hline
Kolonie & \ldots & \ldots \\
\hline
Kolonialismus & \ldots & \ldots \\
\hline
Abhängigkeit & \ldots & \ldots \\
\hline
Ausbeutung & \ldots & \ldots \\
\hline
Unabhängigkeit & \ldots & \ldots \\
\hline
Rohstoff & \ldots & \ldots \\
\hline
Konzern & \ldots & \ldots \\
\hline
Macht & \ldots & \ldots \\
\hline
\end{tabular}
\end{center}
\end{exercice}

\begin{exercice}[Conjugaison : les verbes du thème]
Ergänze die Sätze mit dem Verb im Präsens.
\begin{enumerate}
\item Die Konzerne \ldots\ (ausbeuten) die Rohstoffe des Landes.
\item Viele Länder \ldots\ (bleiben) von den ehemaligen Kolonialmächten abhängig.
\item Der Kontinent \ldots\ (besitzen) große Reichtümer.
\item Die Regierungen \ldots\ (verhandeln) über die Preise des Kakaos.
\item Wir \ldots\ (kämpfen) für eine echte Unabhängigkeit.
\end{enumerate}
\end{exercice}

\courssub{Je m'entraîne}

\begin{exercice}[Nominalisation : du verbe au nom]
Verwandle die verbalen Konstruktionen in nominale Konstruktionen (Nominalisierung). Beispiel : \all{Die Konzerne beuten die Rohstoffe aus.} $\rightarrow$ \all{die Ausbeutung der Rohstoffe durch die Konzerne}
\begin{enumerate}
\item Die Länder hängen von den Konzernen ab.
\item Die Europäer kolonisierten Afrika im 19. Jahrhundert.
\item Die afrikanischen Staaten exportieren Kakao und Kaffee.
\item Die Regierung kontrolliert die Preise nicht.
\item Die ehemaligen Mächte beeinflussen die Politik der Länder.
\item Die Bauern produzieren viel Kakao in der Region von Yaoundé.
\end{enumerate}
\end{exercice}

\begin{exercice}[Nominalisation : du nom au verbe]
Verwandle die nominalen Konstruktionen in vollständige Sätze mit einem Verb.
\begin{enumerate}
\item die Unabhängigkeit des Landes im Jahr 1960
\item die Ausbeutung der Arbeiter durch die Konzerne
\item der Export der Rohstoffe nach Europa
\item die Abhängigkeit der Bauern von den Weltmarktpreisen
\end{enumerate}
\end{exercice}

\begin{exercice}[Le génitif : phrases à trous]
Ergänze die Sätze mit der richtigen Genitivform des Wortes in Klammern.
\begin{enumerate}
\item Trotz \ldots\ (die Unabhängigkeit) bleibt das Land von \ldots\ (die Schulden) abhängig.
\item Wegen \ldots\ (der Kolonialismus) leiden viele Länder noch heute.
\item Während \ldots\ (die Konferenz) diskutierten die Staatschefs über \ldots\ (der Handel).
\item Die Preise \ldots\ (die Rohstoffe) steigen und fallen schnell.
\item Innerhalb \ldots\ (das Jahrzehnt) wurden viele Staaten unabhängig.
\item Die Wirtschaft \ldots\ (der Kontinent) wächst langsam.
\end{enumerate}
\end{exercice}

\begin{exercice}[Appariements : mots et définitions]
Verbinde jedes Wort (1--6) mit der passenden Definition (a--f).
\begin{center}
\begin{tabular}{|l|l|}
\hline
\rowcolor{popblueL} Wort & Definition \\
\hline
1. die Kolonie & a. ein Land ist von einem anderen Land wirtschaftlich abhängig \\
\hline
2. der Rohstoff & b. ein Gebiet, das von einer fremden Macht beherrscht wird \\
\hline
3. die Ausbeutung & c. natürlicher Reichtum, z. B. Kakao, Öl, Gold \\
\hline
4. die Unabhängigkeit & d. jemanden für den eigenen Profit ausnutzen \\
\hline
5. der Neokolonialismus & e. der Zustand, frei und selbstständig zu sein \\
\hline
6. der Konzern & f. ein sehr großes internationales Unternehmen \\
\hline
\end{tabular}
\end{center}
\end{exercice}

\courssub{Je me prépare au Bac}

\begin{exercice}[Texte et compréhension type Bac]
Lies den folgenden Text und beantworte die Fragen mit vollständigen Sätzen auf Deutsch.

\begin{textebox}[Der Kakao zwischen Kamerun und Deutschland]
Kamerun ist einer der wichtigsten Produzenten von Kakao in Afrika. Tausende Bauern in der Region um Yaoundé und im Südwesten des Landes arbeiten jeden Tag auf ihren Plantagen. Doch trotz dieser harten Arbeit bleiben viele von ihnen arm. Der Grund ist einfach : die Preise des Kakaos werden nicht in Kamerun, sondern an den Börsen in Europa bestimmt. Ein großer Teil des kamerunischen Kakaos wird roh nach Europa exportiert, unter anderem nach Deutschland. Dort wird er in Fabriken zu Schokolade verarbeitet. Die Verarbeitung des Kakaos bringt viel mehr Geld als der Verkauf der rohen Bohnen. Deshalb fordern viele Experten die Industrialisierung Afrikas : die Länder sollen ihre Rohstoffe selbst verarbeiten. Erste Schokoladenfabriken in Kamerun zeigen, dass eine echte wirtschaftliche Unabhängigkeit möglich ist.
\end{textebox}

\textbf{Fragen :}
\begin{enumerate}
\item Wo arbeiten die kamerunischen Kakaobauern ?
\item Warum bleiben viele Bauern trotz ihrer harten Arbeit arm ?
\item Was geschieht mit dem Kakao in Deutschland ?
\item Was fordern die Experten, und warum ?
\item Finde im Text zwei Nominalisierungen und verwandle sie in Sätze mit einem Verb.
\item Erkläre mit eigenen Worten : Warum ist die Verarbeitung der Rohstoffe im eigenen Land wichtig für die Unabhängigkeit ?
\end{enumerate}
\end{exercice}

\begin{exercice}[Thème et version]
\textbf{A. Thème.} Übersetze die folgenden Sätze ins Deutsche.
\begin{enumerate}
\item Malgré l'indépendance, beaucoup de pays africains restent dépendants des anciennes puissances coloniales.
\item L'exploitation des matières premières par les grandes entreprises est un problème sérieux.
\item À cause des dettes, le gouvernement ne peut pas investir dans les écoles.
\item Les prix du cacao et du café sont fixés en Europe.
\item Pendant la colonisation, les richesses du continent ont été exportées.
\end{enumerate}

\textbf{B. Version.} Traduis le paragraphe suivant en français.

\begin{textebox}[Afrikas Reichtümer]
Die Ausbeutung der Rohstoffe Afrikas begann während der Kolonialzeit. Trotz der Unabhängigkeit vieler Staaten in den 1960er Jahren blieb die Abhängigkeit des Kontinents von den Weltmärkten bestehen. Wegen der niedrigen Preise für Rohstoffe verlieren die Länder jedes Jahr viel Geld. Die Forderung nach einer gerechten Verteilung der Reichtümer wird deshalb immer lauter.
\end{textebox}
\end{exercice}

\begin{exercice}[Production écrite type Bac]
\textbf{Sujet :} \all{Die Unabhängigkeit Afrikas : eine Illusion ?}

Rédige en allemand un texte argumenté de \textbf{80 à 100 mots}. Tu donneras ton avis sur la question en réutilisant :
\begin{itemize}
\item au moins quatre mots du lexique de la leçon (\all{die Kolonie, der Kolonialismus, die Abhängigkeit, die Ausbeutung, die Unabhängigkeit, der Rohstoff}) ;
\item au moins deux nominalisations avec le génitif ;
\item au moins une préposition avec le génitif (\all{trotz, wegen, während}).
\end{itemize}

Structure ton texte en trois parties : introduction (le problème), développement (arguments et exemples, par exemple le cacao du Cameroun), conclusion (ton avis personnel).
\end{exercice}

\newpage

\coursec{Corrigés des exercices}

\coursec{Corrigés détaillés de la leçon AL27 : Le néocolonialisme}

\courssub{Exercice 1 — QCM : le texte et le thème}
\begin{corrige}[Exercice 1 — QCM : le texte et le thème]
\begin{enumerate}
\item Réponse \textbf{b} : \all{die wirtschaftliche Abhängigkeit trotz politischer Unabhängigkeit}. C'est la définition même du néocolonialisme : les anciennes colonies sont politiquement indépendantes, mais restent économiquement dépendantes.
\item Réponse \textbf{b} : « l'exploitation des matières premières ». Attention au faux ami : \all{die Exploration} signifie « l'exploration » (recherche), \all{die Ausbeutung} signifie « l'exploitation » au sens d'utilisation abusive.
\item Réponse \textbf{c} : \all{der Hangar} (le hangar) n'appartient pas à la famille de \all{abhängig}. La famille comprend \all{die Abhängigkeit} (la dépendance) et \all{abhängen von} (dépendre de).
\item Réponse \textbf{c} : \all{wegen} exige le génitif : \all{wegen des Kolonialismus} (à cause du colonialisme).
\end{enumerate}
\textbf{Points de vigilance :} ne pas confondre \all{die Ausbeutung} et \all{die Exportation} ; retenir que \all{wegen, trotz, während, innerhalb} sont suivies du \cle{génitif}.
\end{corrige}

\courssub{Exercice 2 — Vrai ou faux ? Justifie ta réponse}
\begin{corrige}[Exercice 2 — Vrai ou faux ? Justifie ta réponse]
\begin{enumerate}
\item \textbf{F} (faux). Justification : \all{Die afrikanischen Länder sind zwar politisch unabhängig, aber wirtschaftlich bleiben viele von den ehemaligen Kolonialmächten abhängig.} (Les pays africains sont certes politiquement indépendants, mais beaucoup restent économiquement dépendants des anciennes puissances coloniales.)
\item \textbf{R} (vrai). Justification : \all{Die Preise der Rohstoffe werden oft noch an den Börsen in Europa bestimmt.} (Les prix des matières premières sont souvent encore fixés aux bourses d'Europe.)
\item \textbf{F} (faux). Justification : \all{Der Neokolonialismus ist vor allem ein wirtschaftliches und politisches Problem, kein militärisches.} (Le néocolonialisme est surtout un problème économique et politique, non militaire.)
\item \textbf{F} (faux). Justification : \all{Viele junge Afrikaner lernen in der Schule die Geschichte des Kolonialismus und engagieren sich für ihre Zukunft.} (Beaucoup de jeunes Africains apprennent à l'école l'histoire du colonialisme et s'engagent pour leur avenir.)
\end{enumerate}
\textbf{Points de vigilance :} la justification doit reprendre une idée du texte sous forme de phrase complète (sujet + verbe conjugué en \cle{deuxième position}).
\end{corrige}

\courssub{Exercice 3 — Vocabulaire : articles et pluriels}
\begin{corrige}[Exercice 3 — Vocabulaire : articles et pluriels]
\begin{center}
\begin{tabular}{|l|l|l|}
\hline
\rowcolor{popblueL} Substantiv & Artikel & Plural \\
\hline
Kolonie & die & die Kolonien \\
\hline
Kolonialismus & der & --- (pas de pluriel) \\
\hline
Abhängigkeit & die & die Abhängigkeiten \\
\hline
Ausbeutung & die & die Ausbeutungen \\
\hline
Unabhängigkeit & die & --- (pas de pluriel) \\
\hline
Rohstoff & der & die Rohstoffe \\
\hline
Konzern & der & die Konzerne \\
\hline
Macht & die & die Mächte \\
\hline
\end{tabular}
\end{center}
\textbf{Justifications :} les noms en \all{-ung}, \all{-heit}, \all{-keit}, \all{-ie} sont féminins (\all{die}) ; \all{-ismus} est masculin (\all{der}). \all{die Macht} forme son pluriel avec \cle{Umlaut} : \all{die Mächte}. \all{der Kolonialismus} et \all{die Unabhängigkeit} sont des noms abstraits, en général sans pluriel.
\end{corrige}

\courssub{Exercice 4 — Conjugaison : les verbes du thème}
\begin{corrige}[Exercice 4 — Conjugaison : les verbes du thème]
\begin{enumerate}
\item \all{Die Konzerne \textbf{beuten} die Rohstoffe des Landes \textbf{aus}.} — Verbe à particule séparable : \all{ausbeuten} ; la particule \all{aus} va en fin de phrase.
\item \all{Viele Länder \textbf{bleiben} von den ehemaligen Kolonialmächten abhängig.} — \all{bleiben} est un verbe fort, mais régulier au présent.
\item \all{Der Kontinent \textbf{besitzt} große Reichtümer.} — 3\up{e} personne du singulier : \all{besitzen} $\rightarrow$ \all{besitzt}.
\item \all{Die Regierungen \textbf{verhandeln} über die Preise des Kakaos.} — Verbe faible régulier.
\item \all{Wir \textbf{kämpfen} für eine echte Unabhängigkeit.} — \all{kämpfen für} + accusatif : « lutter pour ».
\end{enumerate}
\textbf{Points de vigilance :} avec les \cle{verbes à particule séparable} (\all{ausbeuten}), le verbe conjugué est en deuxième position et la particule en fin de phrase — erreur classique des francophones.
\end{corrige}

\courssub{Exercice 5 — Nominalisation : du verbe au nom}
\begin{corrige}[Exercice 5 — Nominalisation : du verbe au nom]
\begin{enumerate}
\item \all{Die Länder hängen von den Konzernen ab.} $\rightarrow$ \all{die Abhängigkeit der Länder von den Konzernen} (la dépendance des pays vis-à-vis des groupes industriels).
\item \all{Die Europäer kolonisierten Afrika im 19. Jahrhundert.} $\rightarrow$ \all{die Kolonisierung Afrikas durch die Europäer im 19. Jahrhundert} (la colonisation de l'Afrique par les Européens au XIX\up{e} siècle).
\item \all{Die afrikanischen Staaten exportieren Kakao und Kaffee.} $\rightarrow$ \all{der Export von Kakao und Kaffee durch die afrikanischen Staaten} (l'exportation du cacao et du café par les États africains).
\item \all{Die Regierung kontrolliert die Preise nicht.} $\rightarrow$ \all{die fehlende Kontrolle der Preise durch die Regierung} (l'absence de contrôle des prix par le gouvernement).
\item \all{Die ehemaligen Mächte beeinflussen die Politik der Länder.} $\rightarrow$ \all{die Beeinflussung der Politik der Länder durch die ehemaligen Mächte} (l'influence exercée sur la politique des pays par les anciennes puissances).
\item \all{Die Bauern produzieren viel Kakao in der Region von Yaoundé.} $\rightarrow$ \all{die Produktion von viel Kakao durch die Bauern in der Region von Yaoundé} (la production de beaucoup de cacao par les paysans dans la région de Yaoundé).
\end{enumerate}
\textbf{Règle appliquée :} le sujet de la phrase verbale devient un complément introduit par \all{durch} (+ \cle{accusatif} — construction standard pour l'agent) ou plus rarement par le \cle{génitif} ; le complément d'objet direct devient un \cle{génitif} (\all{Afrikas}, \all{der Preise}).
\end{corrige}

\courssub{Exercice 6 — Nominalisation : du nom au verbe}
\begin{corrige}[Exercice 6 — Nominalisation : du nom au verbe]
\begin{enumerate}
\item \all{die Unabhängigkeit des Landes im Jahr 1960} $\rightarrow$ \all{Das Land wurde im Jahr 1960 unabhängig.} (Le pays est devenu indépendant en 1960.)
\item \all{die Ausbeutung der Arbeiter durch die Konzerne} $\rightarrow$ \all{Die Konzerne beuten die Arbeiter aus.} (Les groupes industriels exploitent les travailleurs.)
\item \all{der Export der Rohstoffe nach Europa} $\rightarrow$ \all{Die Länder exportieren die Rohstoffe nach Europa.} (Les pays exportent les matières premières vers l'Europe.)
\item \all{die Abhängigkeit der Bauern von den Weltmarktpreisen} $\rightarrow$ \all{Die Bauern hängen von den Weltmarktpreisen ab.} (Les paysans dépendent des prix du marché mondial.)
\end{enumerate}
\textbf{Points de vigilance :} dans la phrase verbale, le verbe conjugué revient en \cle{deuxième position} ; \all{abhängen} est séparable : \all{hängen \ldots\ ab}.
\end{corrige}

\courssub{Exercice 7 — Le génitif : phrases à trous}
\begin{corrige}[Exercice 7 — Le génitif : phrases à trous]
\begin{enumerate}
\item \all{Trotz \textbf{der Unabhängigkeit} bleibt das Land von \textbf{den Schulden} abhängig.} — \all{trotz} + génitif ; \all{von} + datif pluriel (\all{den Schulden}).
\item \all{Wegen \textbf{des Kolonialismus} leiden viele Länder noch heute.} — génitif masculin : \all{des Kolonialismus} (sans \all{-s} supplémentaire, car le mot se termine déjà par \all{-s}).
\item \all{Während \textbf{der Konferenz} diskutierten die Staatschefs über \textbf{den Handel}.} — \all{während} + génitif ; \all{über} + accusatif (\all{den Handel}).
\item \all{Die Preise \textbf{der Rohstoffe} steigen und fallen schnell.} — génitif pluriel : \all{der Rohstoffe}.
\item \all{Innerhalb \textbf{eines Jahrzehnts} wurden viele Staaten unabhängig.} — génitif neutre : \all{eines Jahrzehnts}.
\item \all{Die Wirtschaft \textbf{des Kontinents} wächst langsam.} — génitif masculin : \all{des Kontinents}.
\end{enumerate}
\textbf{Rappel :} au génitif, les noms masculins et neutres prennent \all{-(e)s} : \all{des Landes, des Jahres, des Kontinents} ; les noms féminins et les pluriels restent inchangés : \all{der Unabhängigkeit, der Rohstoffe}.
\end{corrige}

\courssub{Exercice 8 — Appariements : mots et définitions}
\begin{corrige}[Exercice 8 — Appariements : mots et définitions]
\begin{center}
\begin{tabular}{|l|l|}
\hline
\rowcolor{popblueL} Wort & Definition \\
\hline
1. die Kolonie & b. ein Gebiet, das von einer fremden Macht beherrscht wird \\
\hline
2. der Rohstoff & c. natürlicher Reichtum, z. B. Kakao, Öl, Gold \\
\hline
3. die Ausbeutung & d. jemanden für den eigenen Profit ausnutzen \\
\hline
4. die Unabhängigkeit & e. der Zustand, frei und selbstständig zu sein \\
\hline
5. der Neokolonialismus & a. ein Land ist von einem anderen Land wirtschaftlich abhängig \\
\hline
6. der Konzern & f. ein sehr großes internationales Unternehmen \\
\hline
\end{tabular}
\end{center}
\textbf{Vérification :} 1-b, 2-c, 3-d, 4-e, 5-a, 6-f. Traductions : la colonie = un territoire dominé par une puissance étrangère ; la matière première = une richesse naturelle ; l'exploitation = utiliser quelqu'un pour son propre profit ; l'indépendance = l'état d'être libre ; le néocolonialisme = la dépendance économique d'un pays envers un autre ; le groupe industriel = une très grande entreprise internationale.
\end{corrige}

\courssub{Exercice 9 — Texte et compréhension type Bac}
\begin{corrige}[Exercice 9 — Texte et compréhension type Bac]
\textbf{Réponses modèles :}
\begin{enumerate}
\item \all{Die kamerunischen Kakaobauern arbeiten in der Region um Yaoundé und im Südwesten des Landes auf ihren Plantagen.} (Les paysans camerounais travaillent dans la région de Yaoundé et dans le sud-ouest du pays, sur leurs plantations.)
\item \all{Viele Bauern bleiben arm, weil die Preise des Kakaos nicht in Kamerun, sondern an den Börsen in Europa bestimmt werden.} (Beaucoup de paysans restent pauvres parce que les prix du cacao ne sont pas fixés au Cameroun, mais aux bourses d'Europe.)
\item \all{In Deutschland wird der Kakao in Fabriken zu Schokolade verarbeitet.} (En Allemagne, le cacao est transformé en chocolat dans des usines.)
\item \all{Die Experten fordern die Industrialisierung Afrikas : die Länder sollen ihre Rohstoffe selbst verarbeiten, weil die Verarbeitung viel mehr Geld bringt als der Verkauf der rohen Bohnen.} (Les experts demandent l'industrialisation de l'Afrique : les pays doivent transformer eux-mêmes leurs matières premières, car la transformation rapporte beaucoup plus d'argent que la vente des fèves brutes.)
\item Deux nominalisations : \all{die Verarbeitung des Kakaos} $\rightarrow$ \all{Man verarbeitet den Kakao.} ; \all{die Industrialisierung Afrikas} $\rightarrow$ \all{Man industrialisiert Afrika.} (ou \all{Afrika wird industrialisiert.})
\item \all{Die Verarbeitung der Rohstoffe im eigenen Land ist wichtig, weil sie mehr Geld und Arbeitsplätze bringt. So wird das Land wirtschaftlich unabhängiger und muss seine Reichtümer nicht roh exportieren.} (La transformation des matières premières dans le pays est importante parce qu'elle rapporte plus d'argent et d'emplois. Ainsi le pays devient économiquement plus indépendant.)
\end{enumerate}
\textbf{Barème indicatif (sur 12 points) :} 2 points par question (1 point pour le contenu exact, 1 point pour la correction grammaticale : place du verbe, cas, orthographe des noms avec majuscule).
\textbf{Points de vigilance :} répondre par des phrases complètes ; attention au \cle{passif} (\all{werden bestimmt}, \all{wird verarbeitet}) ; ne pas oublier la majuscule aux noms (\all{Kakao, Börsen, Fabriken, Arbeitsplätze}).
\end{corrige}

\courssub{Exercice 10 — Thème et version}
\begin{corrige}[Exercice 10 — Thème et version]
\textbf{A. Thème — traductions modèles :}
\begin{enumerate}
\item \all{Trotz der Unabhängigkeit bleiben viele afrikanische Länder von den ehemaligen Kolonialmächten abhängig.} — \all{trotz} + génitif ; \all{abhängig bleiben von} + datif.
\item \all{Die Ausbeutung der Rohstoffe durch die großen Unternehmen ist ein ernstes Problem.} — nominalisation avec double génitif et \all{durch} + accusatif.
\item \all{Wegen der Schulden kann die Regierung nicht in die Schulen investieren.} — \all{wegen} + génitif pluriel ; verbe modal : l'infinitif \all{investieren} en fin de phrase.
\item \all{Die Preise des Kakaos und des Kaffees werden in Europa festgelegt.} — passif au présent : \all{werden} + participe II ; \all{festlegen} est séparable (\all{festgelegt}).
\item \all{Während der Kolonialzeit wurden die Reichtümer des Kontinents exportiert.} — passif au prétérit : \all{wurden} + participe II ; \all{während} + génitif.
\end{enumerate}
\textbf{B. Version — traduction modèle :}

„ L'exploitation des matières premières de l'Afrique a commencé pendant la période coloniale. Malgré l'indépendance de nombreux États dans les années 1960, la dépendance du continent vis-à-vis des marchés mondiaux a persisté. À cause des prix bas des matières premières, les pays perdent chaque année beaucoup d'argent. La revendication d'une répartition équitable des richesses se fait donc de plus en plus forte. “

\textbf{Barème indicatif (sur 10 points) :} 1 point par phrase du thème (5 points) ; 5 points pour la version (fidélité du sens : 3 points ; qualité du français : 2 points).
\textbf{Points de vigilance :} \all{die Forderung nach} = « la revendication de » ; \all{bestehen bleiben} = « persister » ; ne pas traduire \all{gerecht} par « juste » au sens de « exact », mais par « équitable ».
\end{corrige}

\courssub{Exercice 11 — Production écrite type Bac}
\begin{corrige}[Exercice 11 — Production écrite type Bac]
\textbf{Proposition de rédaction modèle (environ 95 mots) :}

\begin{textebox}[Die Unabhängigkeit Afrikas : eine Illusion?]
Seit den 1960er Jahren sind die afrikanischen Staaten politisch unabhängig. Aber ist diese Unabhängigkeit wirklich vollständig?

Trotz der Unabhängigkeit bleibt die Abhängigkeit vieler Länder von den ehemaligen Kolonialmächten groß. Die Ausbeutung der Rohstoffe durch internationale Konzerne ist ein Beispiel dafür. Kamerun exportiert seinen Kakao roh nach Europa, und die Preise werden an den Börsen in Europa bestimmt. Wegen dieser Abhängigkeit verlieren die Länder viel Geld.

Meiner Meinung nach ist die Unabhängigkeit keine Illusion, aber sie ist noch unvollständig. Die Industrialisierung Afrikas, zum Beispiel der Bau von Schokoladenfabriken in Kamerun, ist der Weg zu einer echten wirtschaftlichen Freiheit.
\end{textebox}

\textbf{Traduction du modèle :} « Depuis les années 1960, les États africains sont politiquement indépendants. Mais cette indépendance est-elle vraiment complète ? Malgré l'indépendance, la dépendance de nombreux pays envers les anciennes puissances coloniales reste grande. L'exploitation des matières premières par des groupes internationaux en est un exemple. Le Cameroun exporte son cacao brut vers l'Europe, et les prix sont fixés aux bourses européennes. À cause de cette dépendance, les pays perdent beaucoup d'argent. À mon avis, l'indépendance n'est pas une illusion, mais elle est encore incomplète. L'industrialisation de l'Afrique, par exemple la construction d'usines de chocolat au Cameroun, est le chemin vers une véritable liberté économique. »

\textbf{Vérification des contraintes :}
\begin{itemize}
\item Lexique de la leçon : \all{die Unabhängigkeit, die Abhängigkeit, die Ausbeutung, der Rohstoff, der Kolonialismus} (au moins quatre mots) ;
\item Nominalisations avec génitif : \all{die Abhängigkeit vieler Länder}, \all{die Ausbeutung der Rohstoffe}, \all{die Industrialisierung Afrikas} ;
\item Prépositions avec le génitif : \all{trotz der Unabhängigkeit}, \all{wegen dieser Abhängigkeit}.
\end{itemize}

\textbf{Barème indicatif (sur 20 points) :} respect des consignes et de la structure en trois parties (4 points) ; richesse et exactitude du lexique thématique (4 points) ; grammaire : nominalisations, génitif, place du verbe (6 points) ; cohérence et argumentation (4 points) ; orthographe, majuscules, Umlaute (2 points).

\textbf{Points de vigilance :}
\begin{itemize}
\item Ne pas écrire \all{trotz die Unabhängigkeit} : \all{trotz} exige le génitif (\all{trotz der Unabhängigkeit}).
\item Majuscule à tous les noms : \all{Illusion, Freiheit, Beispiel, Geld, Reichtümer}.
\item Le verbe conjugué reste en deuxième position, même après \all{trotz der Unabhängigkeit} ou \all{meiner Meinung nach}.
\item Ne pas calquer le français « être dépendant de » par \all{sein abhängig von} : la forme correcte est \all{von \ldots\ abhängig sein} ou \all{von \ldots\ abhängen}.
\end{itemize}
\end{corrige}

\newpage

\coursec{Fiche bilan}

\begin{popfigure}
\begin{tikzpicture}[mindmap, grow cyclic, every node/.style={concept, text=white, font=\small},
  root concept/.append style={concept color=popdark, font=\small\bfseries, minimum size=2.6cm},
  level 1/.append style={level distance=3.4cm, sibling angle=51.4},
  level 1 concept/.append style={minimum size=2.1cm, font=\scriptsize\bfseries},
  level 2/.append style={level distance=2.2cm, sibling angle=45},
  level 2 concept/.append style={minimum size=1.7cm, font=\scriptsize}]
\node[root concept, align=center]{Der Neokolonialismus\\ (AL27)}
  child[concept color=popblue]{ node{Texte : compréhension}
    child[concept color=popblue!70]{ node{global puis détaillé} }
    child[concept color=popblue!70]{ node{questions + résumé} }
  }
  child[concept color=poporange]{ node{Lexique : Kolonialismus}
    child[concept color=poporange!70]{ node{die Kolonie, -n} }
    child[concept color=poporange!70]{ node{die Ausbeutung} }
  }
  child[concept color=popgreen]{ node{Nominalisierung}
    child[concept color=popgreen!70]{ node{Verbe $\rightarrow$ Nomen} }
    child[concept color=popgreen!70]{ node{die Unabhängigkeit} }
  }
  child[concept color=poppurple]{ node{Der Genitiv}
    child[concept color=poppurple!70]{ node{des, der, eines} }
    child[concept color=poppurple!70]{ node{die Ausbeutung der Völker} }
  }
  child[concept color=poppink]{ node{Präpositionen + Genitiv}
    child[concept color=poppink!70]{ node{wegen, trotz, während} }
    child[concept color=poppink!70]{ node{anstatt, innerhalb} }
  }
  child[concept color=popgold]{ node{Expression}
    child[concept color=popgold!70]{ node{court paragraphe} }
    child[concept color=popgold!70]{ node{avis argumenté} }
  }
  child[concept color=popmuted]{ node{Méthode}
    child[concept color=popmuted!70]{ node{lire, repérer, reformuler} }
  };
\end{tikzpicture}
\legende{Carte mentale de la leçon : le néocolonialisme}
\end{popfigure}

\begin{aretenir}[L'essentiel de la leçon AL27]
\textbf{1. Lexique clé (à connaître avec article et pluriel)}
\begin{itemize}
  \item \all{die Kolonie, -n} : la colonie ; \all{der Kolonialismus} (sans pluriel) : le colonialisme ; \all{der Neokolonialismus} (sans pluriel) : le néocolonialisme.
  \item \all{die Abhängigkeit, -en} : la dépendance ; \all{die Unabhängigkeit, -en} : l'indépendance ; \all{abhängig (von + Dat.)} : dépendant (de) ; \all{unabhängig} : indépendant.
  \item \all{die Ausbeutung, -en} : l'exploitation (au sens d'abus) ; \all{ausbeuten} : exploiter (qqn) ; \all{der Ausbeuter, -} : l'exploiteur.
  \item \all{der Rohstoff, -e} : la matière première ; \all{die Ressource, -n} : la ressource ; \all{der Reichtum, die Reichtümer} : la richesse.
  \item \all{das Volk, die Völker} : le peuple ; \all{der Handel} (sans pluriel) : le commerce ; \all{die Macht, die Mächte} : le pouvoir / les puissances.
  \item Famille de mots : \all{hängen} $\rightarrow$ \all{abhängen (von + Dat.)} $\rightarrow$ \all{abhängig} $\rightarrow$ \all{die Abhängigkeit}.
\end{itemize}

\textbf{2. Grammaire à connaître par cœur}
\begin{center}
\begin{tabularx}{\linewidth}{|X|X|X|}
\hline
\rowcolor{popblueL} Point & Règle & Exemple \\
\hline
Nominalisation & Verbe ou adjectif $\rightarrow$ nom (souvent en \all{-ung}, \all{-heit}, \all{-keit}, \all{-en}) ; majuscule obligatoire ; genre presque toujours féminin en \all{-ung/-heit/-keit} & \all{ausbeuten $\rightarrow$ die Ausbeutung} ; \all{unabhängig $\rightarrow$ die Unabhängigkeit} \\
\hline
Génitif & \all{des / eines} (m., n.) ; \all{der / einer} (f., pl.) ; \all{-(e)s} au nom m. et n. & \all{die Ausbeutung der Rohstoffe} ; \all{die Unabhängigkeit des Volkes} \\
\hline
Prépositions + génitif & \all{wegen, trotz, während, anstatt, innerhalb, außerhalb, aufgrund} & \all{trotz der Unabhängigkeit} ; \all{wegen der Abhängigkeit} \\
\hline
\end{tabularx}
\end{center}

\textbf{3. Méthodes express}
\begin{itemize}
  \item \textbf{Comprendre le texte :} 1) lire le titre et anticiper ; 2) lecture globale (qui ? quoi ? où ?) ; 3) lecture détaillée avec le lexique ; 4) répondre par des phrases complètes en allemand.
  \item \textbf{Résumer :} repérer les idées principales, remplacer les verbes par des noms (\all{Die Länder werden ausgebeutet} $\rightarrow$ \all{die Ausbeutung der Länder}), rédiger 3 à 5 phrases.
  \item \textbf{Produire :} donner son avis avec \all{Meiner Meinung nach...}, \all{Ich bin der Ansicht, dass...} + verbe en fin de proposition subordonnée.
\end{itemize}
\end{aretenir}

\begin{attention}[Pièges fréquents des francophones]
\begin{itemize}
  \item \textbf{Majuscule :} tout nom issu d'une nominalisation prend une majuscule : on écrit \all{die Ausbeutung}, jamais \emph{die ausbeutung}.
  \item \textbf{Genre :} les noms en \all{-ung}, \all{-heit}, \all{-keit} sont féminins : \all{die Abhängigkeit}, \all{die Unabhängigkeit}, \all{die Ausbeutung}.
  \item \textbf{Génitif :} n'oublie pas le \all{-(e)s} au masculin et au neutre : \all{des Volkes}, \all{des Landes}, \all{des Kontinents} ; mais jamais de \all{-s} au féminin : \all{der Kolonie}, \all{der Macht}.
  \item \textbf{Faux ami :} \all{die Ressource} existe, mais pour « matière première » on dit d'abord \all{der Rohstoff} ; \all{die Ausbeutung} signifie « l'exploitation abusive », pas « l'exploitation agricole » (\all{der Landbau}, \all{die Landwirtschaft}).
  \item \textbf{Registre oral :} à l'oral, on entend souvent \all{wegen} et \all{trotz} avec le datif (\all{wegen dem Wetter}) ; à l'écrit et à l'examen, exige le génitif : \all{wegen des Wetters}.
  \item \textbf{Place du verbe :} après \all{dass}, \all{weil}, \all{obwohl}, le verbe conjugué va en dernière position : \all{Ich bin der Ansicht, dass der Neokolonialismus eine neue Form der Abhängigkeit ist.}
\end{itemize}
\end{attention}

\begin{exemplebox}[Modèles de phrases pour la production]
\begin{itemize}
  \item \all{Der Neokolonialismus ist eine neue Form der wirtschaftlichen Abhängigkeit der afrikanischen Länder.} « Le néocolonialisme est une nouvelle forme de dépendance économique des pays africains. »
  \item \all{Trotz der Unabhängigkeit bleiben viele Länder von den ehemaligen Kolonialmächten abhängig.} « Malgré l'indépendance, beaucoup de pays restent dépendants des anciennes puissances coloniales. »
  \item \all{Während der Kolonialzeit wurden die Rohstoffe Afrikas ausgebeutet.} « Pendant la période coloniale, les matières premières de l'Afrique ont été exploitées. »
  \item \all{Meiner Meinung nach müssen die afrikanischen Staaten ihre Rohstoffe selbst verarbeiten, anstatt sie billig zu exportieren.} « À mon avis, les États africains doivent transformer eux-mêmes leurs matières premières au lieu de les exporter à bas prix. »
\end{itemize}
\end{exemplebox}

\begin{aretenir}[Checklist avant le Bac]
\begin{itemize}
  \item $\square$ Je connais l'article et le pluriel de : \all{die Kolonie, die Abhängigkeit, die Ausbeutung, der Rohstoff, das Volk}.
  \item $\square$ Je sais définir \all{der Neokolonialismus} en une phrase allemande simple.
  \item $\square$ Je sais transformer un verbe en nom : \all{ausbeuten $\rightarrow$ die Ausbeutung}, \all{abhängen $\rightarrow$ die Abhängigkeit}.
  \item $\square$ Je mets une majuscule à tous les noms issus d'une nominalisation.
  \item $\square$ Je connais les articles au génitif : \all{des, der, des, der} / \all{eines, einer, eines, --}.
  \item $\square$ J'ajoute \all{-(e)s} au nom masculin ou neutre au génitif : \all{des Volkes, des Landes}.
  \item $\square$ Je place correctement le complément du nom : \all{die Ausbeutung der Rohstoffe Afrikas}.
  \item $\square$ Je maîtrise les prépositions + génitif : \all{wegen, trotz, während, anstatt, aufgrund}.
  \item $\square$ Je sais répondre aux questions de compréhension par des phrases complètes.
  \item $\square$ Je sais résumer un texte en 3 à 5 phrases avec des nominalisations.
  \item $\square$ Je sais exprimer mon avis : \all{Meiner Meinung nach ist der Neokolonialismus eine neue Form der Abhängigkeit.}
  \item $\square$ Je vérifie la place du verbe : deuxième position en principale, dernière position après \all{dass, weil, obwohl}.
\end{itemize}
\end{aretenir}

\findecours

\end{document}
